%%%%%%%%%%%%%%%%%%%%%%%%%%%%%%%%%%%%%%%%%%%%%%%%%%%%%%%%%%%%%%%%%%%%%%%%%%%%%%%
%% TKS v7.3 COMPILER TRACK — Extended Language and Virtual Machine
%% Agent: tks-compiler
%%
%% This file contains the complete TKS language specification, compiler
%% architecture, effect system, and virtual machine.
%%
%% DEPENDENCIES:
%%   - v7.0 Type System (base types, typing rules)
%%   - v7.2 Algorithm W (Hindley-Milner inference)
%%   - v7.2 Compiler Architecture (lexer, parser, AST, IR, bytecode)
%%
%% STATUS: SKELETON — To be filled by tks-compiler agent
%%%%%%%%%%%%%%%%%%%%%%%%%%%%%%%%%%%%%%%%%%%%%%%%%%%%%%%%%%%%%%%%%%%%%%%%%%%%%%%

%% ============================================================================
%% CHAPTER 1: EXTENDED TKS LANGUAGE SPECIFICATION
%% ============================================================================
\chapter{Extended TKS Language Specification}
\label{ch:extended-language}

\begin{tcolorbox}[colback=blue!5!white,colframe=blue!75!black,title=\vnewthree]
This chapter provides the complete v7.3 TKS language specification, extending v7.2 with transfinite constructs, effect declarations, and advanced pattern matching.
\end{tcolorbox}

%% ----------------------------------------------------------------------------
\section{Review: v7.2 Language Core}
\label{sec:review-v72-language}

The TKS v7.2 language provides the foundational syntax upon which v7.3 builds. We summarize the core components for reference.

\begin{definition}[v7.2 Core Token Categories]
\label{def:v72-token-summary}
The v7.2 lexer produces tokens in the following categories:
\begin{align}
\mathsf{Token}_{7.2} ::=&\; \mathsf{ELEMENT}(W, n) \mid \mathsf{NOETIC}(k) \mid \mathsf{FOUNDATION}(m, w) \\
    &\mid \mathsf{INT}(n) \mid \mathsf{BOOL}(b) \mid \mathsf{IDENT}(s) \\
    &\mid \mathsf{LAMBDA} \mid \mathsf{LET} \mid \mathsf{IN} \mid \mathsf{IF} \mid \mathsf{THEN} \mid \mathsf{ELSE} \\
    &\mid \mathsf{RETURN} \mid \mathsf{BIND} \mid \mathsf{CHECK} \mid \mathsf{ACQUIRE} \mid \mathsf{ACBE} \\
    &\mid \mathsf{FRAC\_OPEN} \mid \mathsf{FRAC\_CLOSE} \mid \mathsf{COLON} \\
    &\mid \mathsf{LPAREN} \mid \mathsf{RPAREN} \mid \mathsf{ARROW} \mid \mathsf{EQUALS} \\
    &\mid \mathsf{PLUS} \mid \mathsf{MINUS} \mid \mathsf{TIMES} \mid \mathsf{DIVIDE} \mid \mathsf{EOF}
\end{align}
\end{definition}

\begin{definition}[v7.2 Core Grammar Summary]
\label{def:v72-grammar-summary}
The v7.2 grammar includes these primary productions:
\begin{lstlisting}[style=bnf,caption={v7.2 Core Grammar (Summary)}]
expr      ::= letexpr | lamexpr | ifexpr | bindexpr
letexpr   ::= 'let' IDENT '=' expr 'in' expr
lamexpr   ::= '\' IDENT '->' expr
ifexpr    ::= 'if' expr 'then' expr 'else' expr
bindexpr  ::= appexpr ('>>=' appexpr)*
appexpr   ::= primary+
primary   ::= ELEMENT | FOUNDATION | INT | BOOL | IDENT
            | '(' expr ')' | noetic | fractal | rpm
noetic    ::= 'nu' DIGIT '(' expr ')'
fractal   ::= '<' DIGIT (':' DIGIT)* '>' '(' expr ')'
rpm       ::= 'return' primary | 'check' primary | 'acquire' primary
\end{lstlisting}
\end{definition}

\begin{remark}[Backward Compatibility Principle]
\label{rem:backward-compat}
\textbf{All v7.2 syntax remains valid in v7.3.} The extensions described in this chapter add new keywords, tokens, and productions without modifying or removing existing ones. Any program valid in v6.1, v7.0, v7.1, or v7.2 remains valid and semantically unchanged in v7.3.
\end{remark}

%% ----------------------------------------------------------------------------
\section{Extended Lexical Structure}
\label{sec:extended-lexical}

Version 7.3 introduces three categories of lexical extensions: \emph{transfinite construct keywords}, \emph{effect system keywords}, and \emph{quantum operation keywords}.

\subsection{New Keywords}
\label{subsec:new-keywords}

\begin{definition}[Transfinite Keywords]
\label{def:transfinite-keywords}
Keywords for ordinal and transfinite constructs:
\begin{center}
\begin{tabular}{lll}
\toprule
\textbf{Keyword} & \textbf{Token} & \textbf{Purpose} \\
\midrule
\texttt{omega} & $\mathsf{OMEGA}$ & First infinite ordinal $\omega$ \\
\texttt{ord} & $\mathsf{ORD}$ & Ordinal type constructor / literal prefix \\
\texttt{limit} & $\mathsf{LIMIT}$ & Limit ordinal constructor / loop keyword \\
\texttt{succ} & $\mathsf{SUCC}$ & Successor ordinal constructor \\
\texttt{transfinite} & $\mathsf{TRANSFINITE}$ & Transfinite iteration marker \\
\bottomrule
\end{tabular}
\end{center}
\end{definition}

\begin{definition}[Effect System Keywords]
\label{def:effect-keywords}
Keywords for algebraic effects and handlers:
\begin{center}
\begin{tabular}{lll}
\toprule
\textbf{Keyword} & \textbf{Token} & \textbf{Purpose} \\
\midrule
\texttt{effect} & $\mathsf{EFFECT}$ & Effect declaration \\
\texttt{handle} & $\mathsf{HANDLE}$ & Effect handler introduction \\
\texttt{with} & $\mathsf{WITH}$ & Handler clause connector \\
\texttt{resume} & $\mathsf{RESUME}$ & Delimited continuation invocation \\
\texttt{perform} & $\mathsf{PERFORM}$ & Effect operation invocation \\
\texttt{handler} & $\mathsf{HANDLER}$ & Named handler declaration \\
\texttt{op} & $\mathsf{OP}$ & Operation declaration in effect \\
\bottomrule
\end{tabular}
\end{center}
\end{definition}

\begin{definition}[Quantum Operation Keywords]
\label{def:quantum-keywords}
Keywords for quantum noetic operations:
\begin{center}
\begin{tabular}{lll}
\toprule
\textbf{Keyword} & \textbf{Token} & \textbf{Purpose} \\
\midrule
\texttt{measure} & $\mathsf{MEASURE}$ & Quantum measurement projection \\
\texttt{superpose} & $\mathsf{SUPERPOSE}$ & Superposition construction \\
\texttt{entangle} & $\mathsf{ENTANGLE}$ & Entanglement creation \\
\texttt{qstate} & $\mathsf{QSTATE}$ & Quantum state type marker \\
\texttt{amplitude} & $\mathsf{AMPLITUDE}$ & Amplitude coefficient accessor \\
\bottomrule
\end{tabular}
\end{center}
\end{definition}

\subsection{Extended Token Set}
\label{subsec:extended-tokens}

\begin{definition}[v7.3 Token Set]
\label{def:v73-token-set}
The complete v7.3 token set extends v7.2:
\begin{align}
\mathsf{Token}_{7.3} ::=&\; \mathsf{Token}_{7.2} \quad \text{(all v7.2 tokens)} \\
    &\mid \mathsf{OMEGA} \mid \mathsf{ORD} \mid \mathsf{LIMIT} \mid \mathsf{SUCC} \mid \mathsf{TRANSFINITE} \\
    &\mid \mathsf{EFFECT} \mid \mathsf{HANDLE} \mid \mathsf{WITH} \mid \mathsf{RESUME} \mid \mathsf{PERFORM} \\
    &\mid \mathsf{HANDLER} \mid \mathsf{OP} \\
    &\mid \mathsf{MEASURE} \mid \mathsf{SUPERPOSE} \mid \mathsf{ENTANGLE} \mid \mathsf{QSTATE} \mid \mathsf{AMPLITUDE} \\
    &\mid \mathsf{ORDINAL\_LIT}(\alpha) \quad \text{(ordinal literals)} \\
    &\mid \mathsf{LBRACE} \mid \mathsf{RBRACE} \mid \mathsf{PIPE} \mid \mathsf{BANG} \\
    &\mid \mathsf{UNDERSCORE\_OMEGA} \quad \text{(subscript omega: } \texttt{\_}\omega \text{)}
\end{align}
\end{definition}

\begin{definition}[Ordinal Literal Lexing]
\label{def:ordinal-literal-lex}
Ordinal literals are recognized by the following patterns:
\begin{enumerate}
  \item $\omega$ alone: matches \texttt{omega} $\to \mathsf{ORDINAL\_LIT}(\omega)$
  \item $\omega + n$: matches \texttt{omega + } $[0\text{-}9]+$ $\to \mathsf{ORDINAL\_LIT}(\omega + n)$
  \item $\omega \cdot n$: matches \texttt{omega * } $[0\text{-}9]+$ $\to \mathsf{ORDINAL\_LIT}(\omega \cdot n)$
  \item $\omega^n$: matches \texttt{omega \^{} } $[0\text{-}9]+$ $\to \mathsf{ORDINAL\_LIT}(\omega^n)$
\end{enumerate}
For programmatic ordinal construction, use \texttt{ord} and \texttt{succ}/\texttt{limit} combinators.
\end{definition}

\begin{definition}[Lexer Extension Rules]
\label{def:lexer-extension}
The extended lexer $\mathsf{lex}_{7.3} : \mathsf{String} \to \mathsf{Token}_{7.3}^*$ adds:
\begin{enumerate}
  \item Match new keywords before identifiers (keyword priority)
  \item Match ordinal literals: \texttt{omega} followed optionally by arithmetic
  \item Match transfinite fractal subscripts: \texttt{\_omega}, \texttt{\_ord(...)}, etc.
  \item Match effect braces: \texttt{\{} $\to \mathsf{LBRACE}$, \texttt{\}} $\to \mathsf{RBRACE}$
  \item Match effect row operators: \texttt{|} $\to \mathsf{PIPE}$, \texttt{!} $\to \mathsf{BANG}$
\end{enumerate}
\end{definition}

%% ----------------------------------------------------------------------------
\section{Extended Grammar}
\label{sec:extended-grammar}

We now present the EBNF grammar extensions for v7.3. These productions extend (not replace) the v7.2 grammar.

\subsection{Top-Level Declarations}
\label{subsec:toplevel-grammar}

\begin{definition}[Extended Top-Level Grammar]
\label{def:toplevel-grammar}
New top-level declaration forms:
\begin{lstlisting}[style=bnf,caption={v7.3 Top-Level Declarations}]
decl        ::= v72_decl                   (* all v7.2 declarations *)
              | effectdecl                 (* effect declaration *)
              | handlerdecl                (* named handler *)

effectdecl  ::= 'effect' IDENT '{' oplist '}'

oplist      ::= opdecl (';' opdecl)*

opdecl      ::= 'op' IDENT '(' paramlist ')' ':' type

handlerdecl ::= 'handler' IDENT ':' effectname '->' type '{'
                clauselist
              '}'

clauselist  ::= clause (';' clause)*

clause      ::= 'return' IDENT '->' expr
              | IDENT '(' paramlist ')' IDENT '->' expr
              (* second IDENT is the continuation variable *)
\end{lstlisting}
\end{definition}

\subsection{Expression Extensions}
\label{subsec:expr-grammar}

\begin{definition}[Extended Expression Grammar]
\label{def:expr-grammar}
New expression forms for v7.3:
\begin{lstlisting}[style=bnf,caption={v7.3 Expression Extensions}]
expr        ::= v72_expr                   (* all v7.2 expressions *)
              | handleexpr                 (* effect handling *)
              | performexpr                (* effect invocation *)
              | transfiniteexpr            (* transfinite iteration *)
              | quantumexpr                (* quantum operations *)

handleexpr  ::= 'handle' expr 'with' '{' clauselist '}'
              | 'handle' expr 'with' IDENT  (* named handler *)

performexpr ::= 'perform' IDENT '(' arglist ')'

transfiniteexpr
            ::= transfinitefractal
              | transfiniteloop
              | ordinalexpr

quantumexpr ::= measureexpr
              | superposexpr
              | entangleexpr
\end{lstlisting}
\end{definition}

%% ----------------------------------------------------------------------------
\section{Transfinite Expression Syntax}
\label{sec:transfinite-syntax}

This section details the syntax for ordinal-indexed fractals, transfinite loops, and ordinal expressions.

\subsection{Transfinite Fractal Literals}
\label{subsec:transfinite-fractal-literals}

\begin{definition}[Transfinite Fractal Syntax]
\label{def:transfinite-fractal-syntax}
Transfinite fractals extend the v7.2 finite fractal syntax $\langle k_0 : k_1 : \cdots : k_{n-1} \rangle$:
\begin{lstlisting}[style=bnf,caption={Transfinite Fractal Grammar}]
transfinitefractal
            ::= '<' fractalseq '>' subscript '(' expr ')'

fractalseq  ::= DIGIT (':' DIGIT)* (':' '...')?
              (* finite prefix, optional ellipsis for infinite *)

subscript   ::= '_' ordinalatom
              | (* empty for finite fractals *)

ordinalatom ::= 'omega'
              | 'omega' '+' INT
              | 'omega' '*' INT
              | 'omega' '^' INT
              | 'ord' '(' ordinalexpr ')'
              | IDENT                      (* ordinal variable *)
\end{lstlisting}
\end{definition}

\begin{example}[Transfinite Fractal Literals]
\label{ex:transfinite-fractal-syntax}
The following illustrate valid v7.3 transfinite fractal syntax:
\begin{enumerate}
  \item \textbf{Finite (v7.2 compatible)}:
  \begin{verbatim}
    <1:4:7>(A3)                  -- Fractal Phi^3 = <1:4:7>, applied to A3
  \end{verbatim}

  \item \textbf{$\omega$-indexed constant fractal}:
  \begin{verbatim}
    <4:4:...>_omega(B5)          -- Phi^omega_4 = <4:4:4:...>, eigenfractal
  \end{verbatim}

  \item \textbf{$\omega$-indexed with finite prefix}:
  \begin{verbatim}
    <1:2:3:5:5:...>_omega(C7)    -- Finite prefix [1,2,3] then constant 5
  \end{verbatim}

  \item \textbf{Beyond $\omega$}:
  \begin{verbatim}
    <1:2:...>_(omega+1)(D2)      -- Phi^{omega+1} with k_omega specified
  \end{verbatim}

  \item \textbf{Programmatic ordinal}:
  \begin{verbatim}
    <3:3:...>_ord(alpha)(x)      -- Ordinal alpha from variable
  \end{verbatim}
\end{enumerate}
\end{example}

\begin{definition}[Ellipsis Interpretation]
\label{def:ellipsis-interp}
The ellipsis \texttt{...} in fractal syntax has precise semantics:
\begin{enumerate}
  \item \texttt{<k:...>\_omega} denotes the constant $\omega$-fractal $\TransFrac{\omega}_k$
  \item \texttt{<k1:k2:...:kn:...>\_omega} denotes the eventually-constant fractal with finite prefix $\langle k_1, \ldots, k_n \rangle$ repeating $k_n$
  \item For user-defined patterns, use \texttt{transfinite} with a generating function (see below)
\end{enumerate}
\end{definition}

\subsection{Transfinite Loop Syntax}
\label{subsec:transfinite-loop}

\begin{definition}[Transfinite Loop Grammar]
\label{def:transfinite-loop-grammar}
Iteration over ordinals:
\begin{lstlisting}[style=bnf,caption={Transfinite Loop Grammar}]
transfiniteloop
            ::= 'transfinite' 'loop' IDENT '<' ordinalexpr
                'from' expr
                'step' '(' IDENT '->' expr ')'
                'limit' '(' IDENT '->' expr ')'

ordinalexpr ::= ordinalatom
              | 'succ' '(' ordinalexpr ')'
              | 'limit' '(' IDENT '->' ordinalexpr ')'
              | ordinalexpr '+' ordinalexpr
              | ordinalexpr '*' ordinalexpr
              | ordinalexpr '^' ordinalexpr
\end{lstlisting}
\end{definition}

\begin{example}[Transfinite Loop]
\label{ex:transfinite-loop}
Computing the $\omega$-limit of iterated noetic application:
\begin{verbatim}
transfinite loop beta < omega
  from A1                                -- initial value
  step (x -> nu3(x))                     -- successor step
  limit (f -> limit_noetic(f))           -- limit construction
\end{verbatim}
This iterates $\nu_3$ transfinitely, computing $\lim_{\beta \to \omega} \nu_3^\beta(\mathsf{A1})$.
\end{example}

%% ----------------------------------------------------------------------------
\section{Effect Declaration Syntax}
\label{sec:effect-declaration-syntax}

Algebraic effects in v7.3 enable modular composition of computational effects including RPM operations, quantum measurements, and I/O.

\subsection{Effect Declarations}
\label{subsec:effect-decl}

\begin{definition}[Effect Declaration Grammar]
\label{def:effect-decl-grammar}
\begin{lstlisting}[style=bnf,caption={Effect Declaration Grammar}]
effectdecl  ::= 'effect' IDENT typeargs? '{' oplist '}'

typeargs    ::= '[' IDENT (',' IDENT)* ']'

oplist      ::= opdecl (';' opdecl)* ';'?

opdecl      ::= 'op' IDENT '(' paramlist ')' ':' type
              | 'op' IDENT ':' type        (* nullary operation *)

paramlist   ::= (* empty *)
              | param (',' param)*

param       ::= IDENT ':' type
\end{lstlisting}
\end{definition}

\begin{example}[RPM Effect Declaration]
\label{ex:rpm-effect}
The RPM effect capturing prerequisite acquisition:
\begin{verbatim}
effect RPMEffect {
  op check(prereq: Element): Bool;
  op acquire(elem: Element): Unit;
  op fail(reason: String): Void
}
\end{verbatim}
Operations return to their continuation except \texttt{fail}, which has return type \texttt{Void} indicating non-resumption.
\end{example}

\begin{example}[Quantum Effect Declaration]
\label{ex:quantum-effect}
Quantum operations as effects:
\begin{verbatim}
effect QuantumEffect[A] {
  op measure(q: QState[A]): A;
  op superpose(states: List[(Complex, A)]): QState[A];
  op entangle(q1: QState[A], q2: QState[A]): QState[(A, A)]
}
\end{verbatim}
\end{example}

%% ----------------------------------------------------------------------------
\section{Handler Syntax}
\label{sec:handler-syntax}

Effect handlers provide interpretations for effect operations, using delimited continuations.

\subsection{Handler Expressions}
\label{subsec:handler-expr}

\begin{definition}[Handler Expression Grammar]
\label{def:handler-expr-grammar}
\begin{lstlisting}[style=bnf,caption={Handler Expression Grammar}]
handleexpr  ::= 'handle' expr 'with' handlerblock
              | 'handle' expr 'with' IDENT

handlerblock::= '{' clauselist '}'

clauselist  ::= clause (';' clause)* ';'?

clause      ::= returnclause | opclause

returnclause::= 'return' IDENT '->' expr

opclause    ::= IDENT '(' paramlist ')' IDENT '->' expr
              (* opname(params) continuation -> body *)
\end{lstlisting}
\end{definition}

\begin{definition}[Resume Syntax]
\label{def:resume-syntax}
Within handler clauses, \texttt{resume} invokes the captured continuation:
\begin{lstlisting}[style=bnf,caption={Resume Syntax}]
resumeexpr  ::= 'resume' '(' expr ')'
              | 'resume' expr             (* without parens for simple args *)
\end{lstlisting}
The argument to \texttt{resume} provides the return value of the operation to the calling context.
\end{definition}

\begin{example}[RPM Handler]
\label{ex:rpm-handler}
A handler interpreting RPM effects with state:
\begin{verbatim}
handle
  let x = perform acquire(A3) in
  let y = perform check(B5) in
  if y then x else perform fail("B5 not acquired")
with {
  return v -> return (v, currentState);

  acquire(elem) k ->
    let newState = addToState(elem, currentState) in
    resume(()) with newState;

  check(prereq) k ->
    let satisfied = hasPrereq(prereq, currentState) in
    resume(satisfied);

  fail(reason) k ->
    return (Error(reason), currentState)   (* no resume *)
}
\end{verbatim}
\end{example}

\subsection{Quantum Operation Syntax}
\label{subsec:quantum-syntax}

\begin{definition}[Quantum Operation Grammar]
\label{def:quantum-op-grammar}
\begin{lstlisting}[style=bnf,caption={Quantum Operation Grammar}]
measureexpr ::= 'measure' '(' expr ')' 'in' basis?
              | 'measure' expr

basis       ::= 'basis' '{' basislist '}'

basislist   ::= basisvec (',' basisvec)*

basisvec    ::= '|' IDENT '>'

superposexpr::= 'superpose' '{' superposelist '}'

superposelist
            ::= superposeterm (',' superposeterm)*

superposeterm
            ::= amplitude ':' '|' expr '>'

amplitude   ::= complexlit | IDENT | '(' expr ')'

complexlit  ::= FLOAT                      (* real *)
              | FLOAT 'i'                  (* imaginary *)
              | '(' FLOAT ',' FLOAT ')'    (* (re, im) *)

entangleexpr::= 'entangle' '(' expr ',' expr ')'
              | 'entangle' '{' entanglelist '}'

entanglelist::= entanglepair (';' entanglepair)*

entanglepair::= '|' expr '>' '<->' '|' expr '>'
\end{lstlisting}
\end{definition}

\begin{example}[Quantum Operations]
\label{ex:quantum-ops}
Quantum noetic operations in v7.3 syntax:
\begin{verbatim}
-- Create superposition of Ideas A1 and A2
let psi = superpose {
  (0.707, 0): |A1>,
  (0.707, 0): |A2>
} in

-- Entangle with another Idea
let entangled = entangle(psi, |B3>) in

-- Measure in computational basis
let result = measure(entangled) in
result
\end{verbatim}
\end{example}

%% ----------------------------------------------------------------------------
\section{Summary: v7.3 Syntax Extensions}
\label{sec:syntax-summary}

\begin{definition}[v7.3 Extension Summary]
\label{def:v73-summary}
The v7.3 language specification adds:
\begin{enumerate}
  \item \textbf{17 new keywords}: 5 transfinite, 7 effect system, 5 quantum
  \item \textbf{Transfinite fractal literals}: $\langle k : k : \ldots \rangle_\omega$ syntax
  \item \textbf{Effect declarations}: \texttt{effect Name \{ op ...; ... \}}
  \item \textbf{Handler expressions}: \texttt{handle e with \{ ... \}}
  \item \textbf{Quantum operations}: \texttt{measure}, \texttt{superpose}, \texttt{entangle}
  \item \textbf{Transfinite loops}: \texttt{transfinite loop} with step/limit clauses
\end{enumerate}
All v6.1--v7.2 syntax remains valid and semantically unchanged.
\end{definition}

\begin{theorem}[Grammar Extension Conservativity]
\label{thm:grammar-conservative}
Let $\mathcal{G}_{7.2}$ denote the v7.2 grammar and $\mathcal{G}_{7.3}$ the v7.3 grammar. Then:
\begin{enumerate}
  \item $\mathcal{L}(\mathcal{G}_{7.2}) \subset \mathcal{L}(\mathcal{G}_{7.3})$ (strict language inclusion)
  \item For any $p \in \mathcal{L}(\mathcal{G}_{7.2})$: $\mathsf{parse}_{7.3}(p) = \mathsf{parse}_{7.2}(p)$ (AST preservation)
  \item $\mathcal{G}_{7.3}$ remains unambiguous if $\mathcal{G}_{7.2}$ is unambiguous
\end{enumerate}
\end{theorem}

\begin{proof}[Proof Sketch]
(1) follows from the extension-only nature of v7.3 productions. (2) follows from new keywords not overlapping with v7.2 keywords or identifier patterns. (3) follows from careful precedence assignment to new constructs and the LALR(1) property being preserved by the extensions.
\end{proof}

%% ----------------------------------------------------------------------------
%% HANDOFF NOTE
%% ----------------------------------------------------------------------------
\begin{tcolorbox}[colback=green!5!white,colframe=green!50!black,title=\textbf{Handoff: COMPILER-TASK-01 Complete},breakable]
\textbf{Completed in this section:}
\begin{itemize}
  \item Review of v7.2 language core (tokens, grammar)
  \item Extended lexical structure with 17 new keywords
  \item EBNF grammar for transfinite constructs, effects, quantum operations
  \item Transfinite fractal literal syntax $\langle k:k:\ldots\rangle_\omega$
  \item Effect declaration and handler syntax
  \item Basic quantum operation syntax
  \item Examples demonstrating new syntactic forms
\end{itemize}

\textbf{Next tasks (COMPILER-TASK-02):}
\begin{itemize}
  \item Extended AST node definitions for v7.3 constructs
  \item Parsing algorithm updates (recursive descent / LR extensions)
  \item Desugaring rules for transfinite fractals to core calculus
  \item Type annotations for effect rows (Section~\ref{sec:effect-types})
\end{itemize}
\end{tcolorbox}

%% ============================================================================
%% CHAPTER 2: EXTENDED TYPE SYSTEM
%% ============================================================================
\chapter{Extended Type System}
\label{ch:extended-type-system}

\begin{tcolorbox}[colback=blue!5!white,colframe=blue!75!black,title=\vnewthree]
This chapter extends the v7.2 type system with effect types, row polymorphism, and ordinal-indexed types.
\end{tcolorbox}

%% ----------------------------------------------------------------------------
\section{Review: v7.2 Type System Core}
\label{sec:review-v72-types}

We begin by summarizing the v7.2 type system upon which v7.3 builds.

\begin{definition}[v7.2 Type Syntax (Recap)]
\label{def:v72-type-recap}
The v7.2 type grammar:
\begin{align}
\tau_{7.2} ::=&\; \alpha \quad \text{(type variable)} \\
    &\mid \mathsf{Int} \mid \mathsf{Bool} \mid \mathsf{Unit} \mid \mathsf{Void} \\
    &\mid \mathsf{Element}[W] \mid \mathsf{Domain} \mid \mathsf{Foundation} \\
    &\mid \mathsf{Frac}[\bar{k}] \mid \mathsf{RPM}[\tau] \\
    &\mid \tau_1 \to \tau_2 \mid \tau_1 \times \tau_2 \mid \tau_1 + \tau_2
\end{align}
Type schemes: $\sigma = \forall \bar{\alpha}.\, \tau$
\end{definition}

\begin{remark}[v7.3 Type System Extension Principle]
\label{rem:type-extension}
The v7.3 type system extends v7.2 in three dimensions:
\begin{enumerate}
  \item \textbf{Effect tracking}: Types annotated with effect rows
  \item \textbf{Ordinal indexing}: Types parameterized by ordinals for transfinite structures
  \item \textbf{World stratification}: Refined world-aware typing with ACBE coercions
\end{enumerate}
All v7.2 types embed into v7.3 via the pure effect row $\{\}$.
\end{remark}

%% ----------------------------------------------------------------------------
\section{Effect Types}
\label{sec:effect-types}

Effect types track which computational effects an expression may perform, enabling static verification of effect usage and handler coverage.

\subsection{Effect Type Syntax}
\label{subsec:effect-type-syntax}

\begin{definition}[Effect Names]
\label{def:effect-names}
An \textbf{effect name} $E$ is an identifier declared via \texttt{effect} (Section~\ref{sec:effect-declaration-syntax}). The set of declared effects is denoted $\mathcal{E}$. Core TKS effects include:
\begin{align}
\mathcal{E}_{\mathrm{core}} = \{&\, \mathsf{RPMEffect}, \mathsf{QuantumEffect}, \mathsf{IOEffect}, \\
    &\, \mathsf{StateEffect}[\tau], \mathsf{ExnEffect}[\tau], \mathsf{NonDetEffect} \,\}
\end{align}
\end{definition}

\begin{definition}[Effect Row Syntax]
\label{def:effect-row-syntax}
An \textbf{effect row} $\varepsilon$ describes a set of effects:
\begin{align}
\varepsilon ::=&\; \{\} \quad \text{(empty row: pure)} \\
    &\mid \{E_1, E_2, \ldots, E_n\} \quad \text{(closed row)} \\
    &\mid \{E_1, \ldots, E_n \mid \rho\} \quad \text{(open row with row variable } \rho \text{)} \\
    &\mid \rho \quad \text{(row variable)}
\end{align}
Row variables $\rho \in \{\varrho, \varrho', \varrho_1, \ldots\}$ enable effect polymorphism.
\end{definition}

\begin{definition}[Effectful Type Syntax]
\label{def:effectful-type}
The v7.3 type grammar extends v7.2 with effect annotations:
\begin{align}
\tau_{7.3} ::=&\; \tau_{7.2} \quad \text{(all v7.2 types)} \\
    &\mid \tau \mathbin{!} \varepsilon \quad \text{(effectful computation type)} \\
    &\mid \tau_1 \xrightarrow{\varepsilon} \tau_2 \quad \text{(effectful function type)} \\
    &\mid \mathsf{Handler}[E, \tau_{\mathrm{in}}, \tau_{\mathrm{out}}] \quad \text{(handler type)}
\end{align}
\end{definition}

\begin{notation}[Effect Type Conventions]
\label{not:effect-conventions}
\begin{enumerate}
  \item $\tau \mathbin{!} \{\}$ is written simply as $\tau$ (pure computation)
  \item $\tau_1 \to \tau_2$ abbreviates $\tau_1 \xrightarrow{\{\}} \tau_2$ (pure function)
  \item $\tau_1 \xrightarrow{E} \tau_2$ abbreviates $\tau_1 \xrightarrow{\{E\}} \tau_2$ (single effect)
  \item $\mathsf{RPM}[\tau]$ from v7.2 is now $\tau \mathbin{!} \{\mathsf{RPMEffect}\}$
\end{enumerate}
\end{notation}

\begin{example}[Effect Types in Practice]
\label{ex:effect-types}
\begin{align}
&\mathsf{acquire} : \mathsf{Element}[W] \xrightarrow{\mathsf{RPMEffect}} \mathsf{Unit} \\
&\mathsf{measure} : \mathsf{QState}[\tau] \xrightarrow{\mathsf{QuantumEffect}} \tau \\
&\mathsf{readFile} : \mathsf{String} \xrightarrow{\mathsf{IOEffect}} \mathsf{String} \\
&\mathsf{pure\_add} : \mathsf{Int} \to \mathsf{Int} \to \mathsf{Int} \quad \text{(no effect annotation)}
\end{align}
\end{example}

%% ----------------------------------------------------------------------------
\section{Row Polymorphism for Effects}
\label{sec:row-polymorphism}

Row polymorphism enables writing functions that are polymorphic over which additional effects may occur, crucial for effect-generic programming.

\subsection{Row Variables and Quantification}
\label{subsec:row-vars}

\begin{definition}[Effect-Polymorphic Type Scheme]
\label{def:effect-poly-scheme}
An \textbf{effect-polymorphic type scheme} quantifies over both type variables and row variables:
\[
\sigma = \forall \bar{\alpha}.\, \forall \bar{\rho}.\, \tau
\]
where $\bar{\alpha}$ are type variables and $\bar{\rho}$ are row variables.
\end{definition}

\begin{definition}[Free Row Variables]
\label{def:free-row-vars}
The free row variables $\FRV$ are defined:
\begin{align}
\FRV(\{\}) &= \varnothing \\
\FRV(\{E_1, \ldots, E_n\}) &= \varnothing \\
\FRV(\{E_1, \ldots, E_n \mid \rho\}) &= \{\rho\} \\
\FRV(\rho) &= \{\rho\} \\
\FRV(\tau \mathbin{!} \varepsilon) &= \FTV(\tau) \cup \FRV(\varepsilon) \\
\FRV(\tau_1 \xrightarrow{\varepsilon} \tau_2) &= \FRV(\tau_1) \cup \FRV(\varepsilon) \cup \FRV(\tau_2)
\end{align}
\end{definition}

\begin{example}[Effect-Polymorphic Functions]
\label{ex:effect-poly}
\begin{enumerate}
  \item \textbf{Monadic bind} (effect-preserving):
  \[
  (\gg\!=) : \forall \alpha\, \beta\, \rho.\; (\alpha \mathbin{!} \rho) \to (\alpha \xrightarrow{\rho} \beta \mathbin{!} \rho) \to (\beta \mathbin{!} \rho)
  \]

  \item \textbf{Effect-polymorphic map}:
  \[
  \mathsf{map} : \forall \alpha\, \beta\, \rho.\; (\alpha \xrightarrow{\rho} \beta) \to \mathsf{List}[\alpha] \xrightarrow{\rho} \mathsf{List}[\beta]
  \]

  \item \textbf{Effect extension} (adding an effect):
  \[
  \mathsf{withRPM} : \forall \alpha\, \rho.\; (\alpha \mathbin{!} \rho) \to (\alpha \mathbin{!} \{\mathsf{RPMEffect} \mid \rho\})
  \]
\end{enumerate}
\end{example}

\subsection{Row Unification}
\label{subsec:row-unification}

\begin{definition}[Row Equivalence]
\label{def:row-equiv}
Effect rows are equivalent up to permutation and idempotence:
\begin{align}
\{E_1, E_2 \mid \rho\} &\equiv \{E_2, E_1 \mid \rho\} \quad \text{(commutativity)} \\
\{E, E \mid \rho\} &\equiv \{E \mid \rho\} \quad \text{(idempotence)}
\end{align}
Rows form a commutative idempotent monoid under union with identity $\{\}$.
\end{definition}

\begin{definition}[Row Unification Algorithm]
\label{def:row-unify}
$\UnifyRow(\varepsilon_1, \varepsilon_2)$ extends standard unification:
\begin{align}
\UnifyRow(\{\}, \{\}) &= [] \\
\UnifyRow(\rho, \varepsilon) &= [\rho \mapsto \varepsilon] \quad \text{if } \rho \notin \FRV(\varepsilon) \\
\UnifyRow(\{E \mid \rho_1\}, \{E \mid \rho_2\}) &= \UnifyRow(\rho_1, \rho_2) \\
\UnifyRow(\{E_1 \mid \rho_1\}, \{E_2 \mid \rho_2\}) &= [\rho_1 \mapsto \{E_2 \mid \rho'\}, \rho_2 \mapsto \{E_1 \mid \rho'\}] \\
    &\quad \text{where } \rho' \text{ is fresh, if } E_1 \neq E_2
\end{align}
The last case handles \textbf{row extension}: unifying rows with different head effects.
\end{definition}

\begin{theorem}[Row Unification Soundness]
\label{thm:row-unify-sound}
If $\UnifyRow(\varepsilon_1, \varepsilon_2) = \theta$, then $\theta(\varepsilon_1) \equiv \theta(\varepsilon_2)$.
\end{theorem}

%% ----------------------------------------------------------------------------
\section{Ordinal-Indexed Types}
\label{sec:ordinal-indexed-types}

Ordinal-indexed types enable static tracking of transfinite structure depths, essential for typing transfinite fractals and ensuring termination of ordinal-bounded computations.

\subsection{Ordinal Type Indices}
\label{subsec:ordinal-indices}

\begin{definition}[Type-Level Ordinals]
\label{def:type-level-ordinals}
The grammar of type-level ordinal expressions:
\begin{align}
\kappa ::=&\; 0 \mid 1 \mid 2 \mid \cdots \quad \text{(finite ordinal literals)} \\
    &\mid \omega \mid \omega_1 \quad \text{(infinite ordinal constants)} \\
    &\mid \xi \quad \text{(ordinal variable)} \\
    &\mid \mathsf{succ}(\kappa) \quad \text{(successor)} \\
    &\mid \kappa_1 + \kappa_2 \mid \kappa_1 \cdot \kappa_2 \mid \kappa_1^{\kappa_2} \quad \text{(ordinal arithmetic)}
\end{align}
Ordinal variables $\xi \in \{\zeta, \zeta', \zeta_1, \ldots\}$ enable ordinal polymorphism.
\end{definition}

\begin{definition}[Ordinal-Indexed Type Constructors]
\label{def:ordinal-indexed-types-def}
Types parameterized by ordinals:
\begin{align}
\tau ::=&\; \cdots \quad \text{(all previous types)} \\
    &\mid \mathsf{Frac}^\kappa[\bar{k}] \quad \text{(fractal of ordinal depth } \kappa \text{)} \\
    &\mid \mathsf{Ord} \quad \text{(ordinal type)} \\
    &\mid \mathsf{OrdBounded}[\kappa] \quad \text{(ordinals } < \kappa \text{)} \\
    &\mid \mathsf{TransIter}[\kappa, \tau] \quad \text{(transfinite iteration result)}
\end{align}
\end{definition}

\begin{example}[Ordinal-Indexed Types]
\label{ex:ordinal-types}
\begin{enumerate}
  \item \textbf{Finite fractal}: $\mathsf{Frac}^3[\langle 1,4,7 \rangle] : \mathsf{Element}[W] \to \mathsf{Element}[W]$
  \item \textbf{$\omega$-fractal}: $\mathsf{Frac}^\omega_k : \mathsf{Element}[W] \to \mathsf{Element}[W]$
  \item \textbf{Transfinite loop bound}: $\mathsf{OrdBounded}[\omega] \cong \mathbb{N}$
  \item \textbf{Iteration at $\omega$}: $\mathsf{TransIter}[\omega, \mathsf{Element}[A]]$
\end{enumerate}
\end{example}

\subsection{Bounded Ordinal Quantification}
\label{subsec:bounded-ordinal}

\begin{definition}[Ordinal-Bounded Quantification]
\label{def:ordinal-bounded-quant}
Ordinal-polymorphic type schemes with bounds:
\[
\sigma = \forall (\xi < \kappa).\, \tau
\]
The bound $\xi < \kappa$ constrains instantiations of $\xi$ to ordinals strictly less than $\kappa$.
\end{definition}

\begin{definition}[Ordinal Constraint System]
\label{def:ordinal-constraints}
During type inference, we collect ordinal constraints:
\begin{align}
C ::=&\; \kappa_1 = \kappa_2 \quad \text{(equality)} \\
    &\mid \kappa_1 < \kappa_2 \quad \text{(strict ordering)} \\
    &\mid \kappa_1 \leq \kappa_2 \quad \text{(non-strict ordering)} \\
    &\mid C_1 \land C_2 \quad \text{(conjunction)}
\end{align}
Constraint satisfaction is decidable for the ordinal fragment we consider (polynomial ordinal expressions below $\omega_1$).
\end{definition}

\begin{example}[Ordinal-Polymorphic Fractal Composition]
\label{ex:ordinal-poly-frac}
\[
\circ_{\mathrm{Ord}} : \forall (\xi_1 < \omega_1).\, \forall (\xi_2 < \omega_1).\;
  \mathsf{Frac}^{\xi_1} \to \mathsf{Frac}^{\xi_2} \to \mathsf{Frac}^{\xi_1 + \xi_2}
\]
The result type precisely tracks that composing fractals of depths $\xi_1$ and $\xi_2$ yields depth $\xi_1 + \xi_2$.
\end{example}

%% ----------------------------------------------------------------------------
\section{World-Stratified Types}
\label{sec:world-stratified-types}

TKS Ideas are organized into four Worlds (A, B, C, D), and the type system tracks World membership to ensure ACBE-compliance.

\subsection{World-Parameterized Types}
\label{subsec:world-param}

\begin{definition}[World Type Parameter]
\label{def:world-type-param}
World parameters in types:
\begin{align}
W ::=&\; A \mid B \mid C \mid D \quad \text{(concrete worlds)} \\
    &\mid \omega \quad \text{(world variable)} \\
    &\mid \mathsf{Any} \quad \text{(any world)}
\end{align}
World variables $\omega \in \{\omega_A, \omega_B, \ldots\}$ (not to be confused with the ordinal $\omega$) enable world polymorphism.
\end{definition}

\begin{definition}[World-Stratified Element Type]
\label{def:world-element}
The element type is parameterized by world:
\[
\mathsf{Element}[W] \quad \text{where } W \in \{A, B, C, D, \omega, \mathsf{Any}\}
\]
with inhabitants:
\begin{align}
\mathsf{Element}[A] &= \{A1, A2, \ldots, A10\} \\
\mathsf{Element}[B] &= \{B1, B2, \ldots, B10\} \\
&\vdots \\
\mathsf{Element}[\mathsf{Any}] &= \bigcup_{W \in \{A,B,C,D\}} \mathsf{Element}[W]
\end{align}
\end{definition}

\begin{definition}[World Polymorphism]
\label{def:world-poly}
World-polymorphic type schemes:
\[
\sigma = \forall \omega.\, \tau
\]
Example: $\mathsf{id}_{\mathsf{Elem}} : \forall \omega.\, \mathsf{Element}[\omega] \to \mathsf{Element}[\omega]$
\end{definition}

\subsection{ACBE-Respecting Type Coercions}
\label{subsec:acbe-coercions}

\begin{definition}[World Ordering for ACBE]
\label{def:world-ordering}
The ACBE functor respects a partial ordering on worlds based on abstractness:
\[
A <_{\mathrm{ACBE}} B <_{\mathrm{ACBE}} C <_{\mathrm{ACBE}} D
\]
This ordering reflects the progression from Abstract (A) to Dense (D).
\end{definition}

\begin{definition}[ACBE Coercion Type]
\label{def:acbe-coercion-type}
The ACBE functor has type:
\[
\mathsf{ACBE} : \forall \omega_1\, \omega_2.\, (\omega_1 <_{\mathrm{ACBE}} \omega_2) \Rightarrow \mathsf{Element}[\omega_1] \to \mathsf{Element}[\omega_2]
\]
The constraint $\omega_1 <_{\mathrm{ACBE}} \omega_2$ is checked statically.
\end{definition}

\begin{definition}[World Constraint System]
\label{def:world-constraints}
World constraints during type inference:
\begin{align}
W_C ::=&\; \omega = W \quad \text{(world equality)} \\
    &\mid \omega_1 <_{\mathrm{ACBE}} \omega_2 \quad \text{(ACBE ordering)} \\
    &\mid \omega \in \{W_1, \ldots, W_n\} \quad \text{(world membership)}
\end{align}
\end{definition}

%% ----------------------------------------------------------------------------
\section{Subtyping for Effects}
\label{sec:effect-subtyping}

Effect subtyping enables safe subsumption: a computation with fewer effects can be used where more effects are allowed.

\subsection{Effect Row Subtyping}
\label{subsec:effect-subtyping-rules}

\begin{definition}[Effect Row Subtyping]
\label{def:effect-row-subtype}
The subtyping relation $\varepsilon_1 <: \varepsilon_2$ on effect rows:
\[
\infer[\textsc{Sub-Empty}]
  {\{\} <: \varepsilon}
  {}
\qquad
\infer[\textsc{Sub-Row}]
  {\{E_1, \ldots, E_n\} <: \{E_1, \ldots, E_n, E_{n+1}, \ldots, E_m\}}
  {}
\]
\[
\infer[\textsc{Sub-Open}]
  {\{E_1, \ldots, E_n \mid \rho\} <: \{E_1, \ldots, E_n, E_{n+1}, \ldots, E_m \mid \rho\}}
  {}
\]
\end{definition}

\begin{proposition}[Subtyping Properties]
\label{prop:subtyping-props}
Effect row subtyping satisfies:
\begin{enumerate}
  \item \textbf{Reflexivity}: $\varepsilon <: \varepsilon$
  \item \textbf{Transitivity}: $\varepsilon_1 <: \varepsilon_2 \land \varepsilon_2 <: \varepsilon_3 \Rightarrow \varepsilon_1 <: \varepsilon_3$
  \item \textbf{Purity is bottom}: $\{\} <: \varepsilon$ for all $\varepsilon$
\end{enumerate}
\end{proposition}

\subsection{Type Subtyping with Effects}
\label{subsec:type-subtyping}

\begin{definition}[Effectful Type Subtyping]
\label{def:effectful-subtype}
Subtyping on effectful types:
\[
\infer[\textsc{Sub-Eff}]
  {\tau \mathbin{!} \varepsilon_1 <: \tau \mathbin{!} \varepsilon_2}
  {\varepsilon_1 <: \varepsilon_2}
\]
\[
\infer[\textsc{Sub-Fun}]
  {\tau_1 \xrightarrow{\varepsilon_1} \tau_2 <: \tau_1' \xrightarrow{\varepsilon_2} \tau_2'}
  {\tau_1' <: \tau_1 & \varepsilon_1 <: \varepsilon_2 & \tau_2 <: \tau_2'}
\]
Note the contravariance in the argument type.
\end{definition}

\begin{definition}[Subsumption Rule]
\label{def:subsumption}
The subsumption typing rule:
\[
\infer[\textsc{T-Sub}]
  {\Gamma \vdash e : \tau_2 \mathbin{!} \varepsilon_2}
  {\Gamma \vdash e : \tau_1 \mathbin{!} \varepsilon_1 & \tau_1 \mathbin{!} \varepsilon_1 <: \tau_2 \mathbin{!} \varepsilon_2}
\]
\end{definition}

%% ----------------------------------------------------------------------------
\section{Typing Rules for v7.3 Constructs}
\label{sec:typing-rules}

We now present the key typing judgments for v7.3 constructs. The judgment form is:
\[
\Gamma \vdash e : \tau \mathbin{!} \varepsilon
\]
meaning ``under context $\Gamma$, expression $e$ has type $\tau$ and may perform effects in $\varepsilon$''.

\subsection{Effect Operation Typing}
\label{subsec:effect-op-typing}

\begin{definition}[Effect Operation Typing Rules]
\label{def:effect-op-rules}
\[
\infer[\textsc{T-Perform}]
  {\Gamma \vdash \mathsf{perform}\; \mathit{op}(e_1, \ldots, e_n) : \tau_{\mathrm{ret}} \mathbin{!} \{E \mid \rho\}}
  {
    \mathit{op} : (\tau_1, \ldots, \tau_n) \to \tau_{\mathrm{ret}} \in E &
    \Gamma \vdash e_i : \tau_i \mathbin{!} \rho \text{ for each } i
  }
\]
\end{definition}

\subsection{Handler Typing}
\label{subsec:handler-typing}

\begin{definition}[Handler Typing Rules]
\label{def:handler-rules}
\[
\infer[\textsc{T-Handle}]
  {\Gamma \vdash \mathsf{handle}\; e\; \mathsf{with}\; h : \tau_{\mathrm{out}} \mathbin{!} \varepsilon}
  {
    \Gamma \vdash e : \tau_{\mathrm{in}} \mathbin{!} \{E \mid \varepsilon\} &
    \Gamma \vdash h : \mathsf{Handler}[E, \tau_{\mathrm{in}}, \tau_{\mathrm{out}}] \mathbin{!} \varepsilon
  }
\]
The handler $h$ removes effect $E$ from the effect row, transforming $\tau_{\mathrm{in}}$ to $\tau_{\mathrm{out}}$.
\end{definition}

\begin{definition}[Handler Clause Typing]
\label{def:handler-clause}
For a handler clause $\mathit{op}(x_1, \ldots, x_n)\; k \to e_{\mathrm{body}}$ handling operation $\mathit{op} : (\tau_1, \ldots, \tau_n) \to \tau_r \in E$:
\[
\infer[\textsc{T-OpClause}]
  {\Gamma \vdash (\mathit{op}(\bar{x})\; k \to e) : \mathsf{OpClause}[E, \tau_{\mathrm{out}}]}
  {
    \Gamma, x_1 : \tau_1, \ldots, x_n : \tau_n, k : \tau_r \to \tau_{\mathrm{out}} \mathbin{!} \varepsilon
    \vdash e : \tau_{\mathrm{out}} \mathbin{!} \varepsilon
  }
\]
The continuation $k$ has type $\tau_r \to \tau_{\mathrm{out}} \mathbin{!} \varepsilon$, representing the delimited continuation.
\end{definition}

\subsection{Transfinite Fractal Typing}
\label{subsec:transfinite-typing}

\begin{definition}[Transfinite Fractal Typing Rules]
\label{def:transfinite-frac-rules}
\[
\infer[\textsc{T-TransFrac}]
  {\Gamma \vdash \langle \bar{k} \rangle_\kappa(e) : \mathsf{Element}[\omega] \mathbin{!} \varepsilon}
  {
    \Gamma \vdash e : \mathsf{Element}[\omega] \mathbin{!} \varepsilon &
    \Gamma \vdash \kappa : \mathsf{Ord} &
    \bar{k} : \kappa \to \{0, \ldots, 9\}
  }
\]

\[
\infer[\textsc{T-TransLoop}]
  {\Gamma \vdash \mathsf{transfinite}\;\mathsf{loop}\; \xi < \kappa\; \mathsf{from}\; e_0\; \mathsf{step}\; f\; \mathsf{limit}\; g : \tau \mathbin{!} \varepsilon}
  {
    \begin{array}{c}
    \Gamma \vdash e_0 : \tau \mathbin{!} \varepsilon \\
    \Gamma \vdash f : \tau \xrightarrow{\varepsilon} \tau \\
    \Gamma \vdash g : (\mathsf{OrdBounded}[\kappa] \to \tau) \xrightarrow{\varepsilon} \tau
    \end{array}
  }
\]
\end{definition}

\subsection{Quantum Operation Typing}
\label{subsec:quantum-typing}

\begin{definition}[Quantum Operation Typing Rules]
\label{def:quantum-rules}
\[
\infer[\textsc{T-Superpose}]
  {\Gamma \vdash \mathsf{superpose}\;\{c_i : |e_i\rangle\}_i : \mathsf{QState}[\tau] \mathbin{!} \varepsilon}
  {
    \Gamma \vdash c_i : \mathbb{C} \mathbin{!} \varepsilon &
    \Gamma \vdash e_i : \tau \mathbin{!} \varepsilon &
    \sum_i |c_i|^2 = 1
  }
\]

\[
\infer[\textsc{T-Measure}]
  {\Gamma \vdash \mathsf{measure}(e) : \tau \mathbin{!} \{\mathsf{QuantumEffect} \mid \varepsilon\}}
  {\Gamma \vdash e : \mathsf{QState}[\tau] \mathbin{!} \varepsilon}
\]

\[
\infer[\textsc{T-Entangle}]
  {\Gamma \vdash \mathsf{entangle}(e_1, e_2) : \mathsf{QState}[\tau_1 \times \tau_2] \mathbin{!} \{\mathsf{QuantumEffect} \mid \varepsilon\}}
  {
    \Gamma \vdash e_1 : \mathsf{QState}[\tau_1] \mathbin{!} \varepsilon &
    \Gamma \vdash e_2 : \mathsf{QState}[\tau_2] \mathbin{!} \varepsilon
  }
\]
\end{definition}

%% ----------------------------------------------------------------------------
\section{Extended Algorithm W}
\label{sec:extended-algorithm-w}

Algorithm W is extended to infer effect annotations and handle row unification.

\subsection{Extended Type Inference State}
\label{subsec:extended-infer-state}

\begin{definition}[Extended Inference Monad]
\label{def:extended-infer-monad}
The inference monad tracks:
\begin{enumerate}
  \item Fresh type variable supply
  \item Fresh row variable supply
  \item Fresh ordinal variable supply
  \item Type substitution $\Subst_\tau$
  \item Row substitution $\Subst_\rho$
  \item Ordinal constraints $C_\kappa$
  \item World constraints $C_W$
\end{enumerate}
\end{definition}

\begin{definition}[Extended Algorithm W]
\label{def:extended-alg-w}
$\AlgW_{7.3}(\Gamma, e) = (\Subst, \tau, \varepsilon)$ where:

\textbf{Variable case:}
\begin{align}
\AlgW_{7.3}(\Gamma, x) &= ([], \Inst(\Gamma(x)), \{\})
\end{align}

\textbf{Abstraction case:}
\begin{align}
\AlgW_{7.3}(\Gamma, \lambda x.\, e) &= \text{let } (\Subst, \tau, \varepsilon) = \AlgW_{7.3}(\Gamma[x \mapsto \alpha], e) \\
    &\quad\; \text{in } (\Subst, \Subst(\alpha) \xrightarrow{\varepsilon} \tau, \{\})
\end{align}

\textbf{Application case:}
\begin{align}
\AlgW_{7.3}(\Gamma, e_1\; e_2) &= \text{let } (\Subst_1, \tau_1, \varepsilon_1) = \AlgW_{7.3}(\Gamma, e_1) \\
    &\quad\; (\Subst_2, \tau_2, \varepsilon_2) = \AlgW_{7.3}(\Subst_1(\Gamma), e_2) \\
    &\quad\; \Subst_3 = \Unify(\Subst_2(\tau_1), \tau_2 \xrightarrow{\rho} \beta) \\
    &\quad\; \Subst_4 = \UnifyRow(\Subst_3(\varepsilon_1), \Subst_3(\varepsilon_2), \Subst_3(\rho)) \\
    &\quad\; \text{in } (\Subst_4 \circ \Subst_3 \circ \Subst_2 \circ \Subst_1, \Subst_4(\beta), \Subst_4(\rho))
\end{align}

\textbf{Perform case:}
\begin{align}
\AlgW_{7.3}(\Gamma, \mathsf{perform}\; \mathit{op}(\bar{e})) &= \text{let } E = \mathsf{effectOf}(\mathit{op}) \\
    &\quad\; (\bar{\tau}, \tau_r) = \mathsf{opType}(\mathit{op}) \\
    &\quad\; \text{infer each } e_i, \text{ unify with } \tau_i \\
    &\quad\; \text{in } (\Subst, \tau_r, \{E \mid \rho\})
\end{align}

\textbf{Handle case:}
\begin{align}
\AlgW_{7.3}(\Gamma, \mathsf{handle}\; e\; \mathsf{with}\; h) &= \text{let } (\Subst_1, \tau, \{E \mid \varepsilon\}) = \AlgW_{7.3}(\Gamma, e) \\
    &\quad\; \text{check } h \text{ handles } E \text{ with output type } \tau' \\
    &\quad\; \text{in } (\Subst, \tau', \varepsilon)
\end{align}
\end{definition}

\subsection{Principal Types with Effects}
\label{subsec:principal-effects}

\begin{theorem}[Principal Effect Types]
\label{thm:principal-effect}
For any well-typed expression $e$ under context $\Gamma$, Algorithm $\AlgW_{7.3}$ computes a principal type $(\sigma, \varepsilon)$ such that:
\begin{enumerate}
  \item $\Gamma \vdash e : \sigma \mathbin{!} \varepsilon$ (soundness)
  \item For any $\sigma', \varepsilon'$ with $\Gamma \vdash e : \sigma' \mathbin{!} \varepsilon'$, there exists a substitution $\theta$ with $\theta(\sigma) = \sigma'$ and $\theta(\varepsilon) <: \varepsilon'$ (principality)
\end{enumerate}
\end{theorem}

%% ----------------------------------------------------------------------------
\section{Type Soundness}
\label{sec:type-soundness}

\begin{theorem}[Type Soundness for TKS v7.3]
\label{thm:type-soundness}
The v7.3 type system is sound with respect to the operational semantics (Chapter~\ref{ch:effect-handlers}):
\begin{enumerate}
  \item \textbf{Progress}: If $\varnothing \vdash e : \tau \mathbin{!} \varepsilon$ and $e$ is not a value, then either:
    \begin{itemize}
      \item $e \to e'$ for some $e'$ (reduction step), or
      \item $e = \mathcal{E}[\mathsf{perform}\; \mathit{op}(\bar{v})]$ where $\mathit{op} \in E$ and $E \in \varepsilon$ (effect invocation)
    \end{itemize}
  \item \textbf{Preservation}: If $\Gamma \vdash e : \tau \mathbin{!} \varepsilon$ and $e \to e'$, then $\Gamma \vdash e' : \tau \mathbin{!} \varepsilon$.
  \item \textbf{Effect Safety}: If $\varnothing \vdash e : \tau \mathbin{!} \{\}$ (pure type), then $e$ evaluates to a value without invoking any effects.
\end{enumerate}
\end{theorem}

\begin{proof}[Proof Sketch]
\textbf{Progress} follows by induction on the typing derivation. The key cases are:
\begin{itemize}
  \item \textsc{T-Perform}: The expression is an effect invocation, covered by the second disjunct
  \item \textsc{T-Handle}: By IH on the handled expression, plus handler clause matching
\end{itemize}

\textbf{Preservation} follows by induction on the reduction relation:
\begin{itemize}
  \item Effect handling reductions preserve types via handler clause typing
  \item Continuation capture and invocation preserve types by the continuation typing rule
\end{itemize}

\textbf{Effect Safety} follows from Progress: a pure expression cannot invoke effects, so it must reduce until reaching a value.
\end{proof}

\begin{theorem}[v7.2 Embedding Soundness]
\label{thm:v72-embedding}
The embedding of v7.2 types into v7.3 via $\iota(\tau_{7.2}) = \tau_{7.2} \mathbin{!} \{\}$ preserves typing:
\[
\Gamma_{7.2} \vdash_{7.2} e : \tau \quad \Longrightarrow \quad \iota(\Gamma_{7.2}) \vdash_{7.3} e : \iota(\tau) \mathbin{!} \{\}
\]
Moreover, $\mathsf{RPM}[\tau]$ from v7.2 corresponds to $\tau \mathbin{!} \{\mathsf{RPMEffect}\}$ in v7.3.
\end{theorem}

%% ----------------------------------------------------------------------------
%% HANDOFF NOTE
%% ----------------------------------------------------------------------------
\begin{tcolorbox}[colback=green!5!white,colframe=green!50!black,title=\textbf{Handoff: COMPILER-TASK-02 Complete},breakable]
\textbf{Completed in this section:}
\begin{itemize}
  \item Effect type syntax with row annotations $\tau \mathbin{!} \varepsilon$
  \item Row polymorphism with row variables and quantification
  \item Row unification algorithm with extension handling
  \item Ordinal-indexed types $\mathsf{Frac}^\kappa$, $\mathsf{OrdBounded}[\kappa]$
  \item Bounded ordinal quantification $\forall (\xi < \kappa).\, \tau$
  \item World-stratified types with ACBE coercion typing
  \item Effect row subtyping rules
  \item Typing rules for perform, handle, transfinite fractals, quantum operations
  \item Extended Algorithm W for effect inference
  \item Type soundness theorem (Progress, Preservation, Effect Safety)
  \item v7.2 embedding soundness
\end{itemize}

\textbf{Next tasks (COMPILER-TASK-03):}
\begin{itemize}
  \item Chapter 3: Algebraic Effect Handler System
    \begin{itemize}
      \item Effect handler theory (free monads, delimited continuations)
      \item RPM effect handlers (acquire, check, fail)
      \item Quantum effect handlers (measure, superpose, entangle)
      \item Handler composition and ordering
      \item Operational semantics of handlers
    \end{itemize}
\end{itemize}
\end{tcolorbox}

%% ============================================================================
%% CHAPTER 3: EFFECT HANDLER SYSTEM
%% ============================================================================
\chapter{Algebraic Effect Handlers}
\label{ch:effect-handlers}

\begin{tcolorbox}[colback=blue!5!white,colframe=blue!75!black,title=\vnewthree]
This chapter develops the full algebraic effect handler system for TKS, enabling modular composition of RPM, quantum, and other effects.
\end{tcolorbox}

%% ----------------------------------------------------------------------------
\section{Effect Handler Theory}
\label{sec:effect-handler-theory}

This section develops the theoretical foundations for TKS effect handlers, drawing on free monads, delimited continuations, and the Plotkin--Pretnar framework for algebraic effects.

\subsection{Algebraic Effects and Operations}
\label{subsec:algebraic-effects}

\begin{definition}[Effect Signature]
\label{def:effect-signature}
An \textbf{effect signature} $\Sigma_E$ for effect $E$ is a set of operation symbols with arities:
\[
\Sigma_E = \{\mathit{op}_1 : A_1 \to B_1, \ldots, \mathit{op}_n : A_n \to B_n\}
\]
where $A_i$ is the \textbf{parameter type} and $B_i$ is the \textbf{return type} of operation $\mathit{op}_i$.
\end{definition}

\begin{example}[TKS Effect Signatures]
\label{ex:tks-effect-sigs}
\begin{align}
\Sigma_{\mathsf{RPMEffect}} &= \left\{
  \begin{array}{l}
  \mathsf{acquire} : \mathsf{Element}[W] \to \mathsf{Unit} \\
  \mathsf{check} : \mathsf{Element}[W] \to \mathsf{Bool} \\
  \mathsf{fail} : \mathsf{Unit} \to \mathsf{Void}
  \end{array}
\right\} \\[2ex]
\Sigma_{\mathsf{QuantumEffect}} &= \left\{
  \begin{array}{l}
  \mathsf{measure} : \mathsf{QState}[\alpha] \to \alpha \\
  \mathsf{superpose} : \mathsf{List}[(\mathbb{C}, \alpha)] \to \mathsf{QState}[\alpha] \\
  \mathsf{entangle} : \mathsf{QState}[\alpha] \times \mathsf{QState}[\beta] \to \mathsf{QState}[\alpha \times \beta]
  \end{array}
\right\} \\[2ex]
\Sigma_{\mathsf{StateEffect}[\sigma]} &= \left\{
  \begin{array}{l}
  \mathsf{get} : \mathsf{Unit} \to \sigma \\
  \mathsf{put} : \sigma \to \mathsf{Unit}
  \end{array}
\right\}
\end{align}
\end{example}

\subsection{Free Monad Interpretation}
\label{subsec:free-monad}

\begin{definition}[Free Monad over Signature]
\label{def:free-monad}
Given an effect signature $\Sigma$, the \textbf{free monad} $\mathsf{Free}[\Sigma, \tau]$ is defined inductively:
\begin{align}
\mathsf{Free}[\Sigma, \tau] ::=&\; \mathsf{Pure}(\tau) \\
    &\mid \mathsf{Op}(\mathit{op}, a, k) \quad \text{where } \mathit{op} : A \to B \in \Sigma, a : A, k : B \to \mathsf{Free}[\Sigma, \tau]
\end{align}
The continuation $k$ represents ``what to do after the operation returns''.
\end{definition}

\begin{definition}[Free Monad Operations]
\label{def:free-monad-ops}
The monad operations for $\mathsf{Free}[\Sigma, -]$:
\begin{align}
\mathsf{return}\; v &= \mathsf{Pure}(v) \\
\mathsf{Pure}(v) \rpmbind f &= f(v) \\
\mathsf{Op}(\mathit{op}, a, k) \rpmbind f &= \mathsf{Op}(\mathit{op}, a, \lambda b.\, k(b) \rpmbind f)
\end{align}
\end{definition}

\begin{proposition}[RPM as Free Monad]
\label{prop:rpm-free}
The v7.2 RPM monad $\RPM[\tau]$ is isomorphic to $\mathsf{Free}[\Sigma_{\mathsf{RPMEffect}}, \tau]$ modulo handler interpretation:
\[
\RPM[\tau] \cong \mathsf{Free}[\Sigma_{\mathsf{RPMEffect}}, \tau]
\]
This isomorphism preserves $\rpmret$ and $\rpmbind$.
\end{proposition}

\subsection{Delimited Continuations}
\label{subsec:delimited-cont}

Effect handlers are implemented using delimited continuations, which capture ``the rest of the computation up to a delimiter''.

\begin{definition}[Delimited Continuation Operators]
\label{def:delimited-ops}
We use Felleisen-style control operators:
\begin{align}
\mathsf{reset}\; e &: \tau \quad \text{(install delimiter)} \\
\mathsf{shift}\; k.\, e &: \tau \quad \text{(capture continuation to delimiter)}
\end{align}
The captured continuation $k : \tau_1 \to \tau_2$ represents the evaluation context up to the nearest enclosing $\mathsf{reset}$.
\end{definition}

\begin{definition}[Handler as Delimited Control]
\label{def:handler-delimited}
A handler $\mathsf{handle}\; e\; \mathsf{with}\; h$ is elaborated to delimited control:
\[
\mathsf{handle}\; e\; \mathsf{with}\; h \quad \leadsto \quad \mathsf{reset}\; (e[\mathsf{perform}\; \mathit{op}(a) \mapsto \mathsf{shift}\; k.\, h_{\mathit{op}}(a, k)])
\]
Each $\mathsf{perform}$ is translated to a $\mathsf{shift}$ that captures the continuation and invokes the appropriate handler clause.
\end{definition}

\begin{definition}[Continuation Answer Type]
\label{def:answer-type}
For a handler with output type $\tau_{\mathrm{out}}$, the captured continuation has type:
\[
k : \tau_{\mathrm{ret}} \to \tau_{\mathrm{out}} \mathbin{!} \varepsilon
\]
where $\tau_{\mathrm{ret}}$ is the return type of the handled operation and $\varepsilon$ is the residual effect row.
\end{definition}

\subsection{Handler Syntax and Structure}
\label{subsec:handler-structure}

\begin{definition}[Handler Definition]
\label{def:handler-def}
A \textbf{handler} for effect $E$ with input type $\tau_{\mathrm{in}}$ and output type $\tau_{\mathrm{out}}$ consists of:
\begin{enumerate}
  \item A \textbf{return clause}: $\mathsf{return}\; x \to e_{\mathrm{ret}}$, where $x : \tau_{\mathrm{in}}$ and $e_{\mathrm{ret}} : \tau_{\mathrm{out}}$
  \item An \textbf{operation clause} for each $\mathit{op} \in \Sigma_E$: $\mathit{op}(a)\; k \to e_{\mathit{op}}$, where $a : A_{\mathit{op}}$, $k : B_{\mathit{op}} \to \tau_{\mathrm{out}}$
\end{enumerate}
\end{definition}

\begin{definition}[Handler Expression Syntax]
\label{def:handler-expr-syntax}
\begin{verbatim}
handle e with {
  return x -> e_ret
  op1(a) k -> e_op1
  op2(a) k -> e_op2
  ...
}
\end{verbatim}
\end{definition}

%% ----------------------------------------------------------------------------
\section{RPM Effect Handlers}
\label{sec:rpm-effect-handlers}

This section defines the concrete handlers for RPM operations, connecting algebraic effects to the v7.2 RPM monad semantics.

\subsection{RPM Effect Operations}
\label{subsec:rpm-effect-ops}

\begin{definition}[RPM Effect Operations (Detailed)]
\label{def:rpm-ops-detailed}
The $\mathsf{RPMEffect}$ signature operations:
\begin{enumerate}
  \item $\mathsf{acquire}(e : \mathsf{Element}[W])$: Attempt to acquire element $e$ as a resource. Succeeds if $e$'s prerequisites are satisfied.

  \item $\mathsf{check}(e : \mathsf{Element}[W])$: Check if element $e$ is currently acquired (returns $\mathsf{Bool}$).

  \item $\mathsf{fail}()$: Abort the RPM computation. Does not return (return type $\mathsf{Void}$).
\end{enumerate}
\end{definition}

\begin{definition}[RPM State]
\label{def:rpm-state}
The RPM handler maintains internal state:
\[
\mathsf{RPMState} = \mathcal{P}(\mathsf{Element}[\mathsf{Any}]) \times \mathsf{Goal} \times \mathsf{PrereqMap}
\]
where:
\begin{itemize}
  \item $\mathcal{P}(\mathsf{Element}[\mathsf{Any}])$ is the set of acquired elements
  \item $\mathsf{Goal}$ is the target goal element(s)
  \item $\mathsf{PrereqMap} : \mathsf{Element}[\mathsf{Any}] \to \mathcal{P}(\mathsf{Element}[\mathsf{Any}])$ maps elements to their prerequisites
\end{itemize}
\end{definition}

\subsection{RPM Handler Definition}
\label{subsec:rpm-handler-def}

\begin{definition}[Standard RPM Handler]
\label{def:rpm-handler}
The standard RPM handler with state threading:
\begin{align}
&\mathsf{handleRPM} : (\tau \mathbin{!} \{\mathsf{RPMEffect} \mid \varepsilon\}) \to \mathsf{RPMState} \to (\mathsf{Option}[\tau] \times \mathsf{RPMState}) \mathbin{!} \varepsilon \\[2ex]
&\mathsf{handleRPM} = \mathsf{handler}\; \{ \\
&\quad \mathsf{return}\; v\; s \to (\mathsf{Some}(v), s) \\[1ex]
&\quad \mathsf{acquire}(e)\; k\; s \to \\
&\qquad \mathsf{let}\; (A, g, P) = s\; \mathsf{in} \\
&\qquad \mathsf{if}\; P(e) \subseteq A \\
&\qquad \mathsf{then}\; k(())(A \cup \{e\}, g, P) \\
&\qquad \mathsf{else}\; (\mathsf{None}, s) \\[1ex]
&\quad \mathsf{check}(e)\; k\; s \to \\
&\qquad \mathsf{let}\; (A, g, P) = s\; \mathsf{in}\; k(e \in A)(s) \\[1ex]
&\quad \mathsf{fail}()\; k\; s \to (\mathsf{None}, s) \\
&\}
\end{align}
\end{definition}

\begin{remark}[State Threading Pattern]
\label{rem:state-threading}
The RPM handler uses the \textbf{state-passing} pattern: each handler clause receives the current state $s$ and passes (possibly modified) state to the continuation $k$. This is equivalent to combining $\mathsf{StateEffect}[\mathsf{RPMState}]$ with the core RPM operations.
\end{remark}

\subsection{Connection to v7.2 RPM Monad}
\label{subsec:rpm-v72-connection}

\begin{theorem}[RPM Handler Equivalence]
\label{thm:rpm-handler-equiv}
For any v7.2 RPM computation $m : \RPM[\tau]$ and corresponding v7.3 effectful computation $m' : \tau \mathbin{!} \{\mathsf{RPMEffect}\}$:
\[
\mathsf{runRPM}_{7.2}(m, s_0) = \pi_1(\mathsf{handleRPM}(m', s_0))
\]
where $\pi_1$ projects the result value and $\mathsf{runRPM}_{7.2}$ is the v7.2 monad runner.
\end{theorem}

\begin{proof}[Proof Sketch]
By structural induction on $m$:
\begin{itemize}
  \item \textbf{Case $\rpmret\; v$}: Both return $\mathsf{Some}(v)$ with unchanged state.
  \item \textbf{Case $\mathsf{acquire}(e) \rpmbind f$}: In v7.2, this checks prerequisites and either continues with $f(())$ or fails. The handler clause for $\mathsf{acquire}$ performs exactly this check, invoking $k(())$ (which continues with $f$) on success.
  \item \textbf{Case $\mathsf{fail}$}: Both immediately return $\mathsf{None}$.
\end{itemize}
\end{proof}

\subsection{Deep vs Shallow RPM Handlers}
\label{subsec:deep-shallow}

\begin{definition}[Deep Handler]
\label{def:deep-handler}
A \textbf{deep handler} automatically reinstalls itself when invoking the continuation:
\[
k(v) \equiv \mathsf{handle}\; (k_{\mathrm{raw}}(v))\; \mathsf{with}\; \mathsf{this\_handler}
\]
The standard RPM handler is deep: subsequent $\mathsf{acquire}$ calls within the continuation are handled.
\end{definition}

\begin{definition}[Shallow Handler]
\label{def:shallow-handler}
A \textbf{shallow handler} invokes the raw continuation without reinstalling:
\[
k(v) \equiv k_{\mathrm{raw}}(v)
\]
Shallow handlers are useful for one-shot effect handling or explicit handler management.
\end{definition}

%% ----------------------------------------------------------------------------
\section{Quantum Effect Handlers}
\label{sec:quantum-effect-handlers}

Quantum effects require specialized handling due to superposition, entanglement, and measurement collapse.

\subsection{Quantum Effect Operations}
\label{subsec:quantum-effect-ops}

\begin{definition}[Quantum Effect Operations (Detailed)]
\label{def:quantum-ops-detailed}
The $\mathsf{QuantumEffect}$ signature:
\begin{enumerate}
  \item $\mathsf{superpose}(\{(c_i, v_i)\}_i)$: Create a quantum superposition $\sum_i c_i |v_i\rangle$ where $\sum_i |c_i|^2 = 1$.

  \item $\mathsf{measure}(q : \mathsf{QState}[\tau])$: Measure quantum state $q$, collapsing superposition and returning a classical value of type $\tau$.

  \item $\mathsf{entangle}(q_1, q_2)$: Create an entangled state from two quantum states.
\end{enumerate}
\end{definition}

\begin{definition}[Quantum State Representation]
\label{def:quantum-state-rep}
The quantum handler maintains state as a \textbf{density matrix} or \textbf{state vector}:
\[
\mathsf{QHState}[\tau] = \sum_{i} c_i \cdot |v_i, k_i\rangle
\]
where each $|v_i, k_i\rangle$ represents a branch with classical value $v_i$ and suspended continuation $k_i$.
\end{definition}

\subsection{Quantum Handler Definition}
\label{subsec:quantum-handler-def}

\begin{definition}[Quantum Handler with Superposition Semantics]
\label{def:quantum-handler}
The quantum handler interprets computations over superpositions:
\begin{align}
&\mathsf{handleQuantum} : (\tau \mathbin{!} \{\mathsf{QuantumEffect} \mid \varepsilon\}) \to \mathsf{QState}[\tau] \mathbin{!} \varepsilon \\[2ex]
&\mathsf{handleQuantum} = \mathsf{handler}\; \{ \\
&\quad \mathsf{return}\; v \to |v\rangle \\[1ex]
&\quad \mathsf{superpose}(\{(c_i, v_i)\}_i)\; k \to \\
&\qquad \sum_{i} c_i \cdot k(|v_i\rangle) \\[1ex]
&\quad \mathsf{measure}(q)\; k \to \\
&\qquad \mathsf{collapse}(q, k) \\[1ex]
&\quad \mathsf{entangle}(q_1, q_2)\; k \to \\
&\qquad k(q_1 \otimes q_2) \\
&\}
\end{align}
\end{definition}

\subsection{Measurement and Collapse}
\label{subsec:measurement-collapse}

\begin{definition}[Measurement Collapse Semantics]
\label{def:collapse-semantics}
The $\mathsf{collapse}$ operation for measurement:
\[
\mathsf{collapse}\left(\sum_i c_i |v_i\rangle, k\right) = \mathsf{Dist}\left[\{(|c_i|^2, k(v_i))\}_i\right]
\]
where $\mathsf{Dist}$ represents a probability distribution over continuations. The runtime selects one branch according to $|c_i|^2$ probabilities.
\end{definition}

\begin{definition}[Deferred Measurement]
\label{def:deferred-measurement}
An alternative handler defers measurement until the end:
\begin{align}
&\mathsf{handleQuantumDeferred} = \mathsf{handler}\; \{ \\
&\quad \mathsf{return}\; v \to |v\rangle \\
&\quad \mathsf{measure}(q)\; k \to q \rpmbind_Q (\lambda v.\, k(v)) \\
&\quad \ldots \\
&\}
\end{align}
where $\rpmbind_Q$ is monadic bind over quantum states, maintaining superposition.
\end{definition}

\begin{proposition}[Measurement Deference Equivalence]
\label{prop:measurement-defer}
For computations without observable side effects between measurements:
\[
\mathsf{handleQuantum}(e) \equiv_{\mathrm{dist}} \mathsf{handleQuantumDeferred}(e)
\]
where $\equiv_{\mathrm{dist}}$ denotes equivalence of probability distributions over final results.
\end{proposition}

\subsection{Entanglement Handling}
\label{subsec:entanglement-handling}

\begin{definition}[Entanglement Semantics]
\label{def:entangle-semantics}
For entangled states, measurement of one component affects the other:
\[
\mathsf{measure}\left(\sum_{i,j} c_{ij} |v_i, w_j\rangle\right) = \left(v_k, \frac{\sum_j c_{kj} |w_j\rangle}{\|\sum_j c_{kj} |w_j\rangle\|}\right)
\]
with probability $\sum_j |c_{kj}|^2$ for outcome $v_k$. The entangled partner collapses to the conditional state.
\end{definition}

%% ----------------------------------------------------------------------------
\section{Handler Composition}
\label{sec:handler-composition}

When computations use multiple effects, handlers must be composed. This section addresses handler nesting, ordering, and commutativity.

\subsection{Handler Nesting}
\label{subsec:handler-nesting}

\begin{definition}[Nested Handlers]
\label{def:nested-handlers}
For a computation $e : \tau \mathbin{!} \{E_1, E_2\}$ with two effects, handlers are nested:
\[
\mathsf{handle}\; (\mathsf{handle}\; e\; \mathsf{with}\; h_2)\; \mathsf{with}\; h_1
\]
The inner handler $h_2$ handles $E_2$, producing $\tau' \mathbin{!} \{E_1\}$. The outer handler $h_1$ then handles $E_1$.
\end{definition}

\begin{definition}[Handler Ordering Convention]
\label{def:handler-ordering}
In TKS, we adopt the convention that handlers are listed inside-out:
\[
\mathsf{handle}\; e\; \mathsf{with}\; [h_1, h_2, h_3]
\]
desugars to:
\[
\mathsf{handle}\; (\mathsf{handle}\; (\mathsf{handle}\; e\; \mathsf{with}\; h_3)\; \mathsf{with}\; h_2)\; \mathsf{with}\; h_1
\]
Effect $E_3$ is handled first (innermost), then $E_2$, then $E_1$ (outermost).
\end{definition}

\subsection{Effect Commutativity}
\label{subsec:effect-commutativity}

\begin{definition}[Commutative Effects]
\label{def:commutative-effects}
Effects $E_1$ and $E_2$ \textbf{commute} if for all handlers $h_1, h_2$ and computations $e$:
\[
\mathsf{handle}\; (\mathsf{handle}\; e\; \mathsf{with}\; h_2)\; \mathsf{with}\; h_1 \equiv \mathsf{handle}\; (\mathsf{handle}\; e\; \mathsf{with}\; h_1)\; \mathsf{with}\; h_2
\]
\end{definition}

\begin{proposition}[TKS Effect Commutativity]
\label{prop:tks-commutativity}
In TKS:
\begin{enumerate}
  \item $\mathsf{StateEffect}[\sigma_1]$ and $\mathsf{StateEffect}[\sigma_2]$ commute when $\sigma_1 \neq \sigma_2$ (independent state cells).
  \item $\mathsf{RPMEffect}$ and $\mathsf{QuantumEffect}$ do \emph{not} commute in general (RPM state may depend on quantum measurement outcomes).
  \item $\mathsf{IOEffect}$ does not commute with most effects due to observable ordering.
\end{enumerate}
\end{proposition}

\subsection{Handler Combinators}
\label{subsec:handler-combinators}

\begin{definition}[Handler Product]
\label{def:handler-product}
For commutative effects, the \textbf{handler product} $h_1 \otimes h_2$ handles both simultaneously:
\[
(h_1 \otimes h_2) : \mathsf{Handler}[E_1 \uplus E_2, \tau, \tau']
\]
where $E_1 \uplus E_2$ is the disjoint union of effect operations.
\end{definition}

\begin{definition}[Handler Composition]
\label{def:handler-comp}
The \textbf{handler composition} $h_1 \cdot h_2$ sequences handlers:
\[
(h_1 \cdot h_2)(e) = h_1(h_2(e))
\]
This is the standard nesting with $h_2$ inner and $h_1$ outer.
\end{definition}

%% ----------------------------------------------------------------------------
\section{Operational Semantics of Handlers}
\label{sec:handler-semantics}

This section provides the small-step operational semantics for effect handlers, defining how $\mathsf{handle}$ and $\mathsf{perform}$ reduce.

\subsection{Values and Evaluation Contexts}
\label{subsec:values-contexts}

\begin{definition}[Values (Extended)]
\label{def:values-extended}
Values in the effectful calculus:
\begin{align}
v ::=&\; n \mid b \mid () \mid \lambda x.\, e \quad \text{(standard values)} \\
    &\mid \langle v_1, v_2 \rangle \mid \mathsf{inl}(v) \mid \mathsf{inr}(v) \quad \text{(data values)} \\
    &\mid A_i \mid B_i \mid C_i \mid D_i \quad \text{(element values)} \\
    &\mid |v\rangle \mid \sum_i c_i |v_i\rangle \quad \text{(quantum values)} \\
    &\mid \mathsf{handler}\{\ldots\} \quad \text{(handler values)}
\end{align}
\end{definition}

\begin{definition}[Evaluation Contexts]
\label{def:eval-contexts}
Evaluation contexts $\mathcal{E}$ with effect holes:
\begin{align}
\mathcal{E} ::=&\; [\cdot] \quad \text{(hole)} \\
    &\mid \mathcal{E}\; e \mid v\; \mathcal{E} \quad \text{(application)} \\
    &\mid \mathsf{let}\; x = \mathcal{E}\; \mathsf{in}\; e \quad \text{(let binding)} \\
    &\mid \langle \mathcal{E}, e \rangle \mid \langle v, \mathcal{E} \rangle \quad \text{(pairs)} \\
    &\mid \mathsf{perform}\; \mathit{op}(\bar{v}, \mathcal{E}, \bar{e}) \quad \text{(effect arguments)} \\
    &\mid \mathsf{handle}\; \mathcal{E}\; \mathsf{with}\; h \quad \text{(handled computation)}
\end{align}
\end{definition}

\begin{definition}[Handler Contexts]
\label{def:handler-contexts}
A \textbf{handler context} $\mathcal{H}$ is a context of the form:
\[
\mathcal{H} ::= \mathsf{handle}\; \mathcal{E}\; \mathsf{with}\; h
\]
where $\mathcal{E}$ does not contain another $\mathsf{handle}$ for the same effect $E$ that $h$ handles.
\end{definition}

\subsection{Small-Step Reduction Rules}
\label{subsec:small-step-rules}

\begin{definition}[Core Reduction Rules]
\label{def:core-reduction}
Standard $\beta$-reduction and let-reduction:
\[
\infer[\textsc{E-Beta}]
  {(\lambda x.\, e)\; v \to e[v/x]}
  {}
\qquad
\infer[\textsc{E-Let}]
  {\mathsf{let}\; x = v\; \mathsf{in}\; e \to e[v/x]}
  {}
\]
\end{definition}

\begin{definition}[Handler Reduction Rules]
\label{def:handler-reduction}
\[
\infer[\textsc{E-HandleReturn}]
  {\mathsf{handle}\; v\; \mathsf{with}\; h \to e_{\mathrm{ret}}[v/x]}
  {h = \mathsf{handler}\{\mathsf{return}\; x \to e_{\mathrm{ret}}, \ldots\}}
\]

\[
\infer[\textsc{E-HandlePerform}]
  {\mathsf{handle}\; \mathcal{E}[\mathsf{perform}\; \mathit{op}(v)]\; \mathsf{with}\; h \to e_{\mathit{op}}[v/a, (\lambda y.\, \mathsf{handle}\; \mathcal{E}[y]\; \mathsf{with}\; h)/k]}
  {
    h = \mathsf{handler}\{\ldots, \mathit{op}(a)\; k \to e_{\mathit{op}}, \ldots\} &
    \mathit{op} \in E &
    h \text{ handles } E
  }
\]

The key insight in \textsc{E-HandlePerform}:
\begin{itemize}
  \item The continuation $k$ is bound to $\lambda y.\, \mathsf{handle}\; \mathcal{E}[y]\; \mathsf{with}\; h$
  \item This captures the context $\mathcal{E}$ up to the handler
  \item The handler $h$ is reinstalled (deep handler semantics)
\end{itemize}
\end{definition}

\begin{definition}[Context Reduction]
\label{def:context-reduction}
\[
\infer[\textsc{E-Context}]
  {\mathcal{E}[e] \to \mathcal{E}[e']}
  {e \to e'}
\]
\end{definition}

\subsection{RPM Handler Reduction Examples}
\label{subsec:rpm-reduction-ex}

\begin{example}[RPM Acquire Reduction]
\label{ex:rpm-acquire-reduction}
Consider:
\[
\mathsf{handle}\; (\mathsf{perform}\; \mathsf{acquire}(A1); \mathsf{perform}\; \mathsf{check}(A1))\; \mathsf{with}\; \mathsf{handleRPM}(s_0)
\]
Reduction sequence (assuming $A1$ has no prerequisites):
\begin{align}
&\to e_{\mathsf{acquire}}[A1/a, k_1/k](s_0) \\
&\quad \text{where } k_1 = \lambda ().\, \mathsf{handle}\; (\mathsf{perform}\; \mathsf{check}(A1))\; \mathsf{with}\; \mathsf{handleRPM} \\
&\to k_1(())(s_0 \cup \{A1\}) \\
&\to \mathsf{handle}\; (\mathsf{perform}\; \mathsf{check}(A1))\; \mathsf{with}\; \mathsf{handleRPM}(s_0 \cup \{A1\}) \\
&\to e_{\mathsf{check}}[A1/a, k_2/k](s_0 \cup \{A1\}) \\
&\to k_2(\mathsf{true})(s_0 \cup \{A1\}) \\
&\to \mathsf{handle}\; \mathsf{true}\; \mathsf{with}\; \mathsf{handleRPM}(s_0 \cup \{A1\}) \\
&\to (\mathsf{Some}(\mathsf{true}), s_0 \cup \{A1\})
\end{align}
\end{example}

\subsection{Quantum Handler Reduction}
\label{subsec:quantum-reduction}

\begin{example}[Superposition and Measurement]
\label{ex:quantum-reduction}
\begin{align}
&\mathsf{handle}\; (\mathsf{let}\; q = \mathsf{perform}\; \mathsf{superpose}[(0.5, 0), (0.5, 1)]\; \mathsf{in}\; \mathsf{perform}\; \mathsf{measure}(q))\; \mathsf{with}\; \mathsf{handleQuantum} \\
&\to \mathsf{handle}\; (\mathsf{let}\; q = \frac{1}{\sqrt{2}}|0\rangle + \frac{1}{\sqrt{2}}|1\rangle\; \mathsf{in}\; \mathsf{perform}\; \mathsf{measure}(q))\; \mathsf{with}\; \mathsf{handleQuantum} \\
&\to \mathsf{handle}\; (\mathsf{perform}\; \mathsf{measure}(\frac{1}{\sqrt{2}}|0\rangle + \frac{1}{\sqrt{2}}|1\rangle))\; \mathsf{with}\; \mathsf{handleQuantum} \\
&\to \mathsf{Dist}[(0.5, 0), (0.5, 1)]
\end{align}
\end{example}

%% ----------------------------------------------------------------------------
\section{Type Soundness for Handlers}
\label{sec:handler-type-soundness}

\begin{theorem}[Handler Composition Preserves Type Soundness]
\label{thm:handler-composition-sound}
Let $e : \tau \mathbin{!} \{E_1, E_2 \mid \varepsilon\}$ be a well-typed expression. Let $h_1 : \mathsf{Handler}[E_1, \tau_1, \tau_1']$ and $h_2 : \mathsf{Handler}[E_2, \tau, \tau_1]$ be well-typed handlers where:
\begin{itemize}
  \item $h_2$ transforms type $\tau$ to $\tau_1$ while handling $E_2$
  \item $h_1$ transforms type $\tau_1$ to $\tau_1'$ while handling $E_1$
\end{itemize}
Then the composed handling:
\[
\mathsf{handle}\; (\mathsf{handle}\; e\; \mathsf{with}\; h_2)\; \mathsf{with}\; h_1 : \tau_1' \mathbin{!} \varepsilon
\]
satisfies:
\begin{enumerate}
  \item \textbf{Progress}: The expression either reduces, is a value, or performs an effect in $\varepsilon$.
  \item \textbf{Preservation}: Reduction preserves the type $\tau_1' \mathbin{!} \varepsilon$.
\end{enumerate}
\end{theorem}

\begin{proof}
By induction on the reduction derivation.

\textbf{Base cases:}
\begin{itemize}
  \item \textsc{E-HandleReturn} for inner handler: By the return clause typing of $h_2$, if $e$ reduces to value $v : \tau$, then $e_{\mathrm{ret}}^{(2)}[v/x] : \tau_1 \mathbin{!} \{E_1 \mid \varepsilon\}$. The outer handler then handles this.
  \item \textsc{E-HandleReturn} for outer handler: By the return clause typing of $h_1$, the final result has type $\tau_1' \mathbin{!} \varepsilon$.
\end{itemize}

\textbf{Inductive cases:}
\begin{itemize}
  \item \textsc{E-HandlePerform} for $E_2$: The continuation $k$ has type $B_{\mathit{op}} \to \tau_1 \mathbin{!} \{E_1 \mid \varepsilon\}$ by construction. The handler clause body $e_{\mathit{op}}$ produces $\tau_1 \mathbin{!} \{E_1 \mid \varepsilon\}$, which the outer handler then processes.

  \item \textsc{E-HandlePerform} for $E_1$: The inner handling has already eliminated $E_2$. The continuation for $E_1$ handling has the reinstalled outer handler, preserving types.

  \item \textsc{E-Context}: By IH on the redex within the context.
\end{itemize}

The key observation is that handler composition respects the effect row transformation: inner handlers remove their effects before outer handlers see the computation, maintaining the type invariants at each nesting level.
\end{proof}

\begin{corollary}[Multi-Handler Composition]
\label{cor:multi-handler}
For any sequence of handlers $h_1, \ldots, h_n$ handling distinct effects $E_1, \ldots, E_n$, the nested composition:
\[
\mathsf{handle}\; (\cdots (\mathsf{handle}\; e\; \mathsf{with}\; h_n) \cdots)\; \mathsf{with}\; h_1
\]
is type sound, with final type $\tau' \mathbin{!} \varepsilon$ where $\varepsilon$ excludes $E_1, \ldots, E_n$.
\end{corollary}

%% ----------------------------------------------------------------------------
%% HANDOFF NOTE
%% ----------------------------------------------------------------------------
\begin{tcolorbox}[colback=green!5!white,colframe=green!50!black,title=\textbf{Handoff: COMPILER-TASK-03 Complete},breakable]
\textbf{Completed in this chapter:}
\begin{itemize}
  \item Effect handler theory: effect signatures, free monad interpretation
  \item Delimited continuations: $\mathsf{reset}$/$\mathsf{shift}$, handler elaboration
  \item Handler syntax and structure (return clause, operation clauses)
  \item RPM effect handlers: $\mathsf{acquire}$, $\mathsf{check}$, $\mathsf{fail}$ with state threading
  \item Connection to v7.2 RPM monad (Theorem~\ref{thm:rpm-handler-equiv})
  \item Deep vs shallow handlers
  \item Quantum effect handlers: $\mathsf{superpose}$, $\mathsf{measure}$, $\mathsf{entangle}$
  \item Measurement collapse semantics, deferred measurement
  \item Entanglement handling
  \item Handler composition: nesting, ordering, commutativity
  \item Handler combinators: product, composition
  \item Small-step operational semantics: values, evaluation contexts
  \item Reduction rules: \textsc{E-HandleReturn}, \textsc{E-HandlePerform}
  \item Worked examples for RPM and quantum reductions
  \item Theorem: Handler composition preserves type soundness
\end{itemize}

\textbf{Next tasks (COMPILER-TASK-04):}
\begin{itemize}
  \item Chapter 4: Extended Compiler Architecture
    \begin{itemize}
      \item Extended lexer for v7.3 tokens
      \item Extended parser for effect/handler/ordinal syntax
      \item Extended AST node types
      \item Effect-aware IR representation
      \item Optimization passes (dead effect elimination, handler fusion)
    \end{itemize}
\end{itemize}
\end{tcolorbox}

%% ============================================================================
%% CHAPTER 4: EXTENDED COMPILER ARCHITECTURE
%% ============================================================================
\chapter{Extended Compiler Architecture}
\label{ch:extended-compiler}

\begin{tcolorbox}[colback=blue!5!white,colframe=blue!75!black,title=\vnewthree]
This chapter extends the v7.2 compiler pipeline to handle v7.3 constructs.
\end{tcolorbox}

%% ----------------------------------------------------------------------------
\section{Extended Lexer}
\label{sec:extended-lexer}

This section specifies the v7.3 lexer extensions, adding new token types for transfinite constructs, algebraic effects, and quantum operations while maintaining full backwards compatibility with v7.2 programs.

\subsection{Token Set Extensions}
\label{subsec:token-extensions}

\begin{definition}[v7.3 Token Set]
\label{def:token-set-v73}
The v7.3 token set extends v7.2 with new categories:
\begin{align}
\mathsf{Token}_{7.3} ::=&\; \mathsf{Token}_{7.2} \quad \text{(all v7.2 tokens preserved)} \\[1ex]
%% Transfinite tokens
    &\mid \mathsf{ORDINAL}(\alpha) \quad \text{(e.g., } \omega, \omega_1, \varepsilon_0\text{)} \\
    &\mid \mathsf{OMEGA} \mid \mathsf{EPSILON} \mid \mathsf{ALEPH} \\
    &\mid \mathsf{SUP} \mid \mathsf{LIMIT} \mid \mathsf{SUCC} \\[1ex]
%% Effect tokens
    &\mid \mathsf{EFFECT} \mid \mathsf{HANDLER} \mid \mathsf{HANDLE} \mid \mathsf{WITH} \\
    &\mid \mathsf{PERFORM} \mid \mathsf{RESUME} \\
    &\mid \mathsf{EFFECTROW} \quad \text{(e.g., } !\{E_1, E_2\}\text{)} \\[1ex]
%% Quantum tokens
    &\mid \mathsf{QSTATE} \mid \mathsf{SUPERPOSE} \mid \mathsf{MEASURE} \mid \mathsf{ENTANGLE} \\
    &\mid \mathsf{KET}(v) \mid \mathsf{BRA}(v) \quad \text{(Dirac notation)} \\[1ex]
%% New punctuation
    &\mid \mathsf{BANG} \quad \text{(!)} \\
    &\mid \mathsf{LBRACE} \mid \mathsf{RBRACE} \quad \text{(\{\})} \\
    &\mid \mathsf{PIPE} \quad \text{(|)} \\
    &\mid \mathsf{DOUBLEARROW} \quad \text{(=>)}
\end{align}
\end{definition}

\begin{definition}[v7.2 Token Set (Preserved)]
\label{def:token-set-v72-preserved}
For reference, the complete v7.2 token set:
\begin{align}
\mathsf{Token}_{7.2} ::=&\; \mathsf{ELEMENT}(W, n) \mid \mathsf{NOETIC}(k) \mid \mathsf{FOUNDATION}(m, w) \\
    &\mid \mathsf{INT}(n) \mid \mathsf{BOOL}(b) \mid \mathsf{IDENT}(s) \\
    &\mid \mathsf{LAMBDA} \mid \mathsf{LET} \mid \mathsf{IN} \mid \mathsf{IF} \mid \mathsf{THEN} \mid \mathsf{ELSE} \\
    &\mid \mathsf{RETURN} \mid \mathsf{BIND} \mid \mathsf{CHECK} \mid \mathsf{ACQUIRE} \\
    &\mid \mathsf{ACBE} \mid \mathsf{FRACTAL} \\
    &\mid \mathsf{LPAREN} \mid \mathsf{RPAREN} \mid \mathsf{LANGLE} \mid \mathsf{RANGLE} \\
    &\mid \mathsf{ARROW} \mid \mathsf{EQUALS} \mid \mathsf{COLON} \mid \mathsf{COMMA}
\end{align}
\end{definition}

\subsection{Unicode Support}
\label{subsec:unicode-support}

\begin{definition}[Unicode Token Mapping]
\label{def:unicode-mapping}
The v7.3 lexer accepts both ASCII and Unicode representations:

\begin{center}
\begin{tabular}{lll}
\toprule
\textbf{Token} & \textbf{ASCII} & \textbf{Unicode} \\
\midrule
$\mathsf{LAMBDA}$ & \texttt{\textbackslash} & $\lambda$ (U+03BB) \\
$\mathsf{ARROW}$ & \texttt{->} & $\to$ (U+2192) \\
$\mathsf{DOUBLEARROW}$ & \texttt{=>} & $\Rightarrow$ (U+21D2) \\
$\mathsf{BIND}$ & \texttt{>>=} & $\gg\!=$ (U+226B U+003D) \\
$\mathsf{OMEGA}$ & \texttt{omega} & $\omega$ (U+03C9) \\
$\mathsf{EPSILON}$ & \texttt{epsilon} & $\varepsilon$ (U+03B5) \\
$\mathsf{ALEPH}$ & \texttt{aleph} & $\aleph$ (U+2135) \\
$\mathsf{NOETIC}(k)$ & \texttt{nu\_k} & $\nu_k$ (U+03BD + subscript) \\
$\mathsf{FORALL}$ & \texttt{forall} & $\forall$ (U+2200) \\
$\mathsf{EXISTS}$ & \texttt{exists} & $\exists$ (U+2203) \\
$\mathsf{KET}(v)$ & \texttt{|v>} & $|v\rangle$ (U+27E9) \\
$\mathsf{BRA}(v)$ & \texttt{<v|}  & $\langle v|$ (U+27E8) \\
$\mathsf{TENSOR}$ & \texttt{*} (in context) & $\otimes$ (U+2297) \\
\bottomrule
\end{tabular}
\end{center}
\end{definition}

\subsection{Lexer State Machine}
\label{subsec:lexer-state-machine}

\begin{definition}[Extended Lexer DFA States]
\label{def:lexer-dfa}
The v7.3 lexer uses a deterministic finite automaton with states:
\begin{align}
Q_{7.3} = Q_{7.2} \cup \{&\, q_{\mathrm{ord}}, q_{\mathrm{eff}}, q_{\mathrm{ket}}, q_{\mathrm{bra}}, q_{\mathrm{effrow}}, \\
    &\, q_{\mathrm{uni}}, q_{\mathrm{uni2}}, q_{\mathrm{uni3}}, q_{\mathrm{uni4}} \,\}
\end{align}
where:
\begin{itemize}
  \item $q_{\mathrm{ord}}$: Recognizing ordinal literals
  \item $q_{\mathrm{eff}}$: Recognizing effect row syntax
  \item $q_{\mathrm{ket}}, q_{\mathrm{bra}}$: Recognizing Dirac notation
  \item $q_{\mathrm{uni}}, q_{\mathrm{uni2}}, q_{\mathrm{uni3}}, q_{\mathrm{uni4}}$: UTF-8 multibyte decoding states
\end{itemize}
\end{definition}

\begin{definition}[Lexer Transition Function Extensions]
\label{def:lexer-transitions}
Key transition extensions (partial specification):
\begin{align}
\delta(q_0, \mathtt{|})\; &= q_{\mathrm{ket}} \\
\delta(q_{\mathrm{ket}}, [a\text{-}zA\text{-}Z0\text{-}9])\; &= q_{\mathrm{ket}} \\
\delta(q_{\mathrm{ket}}, \mathtt{>})\; &= q_{\mathrm{accept}}[\mathsf{KET}] \\[1ex]
\delta(q_0, \mathtt{!})\; &= q_{\mathrm{eff}} \\
\delta(q_{\mathrm{eff}}, \mathtt{\{})\; &= q_{\mathrm{effrow}} \\[1ex]
\delta(q_0, \omega)\; &= q_{\mathrm{accept}}[\mathsf{OMEGA}] \quad \text{(Unicode direct)} \\
\delta(q_0, \mathtt{o})\; &= q_{\mathrm{id}} \quad \text{(may become \texttt{omega})}
\end{align}
\end{definition}

\begin{definition}[Keyword Recognition]
\label{def:keyword-recognition}
After recognizing an identifier in $q_{\mathrm{id}}$, the lexer checks against the keyword table:
\begin{verbatim}
keywords_v73 = keywords_v72 ++ [
  ("effect",    EFFECT),
  ("handler",   HANDLER),
  ("handle",    HANDLE),
  ("with",      WITH),
  ("perform",   PERFORM),
  ("resume",    RESUME),
  ("omega",     OMEGA),
  ("epsilon",   EPSILON),
  ("aleph",     ALEPH),
  ("sup",       SUP),
  ("limit",     LIMIT),
  ("succ",      SUCC),
  ("superpose", SUPERPOSE),
  ("measure",   MEASURE),
  ("entangle",  ENTANGLE),
  ("qstate",    QSTATE)
]
\end{verbatim}
\end{definition}

\subsection{Lexer Pseudocode}
\label{subsec:lexer-pseudocode}

\begin{definition}[Extended Lexer Algorithm]
\label{def:lexer-algorithm}
\begin{verbatim}
lex_v73 : String -> Result [Token] LexError
lex_v73 input = go input [] where
  go ""     acc = Ok (reverse acc)
  go (c:cs) acc
    | isSpace c           = go cs acc
    | c == '-' && head cs == '-' = go (dropLine cs) acc

    -- v7.3 extensions
    | c == '|' && isAlphaNum (head cs) = lexKet cs acc
    | c == '<' && isAlphaNum (head cs) = lexBra cs acc
    | c == '!' && head cs == '{'       = lexEffectRow cs acc
    | isUnicodeStart c    = lexUnicode (c:cs) acc

    -- v7.2 token recognition (unchanged)
    | c `elem` "ABCD" && isDigit (head cs) = lexElement (c:cs) acc
    | c == 'n' && head cs == 'u'           = lexNoetic (c:cs) acc
    | c == 'F' && isDigit (head cs)        = lexFoundation (c:cs) acc
    | isAlpha c           = lexIdentOrKeyword (c:cs) acc
    | isDigit c           = lexNumber (c:cs) acc
    | otherwise           = lexPunct (c:cs) acc
\end{verbatim}
\end{definition}

\begin{proposition}[Backwards Compatibility]
\label{prop:lex-backwards-compat}
For any v7.2 source string $s$:
\[
\mathsf{lex}_{7.3}(s) = \mathsf{lex}_{7.2}(s)
\]
That is, the v7.3 lexer produces identical token streams for v7.2 programs.
\end{proposition}

\begin{proof}
The v7.3 lexer extends v7.2 by adding new transitions from states not reachable by v7.2 input. All v7.2 keywords remain in the keyword table with identical token types. The new transitions for $|$, $!$, and Unicode characters only activate on input patterns that are lexically invalid in v7.2 (and would have produced errors). Therefore, valid v7.2 input follows identical DFA paths and produces identical tokens.
\end{proof}

%% ----------------------------------------------------------------------------
\section{Extended Parser}
\label{sec:extended-parser}

This section extends the v7.2 recursive descent parser with productions for effects, handlers, and transfinite constructs.

\subsection{Extended Grammar}
\label{subsec:extended-grammar}

\begin{definition}[v7.3 Grammar Extensions]
\label{def:grammar-v73}
The v7.3 grammar extends v7.2 with new productions (additions in \textbf{bold}):
\begin{verbatim}
program  ::= topdecl* expr

topdecl  ::= letdecl | typedecl | effectdecl | handlerdecl

effectdecl ::= 'effect' IDENT '{' opdecl* '}'
opdecl     ::= IDENT ':' type '->' type

handlerdecl ::= 'handler' IDENT ':' handlertype '{'
                  retclause opclause*
                '}'
retclause   ::= 'return' IDENT '->' expr
opclause    ::= IDENT '(' IDENT ')' IDENT '->' expr

expr     ::= letexpr | lamexpr | ifexpr | bindexpr
           | handleexpr | performexpr | ordinalexpr

handleexpr  ::= 'handle' expr 'with' handler
handler     ::= IDENT | '{' retclause opclause* '}'

performexpr ::= 'perform' IDENT '(' expr ')'

ordinalexpr ::= ordprimary ordop ordprimary
              | 'limit' IDENT '<' ordinalexpr '.' expr
              | 'succ' '(' ordinalexpr ')'
ordprimary  ::= ORDINAL | 'omega' | 'epsilon' | 'aleph' SUBSCRIPT
              | '(' ordinalexpr ')'
ordop       ::= '+' | '*' | '^'

-- Quantum expressions
quantumexpr ::= superposexp | measureexp | entangleexp | ketexp
superposexp ::= 'superpose' '[' amplist ']'
amplist     ::= '(' expr ',' expr ')' (',' '(' expr ',' expr ')')*
measureexp  ::= 'measure' '(' expr ')'
entangleexp ::= 'entangle' '(' expr ',' expr ')'
ketexp      ::= '|' expr '>'
braexp      ::= '<' expr '|'

-- Types with effects (extends type production)
type     ::= ... | type '!' effectrow
effectrow ::= '{' effectlist '}'
            | '{' effectlist '|' IDENT '}'   -- row variable
effectlist ::= IDENT (',' IDENT)*
             | (* empty *)
\end{verbatim}
\end{definition}

\subsection{Precedence and Associativity}
\label{subsec:precedence}

\begin{definition}[Operator Precedence Table]
\label{def:precedence-table}
The v7.3 operator precedence (highest to lowest):
\begin{center}
\begin{tabular}{clcl}
\toprule
\textbf{Level} & \textbf{Operators} & \textbf{Assoc.} & \textbf{Description} \\
\midrule
10 & function application & left & $f\; x\; y$ \\
9 & $\nu_k$ (noetic application) & right & $\nu_2(x)$ \\
8 & \texttt{\^{}} (ordinal exp) & right & $\omega^{\omega}$ \\
7 & \texttt{*}, $\otimes$ (ordinal/tensor mult) & left & $\omega \cdot 2$ \\
6 & \texttt{+} (ordinal add) & left & $\omega + n$ \\
5 & $\gg\!=$ (monadic bind) & left & $m \gg\!= f$ \\
4 & \texttt{handle ... with} & prefix & handlers \\
3 & \texttt{if...then...else} & prefix & conditionals \\
2 & $\lambda$, \texttt{let...in} & prefix & abstraction \\
1 & \texttt{!} (effect annotation) & postfix & $\tau \mathbin{!} \varepsilon$ \\
0 & \texttt{->} (function type) & right & $\tau_1 \to \tau_2$ \\
\bottomrule
\end{tabular}
\end{center}
\end{definition}

\begin{definition}[Associativity Rules]
\label{def:associativity}
Notable associativity rules:
\begin{enumerate}
  \item \textbf{Handler nesting}: Right-associative. \texttt{handle e with h1 with h2} parses as \texttt{handle (handle e with h2) with h1}.

  \item \textbf{Effect rows}: The $|$ in effect rows is not an operator; it separates concrete effects from the row variable.

  \item \textbf{Ordinal exponentiation}: Right-associative. $\omega^{\omega^\omega}$ parses as $\omega^{(\omega^\omega)}$.

  \item \textbf{Monadic bind}: Left-associative. $m_1 \gg\!= f \gg\!= g$ parses as $(m_1 \gg\!= f) \gg\!= g$.
\end{enumerate}
\end{definition}

\subsection{Parser Implementation}
\label{subsec:parser-impl}

\begin{definition}[Recursive Descent Parser Extensions]
\label{def:parser-extensions}
Key parsing functions for v7.3 constructs:
\begin{verbatim}
parseHandleExpr :: Parser Expr
parseHandleExpr = do
  keyword "handle"
  body <- parseExpr
  keyword "with"
  h <- parseHandler
  return (HandleExpr body h)

parseHandler :: Parser Handler
parseHandler =
      (IdentHandler <$> parseIdent)
  <|> parseInlineHandler

parseInlineHandler :: Parser Handler
parseInlineHandler = do
  symbol "{"
  ret <- parseReturnClause
  ops <- many parseOpClause
  symbol "}"
  return (InlineHandler ret ops)

parsePerformExpr :: Parser Expr
parsePerformExpr = do
  keyword "perform"
  op <- parseIdent
  symbol "("
  arg <- parseExpr
  symbol ")"
  return (PerformExpr op arg)

parseOrdinalExpr :: Parser Expr
parseOrdinalExpr = buildExpressionParser ordinalOps parseOrdPrimary
  where
    ordinalOps = [ [Infix (symbol "^" >> return OrdExp) AssocRight]
                 , [Infix (symbol "*" >> return OrdMul) AssocLeft]
                 , [Infix (symbol "+" >> return OrdAdd) AssocLeft]
                 ]

parseOrdPrimary :: Parser Expr
parseOrdPrimary =
      (OrdLit <$> parseOrdinalLiteral)
  <|> (keyword "omega" >> return OrdOmega)
  <|> parseLimitExpr
  <|> parseSuccExpr
  <|> parens parseOrdinalExpr
\end{verbatim}
\end{definition}

\subsection{Error Recovery}
\label{subsec:error-recovery}

\begin{definition}[Synchronization Points]
\label{def:sync-points}
The parser uses synchronization points for error recovery:
\begin{enumerate}
  \item \textbf{Declaration level}: On error, skip to next \texttt{let}, \texttt{effect}, \texttt{handler}, or end-of-file.

  \item \textbf{Expression level}: On error in expressions, skip to next $)$, $\}$, \texttt{in}, \texttt{then}, \texttt{else}, \texttt{with}, or semicolon.

  \item \textbf{Handler clause level}: On error in handler clause, skip to next clause keyword or $\}$.
\end{enumerate}
\end{definition}

\begin{definition}[Error Messages]
\label{def:error-messages}
Context-aware error messages for v7.3 constructs:
\begin{verbatim}
parseError :: Token -> ParseState -> String
parseError tok st = case currentContext st of
  InHandler ->
    "Expected 'return' clause or operation clause, found: " ++ show tok
  InEffectRow ->
    "Expected effect name or '|' for row variable, found: " ++ show tok
  InOrdinal ->
    "Expected ordinal (omega, epsilon, or literal), found: " ++ show tok
  InQuantum ->
    "Expected quantum expression or ket, found: " ++ show tok
  _ -> "Unexpected token: " ++ show tok
\end{verbatim}
\end{definition}

%% ----------------------------------------------------------------------------
\section{Extended AST}
\label{sec:extended-ast}

This section defines the v7.3 Abstract Syntax Tree, extending v7.2 with nodes for effects, handlers, and transfinite constructs.

\subsection{AST Node Types}
\label{subsec:ast-nodes}

\begin{definition}[v7.3 AST Data Types]
\label{def:ast-types}
The complete v7.3 AST (Haskell-style definition):
\begin{verbatim}
-- Source location for error reporting
data Loc = Loc { file :: String, line :: Int, col :: Int }

-- Top-level declarations
data TopDecl
  = LetDecl Loc Ident (Maybe TypeScheme) Expr
  | TypeDecl Loc Ident [TyVar] Type
  | EffectDecl Loc Ident [OpSig]            -- NEW v7.3
  | HandlerDecl Loc Ident HandlerType HandlerDef  -- NEW v7.3

-- Effect operation signature
data OpSig = OpSig Loc Ident Type Type      -- NEW v7.3

-- Handler definition
data HandlerDef = HandlerDef                -- NEW v7.3
  { hdReturnClause :: (Ident, Expr)
  , hdOpClauses    :: [(Ident, Ident, Ident, Expr)]
  }                  -- op(arg) k -> body

-- Expressions (extending v7.2)
data Expr
  -- v7.2 expressions (preserved)
  = Var Loc Ident
  | Lit Loc Literal
  | Lam Loc Ident (Maybe Type) Expr
  | App Loc Expr Expr
  | Let Loc Ident (Maybe TypeScheme) Expr Expr
  | If Loc Expr Expr Expr
  | Element Loc Wing Int
  | Foundation Loc Int Char
  | Noetic Loc Int Expr
  | Fractal Loc [Int] Expr
  | RPMReturn Loc Expr
  | RPMBind Loc Expr Expr
  | RPMCheck Loc Expr
  | RPMAcquire Loc Expr
  | ACBE Loc Expr Expr
  -- v7.3 effect expressions (NEW)
  | Handle Loc Expr Handler
  | Perform Loc Ident Expr
  -- v7.3 ordinal expressions (NEW)
  | OrdLit Loc Ordinal
  | OrdOmega Loc
  | OrdEpsilon Loc Int
  | OrdAleph Loc Int
  | OrdSucc Loc Expr
  | OrdAdd Loc Expr Expr
  | OrdMul Loc Expr Expr
  | OrdExp Loc Expr Expr
  | OrdLimit Loc Ident Expr Expr
  -- v7.3 quantum expressions (NEW)
  | Superpose Loc [(Expr, Expr)]    -- [(amplitude, value)]
  | Measure Loc Expr
  | Entangle Loc Expr Expr
  | Ket Loc Expr
  | Bra Loc Expr
  | BraKet Loc Expr Expr

-- Handler references
data Handler                                -- NEW v7.3
  = HandlerRef Loc Ident
  | InlineHandler Loc HandlerDef

-- Types (extending v7.2)
data Type
  -- v7.2 types (preserved)
  = TyVar TyVar
  | TyCon TyCon [Type]
  | TyFun Type Type
  | TyElement Wing
  | TyNoetic Type
  | TyFractal Type
  | TyRPM Type
  -- v7.3 types (NEW)
  | TyEffectful Type EffectRow          -- tau ! {E1, E2 | r}
  | TyHandler EffectRow Type Type       -- Handler[E, tau_in, tau_out]
  | TyOrdinal                           -- Ord
  | TyQState Type                       -- QState[tau]

-- Effect rows
data EffectRow                              -- NEW v7.3
  = EffectRowEmpty
  | EffectRowCons Ident EffectRow
  | EffectRowVar TyVar
\end{verbatim}
\end{definition}

\subsection{AST Annotations}
\label{subsec:ast-annotations}

\begin{definition}[Effect Annotations]
\label{def:effect-annotations}
After type checking, AST nodes carry effect annotations:
\begin{verbatim}
data AnnotatedExpr = AExpr
  { aeExpr   :: Expr
  , aeLoc    :: Loc
  , aeType   :: Type
  , aeEffect :: EffectRow           -- NEW v7.3
  }

-- Effect inference produces annotations
inferEffects :: TypeEnv -> Expr -> (AnnotatedExpr, Constraints)
\end{verbatim}
\end{definition}

\begin{definition}[Source Location Tracking]
\label{def:source-loc}
Every AST node carries source location for error reporting:
\begin{verbatim}
class HasLoc a where
  getLoc :: a -> Loc

instance HasLoc Expr where
  getLoc (Var loc _)      = loc
  getLoc (Handle loc _ _) = loc
  getLoc (Perform loc _ _) = loc
  -- ... etc.

-- Error reporting uses locations
reportError :: Loc -> String -> CompilerError
reportError loc msg = CompilerError
  { errLoc = loc
  , errMsg = msg
  , errContext = getSourceLine (file loc) (line loc)
  }
\end{verbatim}
\end{definition}

\subsection{AST Well-formedness}
\label{subsec:ast-wellformed}

\begin{definition}[AST Well-formedness Conditions]
\label{def:ast-wellformed}
An AST is well-formed if:
\begin{enumerate}
  \item All variable references have corresponding bindings in scope.
  \item All effect references in \texttt{Perform} nodes correspond to declared effects.
  \item All handler references correspond to declared or inline handlers.
  \item Effect rows are well-formed (no duplicate effect names).
  \item Ordinal expressions use only ordinal-typed subexpressions where required.
\end{enumerate}
\end{definition}

\begin{verbatim}
checkWellFormed :: Program -> Either [WFError] ()
checkWellFormed prog = do
  checkScoping prog
  checkEffectRefs prog
  checkHandlerRefs prog
  checkEffectRows prog
  checkOrdinalTypes prog
\end{verbatim}

%% ----------------------------------------------------------------------------
\section{Effect-Aware Intermediate Representation}
\label{sec:effect-aware-ir}

This section defines the TKS Intermediate Representation (IR), which makes effects explicit and prepares for handler compilation via CPS transformation.

\subsection{IR Design}
\label{subsec:ir-design}

\begin{definition}[TKS IR Syntax]
\label{def:ir-syntax}
The TKS IR uses A-Normal Form (ANF) with explicit effect tracking:
\begin{verbatim}
-- IR values (trivial expressions)
data IRVal
  = IRVar Ident
  | IRLit Literal
  | IRLam Ident IRTerm
  | IRElement Wing Int
  | IRNoetic Int
  | IROrdinal Ordinal
  | IRKet IRVal

-- IR terms (complex expressions)
data IRTerm
  -- Core terms
  = IRReturn IRVal                      -- return v (pure)
  | IRLet Ident IRComp IRTerm           -- let x = comp in term
  | IRApp IRVal IRVal                   -- function application
  | IRIf IRVal IRTerm IRTerm            -- conditional

  -- Effect terms (NEW v7.3)
  | IRPerform EffectName OpName IRVal   -- perform E.op(v)
  | IRHandle IRTerm IRHandler           -- handle term with handler

  -- RPM terms (v7.2, now as effects)
  | IRRPMAcquire IRVal                  -- acquire element
  | IRRPMCheck IRVal                    -- check element
  | IRRPMFail                           -- fail computation

  -- Ordinal terms (NEW v7.3)
  | IROrdLimit Ident IRVal IRTerm       -- limit i < alpha. body
  | IROrdSucc IRVal
  | IROrdOp OrdOp IRVal IRVal

  -- Quantum terms (NEW v7.3)
  | IRSuperpose [(IRVal, IRVal)]        -- superpose
  | IRMeasure IRVal                     -- measure
  | IREntangle IRVal IRVal              -- entangle

-- Computations (for ANF)
data IRComp
  = IRPure IRVal                        -- pure value
  | IRBind IRTerm (Ident -> IRTerm)     -- sequencing
  | IRTerm IRTerm                       -- effectful term

-- IR Handlers
data IRHandler = IRHandler
  { irhEffect  :: EffectName
  , irhReturn  :: Ident -> IRTerm
  , irhOps     :: [(OpName, Ident, Ident, IRTerm)]  -- op(a) k -> body
  }

-- Effect annotations on IR
data IRType = IRType
  { irtBase   :: BaseType
  , irtEffect :: EffectRow
  }
\end{verbatim}
\end{definition}

\subsection{ANF Transformation}
\label{subsec:anf-transform}

\begin{definition}[ANF Transformation Rules]
\label{def:anf-rules}
The ANF transformation $\mathcal{A}\llbracket - \rrbracket$ converts AST to IR:
\begin{align}
\mathcal{A}\llbracket x \rrbracket &= \mathsf{IRReturn}\; (\mathsf{IRVar}\; x) \\
\mathcal{A}\llbracket \lambda x.\, e \rrbracket &= \mathsf{IRReturn}\; (\mathsf{IRLam}\; x\; \mathcal{A}\llbracket e \rrbracket) \\
\mathcal{A}\llbracket e_1\; e_2 \rrbracket &= \mathsf{IRLet}\; x_1\; (\mathcal{A}\llbracket e_1 \rrbracket)\; (\mathsf{IRLet}\; x_2\; (\mathcal{A}\llbracket e_2 \rrbracket)\; (\mathsf{IRApp}\; x_1\; x_2)) \\
\mathcal{A}\llbracket \mathsf{let}\; x = e_1\; \mathsf{in}\; e_2 \rrbracket &= \mathsf{IRLet}\; x\; (\mathcal{A}\llbracket e_1 \rrbracket)\; (\mathcal{A}\llbracket e_2 \rrbracket) \\
\mathcal{A}\llbracket \mathsf{perform}\; \mathit{op}(e) \rrbracket &= \mathsf{IRLet}\; x\; (\mathcal{A}\llbracket e \rrbracket)\; (\mathsf{IRPerform}\; E\; \mathit{op}\; x) \\
\mathcal{A}\llbracket \mathsf{handle}\; e\; \mathsf{with}\; h \rrbracket &= \mathsf{IRHandle}\; (\mathcal{A}\llbracket e \rrbracket)\; (\mathcal{H}\llbracket h \rrbracket)
\end{align}
where $x, x_1, x_2$ are fresh variables.
\end{definition}

\begin{definition}[Handler ANF Transformation]
\label{def:handler-anf}
Handler transformation $\mathcal{H}\llbracket - \rrbracket$:
\begin{align}
\mathcal{H}\llbracket \mathsf{handler}\{\mathsf{return}\; x \to e_r, \overline{\mathit{op}(a)\; k \to e_{op}}\} \rrbracket &= \mathsf{IRHandler}\; \{\\
    &\quad \mathsf{irhReturn} = \lambda x.\, \mathcal{A}\llbracket e_r \rrbracket \\
    &\quad \mathsf{irhOps} = [(\mathit{op}, a, k, \mathcal{A}\llbracket e_{op} \rrbracket)]_{\mathit{op}} \\
    &\}
\end{align}
\end{definition}

\subsection{CPS Transformation for Handlers}
\label{subsec:cps-transform}

\begin{definition}[CPS Transformation]
\label{def:cps-transform}
For compilation to the VM, handlers are transformed to continuation-passing style:
\begin{align}
\mathcal{C}\llbracket \mathsf{IRReturn}\; v \rrbracket\; k &= k(v) \\
\mathcal{C}\llbracket \mathsf{IRLet}\; x\; c\; t \rrbracket\; k &= \mathcal{C}\llbracket c \rrbracket\; (\lambda v.\, \mathcal{C}\llbracket t[v/x] \rrbracket\; k) \\
\mathcal{C}\llbracket \mathsf{IRPerform}\; E\; \mathit{op}\; v \rrbracket\; k &= \mathsf{IRPerformCPS}\; E\; \mathit{op}\; v\; k \\
\mathcal{C}\llbracket \mathsf{IRHandle}\; t\; h \rrbracket\; k &= \mathcal{C}\llbracket t \rrbracket\; (\mathsf{wrapHandler}\; h\; k)
\end{align}
where $\mathsf{wrapHandler}$ installs handler clauses around the continuation.
\end{definition}

\begin{definition}[CPS Handler Installation]
\label{def:cps-handler-install}
\begin{verbatim}
wrapHandler :: IRHandler -> (IRVal -> IRTerm) -> (IRVal -> IRTerm)
wrapHandler h outerK = innerK where
  innerK v = case v of
    -- Normal return: use handler's return clause
    IRReturnVal x ->
      let retBody = irhReturn h x
      in cpsTerm retBody outerK

    -- Effect operation: check if handler handles it
    IRPerformVal eff op arg k' ->
      if eff == irhEffect h
      then -- Handle the operation
        let (_, argName, kName, opBody) =
              findOp op (irhOps h)
            resumeK = \v -> wrapHandler h outerK (k' v)
        in opBody[arg/argName, resumeK/kName]
      else -- Pass through to outer handler
        IRPerformCPS eff op arg (\v -> innerK (k' v))
\end{verbatim}
\end{definition}

\subsection{IR Typing Preservation}
\label{subsec:ir-typing}

\begin{theorem}[IR Preserves Typing Judgments]
\label{thm:ir-typing-preservation}
If $\Gamma \vdash e : \tau \mathbin{!} \varepsilon$ in the surface language, then:
\[
\mathcal{A}\llbracket e \rrbracket : \mathsf{IRType}(\tau, \varepsilon)
\]
and for all $\Gamma \vdash e : \tau \mathbin{!} \varepsilon$ where $e$ reduces to $e'$ in the surface semantics:
\[
\mathcal{A}\llbracket e \rrbracket \to^* \mathcal{A}\llbracket e' \rrbracket \text{ in IR reduction}
\]
\end{theorem}

\begin{proof}
By structural induction on the typing derivation.

\textbf{Base cases:}
\begin{itemize}
  \item \textbf{Variables}: $\Gamma \vdash x : \tau$ transforms to $\mathsf{IRReturn}\;(\mathsf{IRVar}\; x)$. The IR typing rule for $\mathsf{IRReturn}$ assigns type $\mathsf{IRType}(\tau, \emptyset)$, which matches since variables are pure.

  \item \textbf{Literals}: Similar to variables; literals are pure.
\end{itemize}

\textbf{Inductive cases:}
\begin{itemize}
  \item \textbf{Application} $\Gamma \vdash e_1\; e_2 : \tau \mathbin{!} \varepsilon$:
  By IH, $\mathcal{A}\llbracket e_1 \rrbracket : \mathsf{IRType}(\tau_1 \to \tau, \varepsilon_1)$ and $\mathcal{A}\llbracket e_2 \rrbracket : \mathsf{IRType}(\tau_1, \varepsilon_2)$.
  The ANF transformation produces:
  \[
  \mathsf{IRLet}\; x_1\; (\mathcal{A}\llbracket e_1 \rrbracket)\; (\mathsf{IRLet}\; x_2\; (\mathcal{A}\llbracket e_2 \rrbracket)\; (\mathsf{IRApp}\; x_1\; x_2))
  \]
  By the IR typing rule for $\mathsf{IRLet}$, this has type $\mathsf{IRType}(\tau, \varepsilon_1 \cup \varepsilon_2)$, which equals $\varepsilon$ by the surface typing rule for application.

  \item \textbf{Perform} $\Gamma \vdash \mathsf{perform}\; \mathit{op}(e) : B \mathbin{!} \{E \mid \varepsilon\}$:
  By IH, $\mathcal{A}\llbracket e \rrbracket : \mathsf{IRType}(A, \varepsilon')$.
  The transformation produces $\mathsf{IRPerform}\; E\; \mathit{op}\; x$ after binding $e$ to $x$.
  The IR typing rule for $\mathsf{IRPerform}$ adds effect $E$ to the row, giving type $\mathsf{IRType}(B, \{E\} \cup \varepsilon')$.

  \item \textbf{Handle} $\Gamma \vdash \mathsf{handle}\; e\; \mathsf{with}\; h : \tau' \mathbin{!} \varepsilon'$:
  By IH, $\mathcal{A}\llbracket e \rrbracket : \mathsf{IRType}(\tau, \{E \mid \varepsilon'\})$.
  The handler $h$ transforms effects of type $E$, removing $E$ from the effect row.
  The resulting $\mathsf{IRHandle}$ has type $\mathsf{IRType}(\tau', \varepsilon')$ by the IR handler typing rule.
\end{itemize}

\textbf{Simulation:}
For each surface reduction step $e \to e'$, there exists a corresponding IR reduction sequence $\mathcal{A}\llbracket e \rrbracket \to^* \mathcal{A}\llbracket e' \rrbracket$. This follows because:
\begin{enumerate}
  \item Beta reduction in the surface corresponds to IR application reduction.
  \item Handler reduction (E-HandleReturn, E-HandlePerform) corresponds to CPS handler execution.
  \item ANF introduces only administrative reductions (let-bindings) that preserve semantics.
\end{enumerate}
\end{proof}

\begin{corollary}[v7.2 Compilation Unchanged]
\label{cor:v72-compilation}
For any v7.2 expression $e$ not using v7.3 constructs:
\[
\mathcal{A}_{7.3}\llbracket e \rrbracket = \mathcal{A}_{7.2}\llbracket e \rrbracket
\]
The v7.3 IR transformation is backwards compatible.
\end{corollary}

%% ----------------------------------------------------------------------------
\section{Optimization Passes}
\label{sec:optimization-passes}

This section describes effect-aware optimization passes for the TKS compiler.

\subsection{Effect-Aware Optimization Framework}
\label{subsec:effect-opt-framework}

\begin{definition}[Effect Purity Analysis]
\label{def:purity-analysis}
An IR term $t$ is \textbf{pure} if its effect row is empty: $\mathsf{effects}(t) = \emptyset$.
\begin{verbatim}
isPure :: IRTerm -> Bool
isPure t = effectRow t == EffectRowEmpty

-- Purity allows reordering, inlining, and elimination
canReorder :: IRTerm -> IRTerm -> Bool
canReorder t1 t2 = isPure t1 || isPure t2 || effectsCommute t1 t2
\end{verbatim}
\end{definition}

\subsection{Dead Effect Elimination}
\label{subsec:dead-effect-elim}

\begin{definition}[Dead Effect Elimination]
\label{def:dead-effect-elim}
An effect operation is \textbf{dead} if:
\begin{enumerate}
  \item Its result is unused, AND
  \item It has no observable side effects beyond its declared effect
\end{enumerate}
\begin{verbatim}
deadEffectElim :: IRTerm -> IRTerm
deadEffectElim term = case term of
  IRLet x (IRTerm (IRPerform eff op v)) body
    | x `notFreeIn` body && isDeadEffect eff op ->
        deadEffectElim body
  IRLet x comp body ->
    IRLet x (deadEffectElim comp) (deadEffectElim body)
  _ -> term

-- Only certain effects can be eliminated
isDeadEffect :: EffectName -> OpName -> Bool
isDeadEffect "State" "get" = True   -- Reading state is dead if unused
isDeadEffect "RPM" "check" = True   -- Checking is dead if unused
isDeadEffect _ _ = False            -- Most effects are not dead
\end{verbatim}
\end{definition}

\subsection{Handler Fusion}
\label{subsec:handler-fusion}

\begin{definition}[Handler Fusion]
\label{def:handler-fusion}
When handlers for commutative effects are nested, they can be fused:
\[
\mathsf{handle}\; (\mathsf{handle}\; e\; \mathsf{with}\; h_2)\; \mathsf{with}\; h_1 \quad \Rightarrow \quad \mathsf{handle}\; e\; \mathsf{with}\; (h_1 \otimes h_2)
\]
if $\mathsf{effect}(h_1)$ and $\mathsf{effect}(h_2)$ commute.
\begin{verbatim}
handlerFusion :: IRTerm -> IRTerm
handlerFusion term = case term of
  IRHandle (IRHandle inner h2) h1
    | effectsCommute (irhEffect h1) (irhEffect h2) ->
        IRHandle inner (fuseHandlers h1 h2)
  _ -> term

fuseHandlers :: IRHandler -> IRHandler -> IRHandler
fuseHandlers h1 h2 = IRHandler
  { irhEffect = irhEffect h1 `effectUnion` irhEffect h2
  , irhReturn = \x -> irhReturn h1 (irhReturn h2 x)
  , irhOps = irhOps h1 ++ irhOps h2
  }
\end{verbatim}
\end{definition}

\subsection{Handler Inlining}
\label{subsec:handler-inlining}

\begin{definition}[Static Handler Inlining]
\label{def:handler-inlining}
When a handler's effect operations are statically known, the handler can be inlined:
\begin{verbatim}
inlineHandler :: IRTerm -> IRTerm
inlineHandler term = case term of
  IRHandle body h | canInline h body ->
    inlineHandlerClauses h body
  _ -> term

canInline :: IRHandler -> IRTerm -> Bool
canInline h body =
  -- All perform calls in body are to known operations
  all (isKnownOp h) (collectPerforms body) &&
  -- No dynamic handler references
  not (hasDynamicHandlerRefs body)

inlineHandlerClauses :: IRHandler -> IRTerm -> IRTerm
inlineHandlerClauses h = transform go where
  go (IRPerform eff op v)
    | eff == irhEffect h =
        let (_, argName, kName, opBody) = findOp op (irhOps h)
        in opBody[v/argName, identity/kName]  -- simplified
  go t = t
\end{verbatim}
\end{definition}

\subsection{Ordinal Loop Optimization}
\label{subsec:ordinal-opt}

\begin{definition}[Finite Ordinal Detection]
\label{def:finite-ordinal-detect}
Ordinal loops with finite bounds can be converted to standard loops:
\begin{verbatim}
optimizeOrdinalLoop :: IRTerm -> IRTerm
optimizeOrdinalLoop term = case term of
  IROrdLimit i alpha body
    | isFiniteOrdinal alpha ->
        let n = toNat alpha
        in IRForLoop i 0 n body
  _ -> term

isFiniteOrdinal :: IRVal -> Bool
isFiniteOrdinal (IROrdinal (OrdNat n)) = True
isFiniteOrdinal _ = False
\end{verbatim}
\end{definition}

\begin{definition}[Limit Ordinal Approximation]
\label{def:limit-approx}
For transfinite loops at limit ordinals, the optimizer inserts convergence checks:
\begin{verbatim}
approximateLimitLoop :: IRTerm -> IRTerm
approximateLimitLoop term = case term of
  IROrdLimit i alpha body
    | isLimitOrdinal alpha ->
        IRWhileConverge i body (convergenceCheck alpha)
  _ -> term

convergenceCheck :: IRVal -> IRVal -> IRVal -> Bool
convergenceCheck alpha prev curr =
  -- Check if iteration has stabilized
  prev == curr || iterationCount > maxIterations alpha
\end{verbatim}
\end{definition}

%% ----------------------------------------------------------------------------
%% HANDOFF NOTE
%% ----------------------------------------------------------------------------
\begin{tcolorbox}[colback=green!5!white,colframe=green!50!black,title=\textbf{Handoff: COMPILER-TASK-04 Complete},breakable]
\textbf{Completed in this chapter:}
\begin{itemize}
  \item \textbf{Extended Lexer} (Section~\ref{sec:extended-lexer}):
    \begin{itemize}
      \item v7.3 token set extensions (transfinite, effect, quantum)
      \item Unicode support with ASCII fallbacks
      \item Lexer DFA state extensions
      \item Keyword recognition table
      \item Backwards compatibility proposition
    \end{itemize}
  \item \textbf{Extended Parser} (Section~\ref{sec:extended-parser}):
    \begin{itemize}
      \item Extended grammar productions for all v7.3 constructs
      \item Precedence and associativity table
      \item Recursive descent parser extensions
      \item Error recovery with synchronization points
    \end{itemize}
  \item \textbf{Extended AST} (Section~\ref{sec:extended-ast}):
    \begin{itemize}
      \item Complete AST data types with v7.3 nodes
      \item Effect annotations on expressions
      \item Source location tracking
      \item Well-formedness conditions
    \end{itemize}
  \item \textbf{Effect-Aware IR} (Section~\ref{sec:effect-aware-ir}):
    \begin{itemize}
      \item ANF-based IR with explicit effects
      \item ANF transformation rules
      \item CPS transformation for handlers
      \item \textbf{Theorem~\ref{thm:ir-typing-preservation}}: IR preserves typing judgments
      \item v7.2 backwards compatibility corollary
    \end{itemize}
  \item \textbf{Optimization Passes} (Section~\ref{sec:optimization-passes}):
    \begin{itemize}
      \item Effect purity analysis
      \item Dead effect elimination
      \item Handler fusion for commutative effects
      \item Handler inlining
      \item Ordinal loop optimization (finite detection, limit approximation)
    \end{itemize}
\end{itemize}

\textbf{Next tasks (COMPILER-TASK-05):}
\begin{itemize}
  \item Chapter 5: TKS Virtual Machine
    \begin{itemize}
      \item VM architecture (stack-based design, registers)
      \item Complete instruction set specification
      \item RPM-aware execution model
      \item Effect handling in VM (continuation capture, handler stack)
      \item Transfinite loop execution
      \item Garbage collection for TKS values
    \end{itemize}
\end{itemize}
\end{tcolorbox}

%% ============================================================================
%% CHAPTER 5: TKS VIRTUAL MACHINE
%% ============================================================================
\chapter{TKS Virtual Machine}
\label{ch:tks-vm}

\begin{tcolorbox}[colback=blue!5!white,colframe=blue!75!black,title=\vnewthree]
This chapter specifies the TKS Virtual Machine with support for RPM-aware execution, effect handling, and transfinite iteration.
\end{tcolorbox}

%% ----------------------------------------------------------------------------
\section{VM Architecture}
\label{sec:vm-architecture}

This section specifies the TKS Virtual Machine architecture, extending the v7.2 stack-based design with support for effect handlers, transfinite iteration, and quantum state management.

\subsection{Design Overview}
\label{subsec:vm-design-overview}

\begin{definition}[TKS VM Design Principles]
\label{def:vm-principles}
The TKS v7.3 VM is a \textbf{stack-based} virtual machine with the following design principles:
\begin{enumerate}
  \item \textbf{Operand stack}: All computations use an operand stack for intermediate values.
  \item \textbf{Handler stack}: A separate stack for installed effect handlers (v7.3).
  \item \textbf{Call stack}: Frames for function calls with local environments.
  \item \textbf{Heap}: Dynamically allocated objects (closures, Ideas, quantum states).
  \item \textbf{RPM context}: Persistent acquisition state across computations.
  \item \textbf{Ordinal register}: Special register for transfinite loop indices (v7.3).
\end{enumerate}
\end{definition}

\subsection{VM State}
\label{subsec:vm-state}

\begin{definition}[Extended VM State]
\label{def:vm-state-v73}
The v7.3 VM state extends v7.2:
\[
\VM_{7.3} = (\mathsf{code}, \mathsf{pc}, \mathsf{stack}, \mathsf{env}, \mathsf{frames}, \mathsf{handlers}, \mathsf{heap}, \alpha, \mathsf{ord}, \mathsf{qreg})
\]
where:
\begin{itemize}
  \item $\mathsf{code} : \BC^*$ --- bytecode program
  \item $\mathsf{pc} : \mathbb{N}$ --- program counter
  \item $\mathsf{stack} : \mathsf{Val}^*$ --- operand stack
  \item $\mathsf{env} : \mathsf{Ident} \rightharpoonup \mathsf{Val}$ --- local environment
  \item $\mathsf{frames} : \mathsf{Frame}^*$ --- call stack
  \item $\mathsf{handlers} : \mathsf{Handler}^*$ --- handler stack \textbf{(NEW v7.3)}
  \item $\mathsf{heap} : \mathsf{Addr} \rightharpoonup \mathsf{HeapObj}$ --- heap memory
  \item $\alpha : \mathcal{P}(\mathcal{E})$ --- RPM acquisition state
  \item $\mathsf{ord} : \mathsf{Ordinal}$ --- ordinal loop register \textbf{(NEW v7.3)}
  \item $\mathsf{qreg} : \mathsf{QState}$ --- quantum state register \textbf{(NEW v7.3)}
\end{itemize}
\end{definition}

\begin{definition}[Stack Frame Structure]
\label{def:stack-frame}
Each call stack frame contains:
\begin{verbatim}
data Frame = Frame
  { fReturnAddr  :: Int           -- Return program counter
  , fEnv         :: Env           -- Saved local environment
  , fStackBase   :: Int           -- Operand stack base pointer
  , fHandlerMark :: Int           -- Handler stack depth at call
  , fContinuation :: Maybe Cont   -- For effect resumption (NEW v7.3)
  }
\end{verbatim}
\end{definition}

\begin{definition}[Handler Stack Entry]
\label{def:handler-entry}
Each handler stack entry (v7.3):
\begin{verbatim}
data HandlerEntry = HandlerEntry
  { heEffect     :: EffectName        -- Effect being handled
  , heReturnAddr :: Int               -- Address of return clause code
  , heOpAddrs    :: Map OpName Int    -- Addresses of operation clauses
  , heContinuation :: Cont            -- Captured continuation
  , heEnv        :: Env               -- Environment at handler installation
  , heStackMark  :: Int               -- Operand stack depth at installation
  }
\end{verbatim}
\end{definition}

\subsection{Memory Model}
\label{subsec:memory-model}

\begin{definition}[VM Values]
\label{def:vm-values-v73}
Runtime values (extending v7.2):
\begin{align}
\mathsf{Val} ::=&\; \mathsf{VInt}(n) \mid \mathsf{VBool}(b) \mid \mathsf{VUnit} \\
    &\mid \mathsf{VElement}(W, n) \mid \mathsf{VFoundation}(m, w) \\
    &\mid \mathsf{VNoetic}(k) \mid \mathsf{VIdea}(\mathsf{addr}) \\
    &\mid \mathsf{VClosure}(\mathsf{addr}) \mid \mathsf{VFractal}(\mathsf{addr}) \\
    &\mid \mathsf{VRPMResult}(r) \quad \text{where } r \in \{\mathsf{Success}(v), \mathsf{Fail}\} \\
    %% v7.3 additions
    &\mid \mathsf{VOrdinal}(\alpha) \quad \text{(NEW v7.3)} \\
    &\mid \mathsf{VHandler}(\mathsf{addr}) \quad \text{(NEW v7.3)} \\
    &\mid \mathsf{VContinuation}(\mathsf{addr}) \quad \text{(NEW v7.3)} \\
    &\mid \mathsf{VQState}(\mathsf{addr}) \quad \text{(NEW v7.3)} \\
    &\mid \mathsf{VKet}(\mathsf{addr}) \quad \text{(NEW v7.3)}
\end{align}
\end{definition}

\begin{definition}[Heap Objects]
\label{def:heap-objects}
Heap-allocated objects:
\begin{verbatim}
data HeapObj
  -- v7.2 objects
  = HOClosure { hocCode :: Int, hocEnv :: Env }
  | HOIdea { hoiContent :: Val, hoiRefs :: [Addr] }
  | HOFractal { hofLevels :: [Int], hofContent :: Val }
  -- v7.3 objects
  | HOHandler                                   -- NEW
      { hohEffect :: EffectName
      , hohRetCode :: Int
      , hohOps :: Map OpName Int
      , hohEnv :: Env
      }
  | HOContinuation                              -- NEW
      { hocFrames :: [Frame]
      , hocStack :: [Val]
      , hocHandlers :: [HandlerEntry]
      , hocPC :: Int
      }
  | HOQState                                    -- NEW
      { hoqAmplitudes :: Map Val Complex
      , hoqEntangled :: [Addr]
      }
\end{verbatim}
\end{definition}

\begin{definition}[Value Size and Alignment]
\label{def:value-size}
Value representation in memory:
\begin{center}
\begin{tabular}{lcc}
\toprule
\textbf{Value Type} & \textbf{Size (bytes)} & \textbf{Alignment} \\
\midrule
\textsf{VInt} & 8 & 8 \\
\textsf{VBool} & 1 & 1 \\
\textsf{VUnit} & 0 & 1 \\
\textsf{VElement} & 2 & 2 \\
\textsf{VNoetic} & 1 & 1 \\
\textsf{VOrdinal} & 16 & 8 \\
Heap pointer & 8 & 8 \\
\bottomrule
\end{tabular}
\end{center}
\end{definition}

%% ----------------------------------------------------------------------------
\section{Instruction Set}
\label{sec:instruction-set}

This section specifies the complete TKS v7.3 bytecode instruction set.

\subsection{Instruction Encoding}
\label{subsec:instruction-encoding}

\begin{definition}[Instruction Format]
\label{def:instruction-format}
Each instruction is encoded as:
\begin{center}
\begin{tabular}{|c|c|c|c|}
\hline
\textbf{Opcode} & \textbf{Flags} & \textbf{Operand 1} & \textbf{Operand 2} \\
\hline
1 byte & 1 byte & 0--8 bytes & 0--8 bytes \\
\hline
\end{tabular}
\end{center}
The flags byte encodes:
\begin{itemize}
  \item Bits 0--1: Operand 1 size (0=none, 1=1 byte, 2=4 bytes, 3=8 bytes)
  \item Bits 2--3: Operand 2 size
  \item Bits 4--7: Instruction-specific flags
\end{itemize}
\end{definition}

\subsection{Core Instructions}
\label{subsec:core-instructions}

\begin{definition}[Core Instruction Set]
\label{def:core-instructions}
Basic stack and control flow instructions:
\begin{center}
\begin{tabular}{llp{7cm}}
\toprule
\textbf{Opcode} & \textbf{Mnemonic} & \textbf{Operation} \\
\midrule
\texttt{0x00} & \texttt{NOP} & No operation \\
\texttt{0x01} & \texttt{PUSH\_INT n} & Push integer $n$ onto stack \\
\texttt{0x02} & \texttt{PUSH\_BOOL b} & Push boolean $b$ onto stack \\
\texttt{0x03} & \texttt{PUSH\_UNIT} & Push unit value \\
\texttt{0x04} & \texttt{POP} & Pop and discard top of stack \\
\texttt{0x05} & \texttt{DUP} & Duplicate top of stack \\
\texttt{0x06} & \texttt{SWAP} & Swap top two stack elements \\
\texttt{0x07} & \texttt{ROT} & Rotate top three elements \\
\midrule
\texttt{0x10} & \texttt{LOAD x} & Push value of variable $x$ \\
\texttt{0x11} & \texttt{STORE x} & Pop and store in variable $x$ \\
\texttt{0x12} & \texttt{LOAD\_GLOBAL x} & Load from global environment \\
\texttt{0x13} & \texttt{STORE\_GLOBAL x} & Store to global environment \\
\midrule
\texttt{0x20} & \texttt{JMP addr} & Unconditional jump \\
\texttt{0x21} & \texttt{JMP\_IF addr} & Jump if top of stack is true \\
\texttt{0x22} & \texttt{JMP\_UNLESS addr} & Jump if top of stack is false \\
\texttt{0x23} & \texttt{CALL n} & Call function with $n$ arguments \\
\texttt{0x24} & \texttt{RET} & Return from function \\
\texttt{0x25} & \texttt{TAIL\_CALL n} & Tail call optimization \\
\midrule
\texttt{0x30} & \texttt{ADD} & Integer addition \\
\texttt{0x31} & \texttt{SUB} & Integer subtraction \\
\texttt{0x32} & \texttt{MUL} & Integer multiplication \\
\texttt{0x33} & \texttt{DIV} & Integer division \\
\texttt{0x34} & \texttt{EQ} & Equality comparison \\
\texttt{0x35} & \texttt{LT} & Less than comparison \\
\texttt{0x36} & \texttt{AND} & Boolean and \\
\texttt{0x37} & \texttt{OR} & Boolean or \\
\texttt{0x38} & \texttt{NOT} & Boolean negation \\
\bottomrule
\end{tabular}
\end{center}
\end{definition}

\subsection{TKS Domain Instructions}
\label{subsec:tks-instructions}

\begin{definition}[Element and Foundation Instructions]
\label{def:element-instructions}
\begin{center}
\begin{tabular}{llp{7cm}}
\toprule
\textbf{Opcode} & \textbf{Mnemonic} & \textbf{Operation} \\
\midrule
\texttt{0x40} & \texttt{PUSH\_ELEMENT W n} & Push Element $W_n$ \\
\texttt{0x41} & \texttt{PUSH\_FOUNDATION m w} & Push Foundation $F_{m,w}$ \\
\texttt{0x42} & \texttt{ELEMENT\_WING} & Extract wing from element \\
\texttt{0x43} & \texttt{ELEMENT\_INDEX} & Extract index from element \\
\texttt{0x44} & \texttt{FOUNDATION\_LEVEL} & Extract level from foundation \\
\texttt{0x45} & \texttt{FOUNDATION\_ASPECT} & Extract aspect from foundation \\
\bottomrule
\end{tabular}
\end{center}
\end{definition}

\begin{definition}[Noetic Instructions]
\label{def:noetic-instructions}
\begin{center}
\begin{tabular}{llp{7cm}}
\toprule
\textbf{Opcode} & \textbf{Mnemonic} & \textbf{Operation} \\
\midrule
\texttt{0x50} & \texttt{PUSH\_NOETIC k} & Push noetic operator $\nu_k$ \\
\texttt{0x51} & \texttt{APPLY\_NOETIC} & Apply noetic: pop $\nu_k$, pop $v$, push $\nu_k(v)$ \\
\texttt{0x52} & \texttt{NOETIC\_COMPOSE} & Compose two noetics: $\nu_j \circ \nu_k$ \\
\texttt{0x53} & \texttt{NOETIC\_LEVEL} & Get noetic level $k$ from $\nu_k$ \\
\bottomrule
\end{tabular}
\end{center}
\end{definition}

\begin{definition}[Fractal and ACBE Instructions]
\label{def:fractal-instructions}
\begin{center}
\begin{tabular}{llp{7cm}}
\toprule
\textbf{Opcode} & \textbf{Mnemonic} & \textbf{Operation} \\
\midrule
\texttt{0x58} & \texttt{MAKE\_FRACTAL n} & Create fractal from top $n$ levels \\
\texttt{0x59} & \texttt{FRACTAL\_LEVEL k} & Extract level $k$ from fractal \\
\texttt{0x5A} & \texttt{ACBE\_INIT} & Initialize ACBE with goal Idea \\
\texttt{0x5B} & \texttt{ACBE\_DESCEND} & One ACBE descent step \\
\texttt{0x5C} & \texttt{ACBE\_COMPLETE} & Check if ACBE is complete \\
\texttt{0x5D} & \texttt{ACBE\_RESULT} & Extract ACBE result \\
\bottomrule
\end{tabular}
\end{center}
\end{definition}

\subsection{RPM Instructions}
\label{subsec:rpm-instructions}

\begin{definition}[RPM Monad Instructions]
\label{def:rpm-instructions}
\begin{center}
\begin{tabular}{llp{7cm}}
\toprule
\textbf{Opcode} & \textbf{Mnemonic} & \textbf{Operation} \\
\midrule
\texttt{0x60} & \texttt{RPM\_RETURN} & Wrap value in RPM success \\
\texttt{0x61} & \texttt{RPM\_BIND} & Monadic bind for RPM \\
\texttt{0x62} & \texttt{RPM\_CHECK} & Check element prerequisite \\
\texttt{0x63} & \texttt{RPM\_ACQUIRE} & Acquire element \\
\texttt{0x64} & \texttt{RPM\_FAIL} & RPM computation failure \\
\texttt{0x65} & \texttt{RPM\_IS\_SUCCESS} & Test if RPM result is success \\
\texttt{0x66} & \texttt{RPM\_UNWRAP} & Extract value from success \\
\texttt{0x67} & \texttt{RPM\_STATE} & Push current acquisition state \\
\bottomrule
\end{tabular}
\end{center}
\end{definition}

\subsection{Effect Instructions (v7.3)}
\label{subsec:effect-instructions}

\begin{definition}[Effect Handler Instructions]
\label{def:effect-instructions}
New instructions for algebraic effects:
\begin{center}
\begin{tabular}{llp{6.5cm}}
\toprule
\textbf{Opcode} & \textbf{Mnemonic} & \textbf{Operation} \\
\midrule
\texttt{0x70} & \texttt{INSTALL\_HANDLER addr} & Install handler at \texttt{addr} onto handler stack \\
\texttt{0x71} & \texttt{REMOVE\_HANDLER} & Remove top handler from stack \\
\texttt{0x72} & \texttt{PERFORM\_EFFECT E op} & Perform effect operation $E.\mathit{op}$ \\
\texttt{0x73} & \texttt{RESUME} & Resume captured continuation \\
\texttt{0x74} & \texttt{RESUME\_WITH v} & Resume with value $v$ \\
\texttt{0x75} & \texttt{CAPTURE\_CONT} & Capture current continuation \\
\texttt{0x76} & \texttt{ABORT\_CONT} & Abort to nearest handler \\
\texttt{0x77} & \texttt{HANDLER\_RETURN} & Execute handler's return clause \\
\bottomrule
\end{tabular}
\end{center}
\end{definition}

\subsection{Transfinite Instructions (v7.3)}
\label{subsec:transfinite-instructions}

\begin{definition}[Ordinal and Transfinite Instructions]
\label{def:transfinite-instructions}
New instructions for transfinite computation:
\begin{center}
\begin{tabular}{llp{6.5cm}}
\toprule
\textbf{Opcode} & \textbf{Mnemonic} & \textbf{Operation} \\
\midrule
\texttt{0x80} & \texttt{PUSH\_ORD $\alpha$} & Push ordinal literal \\
\texttt{0x81} & \texttt{PUSH\_OMEGA} & Push $\omega$ \\
\texttt{0x82} & \texttt{PUSH\_EPSILON n} & Push $\varepsilon_n$ \\
\texttt{0x83} & \texttt{ORD\_SUCC} & Ordinal successor \\
\texttt{0x84} & \texttt{ORD\_ADD} & Ordinal addition \\
\texttt{0x85} & \texttt{ORD\_MUL} & Ordinal multiplication \\
\texttt{0x86} & \texttt{ORD\_EXP} & Ordinal exponentiation \\
\texttt{0x87} & \texttt{ORD\_LT} & Ordinal comparison $<$ \\
\texttt{0x88} & \texttt{ORD\_IS\_LIMIT} & Test if ordinal is limit \\
\texttt{0x89} & \texttt{ORD\_IS\_FINITE} & Test if ordinal is finite \\
\midrule
\texttt{0x90} & \texttt{TRANS\_LOOP\_INIT} & Initialize transfinite loop \\
\texttt{0x91} & \texttt{TRANS\_LOOP\_STEP} & One loop iteration step \\
\texttt{0x92} & \texttt{TRANS\_LOOP\_CHECK} & Check loop termination/convergence \\
\texttt{0x93} & \texttt{TRANS\_LOOP\_END} & Finalize transfinite loop \\
\texttt{0x94} & \texttt{ORD\_LIMIT} & Compute limit ordinal \\
\bottomrule
\end{tabular}
\end{center}
\end{definition}

\subsection{Quantum Instructions (v7.3)}
\label{subsec:quantum-instructions}

\begin{definition}[Quantum State Instructions]
\label{def:quantum-instructions}
Instructions for quantum-inspired computation:
\begin{center}
\begin{tabular}{llp{6.5cm}}
\toprule
\textbf{Opcode} & \textbf{Mnemonic} & \textbf{Operation} \\
\midrule
\texttt{0xA0} & \texttt{MAKE\_KET} & Create ket state $|v\rangle$ \\
\texttt{0xA1} & \texttt{MAKE\_BRA} & Create bra state $\langle v|$ \\
\texttt{0xA2} & \texttt{BRAKET} & Inner product $\langle u|v\rangle$ \\
\texttt{0xA3} & \texttt{SUPERPOSE n} & Create superposition of $n$ states \\
\texttt{0xA4} & \texttt{MEASURE} & Measure quantum state \\
\texttt{0xA5} & \texttt{ENTANGLE} & Entangle two states \\
\texttt{0xA6} & \texttt{TENSOR} & Tensor product $\otimes$ \\
\texttt{0xA7} & \texttt{Q\_APPLY} & Apply quantum operation \\
\bottomrule
\end{tabular}
\end{center}
\end{definition}

\subsection{Closure and Allocation Instructions}
\label{subsec:closure-instructions}

\begin{definition}[Closure and Heap Instructions]
\label{def:closure-heap-instructions}
\begin{center}
\begin{tabular}{llp{6.5cm}}
\toprule
\textbf{Opcode} & \textbf{Mnemonic} & \textbf{Operation} \\
\midrule
\texttt{0xB0} & \texttt{MAKE\_CLOSURE addr n} & Create closure capturing $n$ variables \\
\texttt{0xB1} & \texttt{CLOSURE\_APPLY} & Apply closure to argument \\
\texttt{0xB2} & \texttt{ALLOC n} & Allocate $n$ bytes on heap \\
\texttt{0xB3} & \texttt{LOAD\_HEAP} & Load from heap address \\
\texttt{0xB4} & \texttt{STORE\_HEAP} & Store to heap address \\
\texttt{0xB5} & \texttt{MAKE\_IDEA} & Create Idea object \\
\texttt{0xB6} & \texttt{IDEA\_CONTENT} & Extract Idea content \\
\bottomrule
\end{tabular}
\end{center}
\end{definition}

\subsection{Instruction Examples}
\label{subsec:instruction-examples}

\begin{example}[Bytecode for Noetic Application]
\label{ex:bytecode-noetic}
The expression $\nu_2(\mathsf{A3})$ compiles to:
\begin{verbatim}
  PUSH_ELEMENT A 3     ; stack: [A3]
  PUSH_NOETIC 2        ; stack: [A3, nu_2]
  APPLY_NOETIC         ; stack: [nu_2(A3)]
\end{verbatim}
\end{example}

\begin{example}[Bytecode for Effect Handler]
\label{ex:bytecode-handler}
The expression \texttt{handle e with h} compiles to:
\begin{verbatim}
  INSTALL_HANDLER @h   ; Install handler h
  ; ... code for e ...
  HANDLER_RETURN       ; Execute return clause
  REMOVE_HANDLER       ; Clean up handler
\end{verbatim}
\end{example}

\begin{example}[Bytecode for Transfinite Loop]
\label{ex:bytecode-transfinite}
The expression \texttt{limit i < omega. body} compiles to:
\begin{verbatim}
  PUSH_OMEGA           ; Push limit ordinal
  TRANS_LOOP_INIT      ; Initialize loop
loop_start:
  TRANS_LOOP_CHECK     ; Check termination
  JMP_IF loop_end      ; Exit if converged
  ; ... code for body ...
  TRANS_LOOP_STEP      ; Increment ordinal index
  JMP loop_start
loop_end:
  TRANS_LOOP_END       ; Finalize and get result
\end{verbatim}
\end{example}

%% ----------------------------------------------------------------------------
\section{RPM-Aware Execution}
\label{sec:rpm-aware-execution}

This section specifies how the TKS VM maintains and operates on RPM (Requisite Progression Model) state.

\subsection{RPM State in VM Context}
\label{subsec:rpm-state}

\begin{definition}[RPM State Structure]
\label{def:rpm-state-structure}
The RPM state $\alpha$ in the VM maintains:
\begin{verbatim}
data RPMState = RPMState
  { rpmAcquired :: Set Element       -- Acquired elements
  , rpmGoals    :: [Goal]            -- Active goal stack
  , rpmPrereqs  :: Map Element [Element]  -- Prerequisite graph
  , rpmHistory  :: [RPMEvent]        -- Acquisition history
  }

data Goal = Goal
  { goalTarget  :: Element           -- Target element
  , goalPath    :: [Element]         -- Planned acquisition path
  , goalStatus  :: GoalStatus        -- Active, Blocked, Complete
  }

data RPMEvent
  = AcquiredEvent Element Timestamp
  | CheckedEvent Element Bool Timestamp
  | FailedEvent Element String Timestamp
\end{verbatim}
\end{definition}

\begin{definition}[Prerequisite Graph]
\label{def:prereq-graph}
The prerequisite graph encodes TKS progression rules:
\begin{itemize}
  \item \textbf{Wing progression}: $A_n \to A_{n+1}$ requires $A_n$ acquired.
  \item \textbf{Foundation dependencies}: $F_{m,w}$ requires elements in wing $W$ at appropriate levels.
  \item \textbf{Noetic prerequisites}: $\nu_k$ application requires specific acquisition patterns.
\end{itemize}
The graph is initialized from the TKS element catalog and can be extended by user declarations.
\end{definition}

\subsection{RPM Instruction Semantics}
\label{subsec:rpm-semantics}

\begin{definition}[RPM\_CHECK Semantics]
\label{def:rpm-check-semantics}
The \texttt{RPM\_CHECK} instruction:
\begin{align}
&\mathsf{RPM\_CHECK} : \\
&\quad \text{Pre: } \mathsf{stack} = e :: \mathsf{rest}, \; e : \mathsf{Element} \\
&\quad \text{Post: } \mathsf{stack} = \mathsf{VBool}(b) :: \mathsf{rest} \\
&\quad \text{where } b = (\mathsf{prereqs}(e) \subseteq \alpha.\mathsf{acquired})
\end{align}
The instruction checks if all prerequisites of element $e$ are satisfied by the current acquisition state.
\end{definition}

\begin{definition}[RPM\_ACQUIRE Semantics]
\label{def:rpm-acquire-semantics}
The \texttt{RPM\_ACQUIRE} instruction:
\begin{verbatim}
RPM_ACQUIRE:
  e <- pop stack
  if prereqs(e) not subset of alpha.acquired:
    push VRPMResult(Fail)
    record FailedEvent(e, "Prerequisites not met")
  else if e in alpha.acquired:
    push VRPMResult(Success(VUnit))  -- Already acquired
  else:
    alpha.acquired <- alpha.acquired + {e}
    record AcquiredEvent(e, now())
    push VRPMResult(Success(VUnit))
\end{verbatim}
\end{definition}

\begin{definition}[RPM\_BIND Semantics]
\label{def:rpm-bind-semantics}
The \texttt{RPM\_BIND} instruction implements monadic sequencing:
\begin{verbatim}
RPM_BIND:
  f <- pop stack        -- continuation function
  m <- pop stack        -- RPM computation result
  case m of
    VRPMResult(Fail) ->
      push VRPMResult(Fail)  -- Propagate failure
    VRPMResult(Success(v)) ->
      push v
      push f
      CALL 1            -- Apply f to v
\end{verbatim}
\end{definition}

\subsection{Goal and Prerequisite Tracking}
\label{subsec:goal-tracking}

\begin{definition}[Goal Management Instructions]
\label{def:goal-instructions}
Additional instructions for goal tracking:
\begin{center}
\begin{tabular}{llp{6.5cm}}
\toprule
\textbf{Opcode} & \textbf{Mnemonic} & \textbf{Operation} \\
\midrule
\texttt{0x68} & \texttt{RPM\_SET\_GOAL} & Set acquisition goal \\
\texttt{0x69} & \texttt{RPM\_GOAL\_STATUS} & Query goal status \\
\texttt{0x6A} & \texttt{RPM\_PLAN\_PATH} & Compute acquisition path \\
\texttt{0x6B} & \texttt{RPM\_NEXT\_STEP} & Get next element to acquire \\
\bottomrule
\end{tabular}
\end{center}
\end{definition}

\begin{example}[RPM Goal Planning]
\label{ex:rpm-goal}
Setting a goal to acquire D5 (assuming prerequisites):
\begin{verbatim}
  PUSH_ELEMENT D 5     ; Target element
  RPM_SET_GOAL         ; Set as goal
  RPM_PLAN_PATH        ; Compute: [A1,A2,B1,B2,C1,...,D5]
loop:
  RPM_NEXT_STEP        ; Get next element
  DUP
  PUSH_UNIT
  EQ
  JMP_IF done          ; Exit if no more steps
  RPM_ACQUIRE          ; Acquire element
  RPM_IS_SUCCESS
  JMP_UNLESS fail      ; Handle failure
  JMP loop
done:
  PUSH_UNIT
  RPM_RETURN
fail:
  RPM_FAIL
\end{verbatim}
\end{example}

%% ----------------------------------------------------------------------------
\section{Effect Handling in VM}
\label{sec:effect-handling-vm}

This section specifies the runtime implementation of algebraic effect handlers in the TKS VM.

\subsection{Handler Stack Management}
\label{subsec:handler-stack}

\begin{definition}[Handler Stack Operations]
\label{def:handler-stack-ops}
The handler stack supports the following operations:
\begin{verbatim}
pushHandler :: HandlerEntry -> VM -> VM
pushHandler h vm = vm { handlers = h : handlers vm }

popHandler :: VM -> (HandlerEntry, VM)
popHandler vm = case handlers vm of
  []     -> error "Handler stack underflow"
  (h:hs) -> (h, vm { handlers = hs })

findHandler :: EffectName -> VM -> Maybe (HandlerEntry, Int)
findHandler eff vm = go (handlers vm) 0 where
  go [] _ = Nothing
  go (h:hs) depth
    | heEffect h == eff = Just (h, depth)
    | otherwise         = go hs (depth + 1)
\end{verbatim}
\end{definition}

\begin{definition}[INSTALL\_HANDLER Semantics]
\label{def:install-handler-semantics}
The \texttt{INSTALL\_HANDLER addr} instruction:
\begin{verbatim}
INSTALL_HANDLER addr:
  -- Read handler object from heap
  hobj <- readHeap(addr)
  -- Create handler entry
  entry <- HandlerEntry
    { heEffect = hohEffect hobj
    , heReturnAddr = hohRetCode hobj
    , heOpAddrs = hohOps hobj
    , heContinuation = emptyCont
    , heEnv = env
    , heStackMark = length stack
    }
  -- Push onto handler stack
  handlers <- entry : handlers
  pc <- pc + 1
\end{verbatim}
\end{definition}

\subsection{Continuation Capture and Restoration}
\label{subsec:continuation-capture}

\begin{definition}[Continuation Structure]
\label{def:continuation-structure}
A captured continuation represents suspended computation state:
\begin{verbatim}
data Continuation = Cont
  { contFrames   :: [Frame]        -- Saved call stack
  , contStack    :: [Val]          -- Saved operand stack
  , contHandlers :: [HandlerEntry] -- Saved handler stack
  , contPC       :: Int            -- Resume address
  , contEnv      :: Env            -- Saved environment
  }
\end{verbatim}
\end{definition}

\begin{definition}[PERFORM\_EFFECT Semantics]
\label{def:perform-effect-semantics}
The \texttt{PERFORM\_EFFECT E op} instruction:
\begin{verbatim}
PERFORM_EFFECT E op:
  arg <- pop stack
  case findHandler E of
    Nothing ->
      -- Unhandled effect: runtime error
      error ("Unhandled effect: " ++ E ++ "." ++ op)
    Just (handler, depth) ->
      -- Capture continuation up to handler
      k <- captureContinuation depth
      -- Restore to handler's installation point
      stack <- take (heStackMark handler) stack
      frames <- take (length frames - depth) frames
      handlers <- drop (depth + 1) handlers
      env <- heEnv handler
      -- Push argument and continuation for op clause
      push arg
      push (VContinuation k)
      -- Jump to operation clause
      pc <- heOpAddrs handler ! op

captureContinuation :: Int -> VM -> Continuation
captureContinuation depth vm = Cont
  { contFrames   = take depth (frames vm)
  , contStack    = drop (heStackMark h) (stack vm)
  , contHandlers = take depth (handlers vm)
  , contPC       = pc vm + 1
  , contEnv      = env vm
  }
  where h = handlers vm !! depth
\end{verbatim}
\end{definition}

\begin{definition}[RESUME Semantics]
\label{def:resume-semantics}
The \texttt{RESUME} and \texttt{RESUME\_WITH} instructions:
\begin{verbatim}
RESUME:
  -- Resume with unit value
  k <- pop stack
  resumeContinuation k VUnit

RESUME_WITH:
  v <- pop stack
  k <- pop stack
  resumeContinuation k v

resumeContinuation :: Continuation -> Val -> VM -> VM
resumeContinuation k v vm = vm
  { frames   = contFrames k ++ frames vm
  , stack    = v : contStack k
  , handlers = contHandlers k ++ handlers vm
  , pc       = contPC k
  , env      = contEnv k
  }
\end{verbatim}
\end{definition}

\subsection{Handler Return Clause}
\label{subsec:handler-return}

\begin{definition}[HANDLER\_RETURN Semantics]
\label{def:handler-return-semantics}
When a handled computation completes normally:
\begin{verbatim}
HANDLER_RETURN:
  v <- pop stack
  (handler, _) <- popHandler
  -- Jump to return clause with value bound
  env <- heEnv handler
  push v
  pc <- heReturnAddr handler
\end{verbatim}
\end{definition}

\subsection{Correctness Theorem}
\label{subsec:vm-correctness}

\begin{theorem}[VM Execution Preserves Operational Semantics]
\label{thm:vm-correctness-v73}
For any well-typed TKS expression $e$ with $\Gamma \vdash e : \tau \mathbin{!} \varepsilon$, let $\mathsf{compile}(e) = \mathsf{code}$. Then:

\begin{enumerate}
  \item \textbf{Type Preservation}: If the initial VM state has stack type $\sigma$ and the VM terminates normally, the final stack has type $\tau :: \sigma$.

  \item \textbf{Semantic Equivalence}: If $e \Downarrow v$ in the surface operational semantics, then running $\mathsf{code}$ from initial state $\VM_0$ terminates with $v$ on top of the stack:
  \[
  \VM_0 \Rightarrow^* \VM_f \quad \text{and} \quad \mathsf{stack}(\VM_f) = v :: \mathsf{stack}(\VM_0)
  \]

  \item \textbf{Effect Correspondence}: If $e$ performs effect operation $E.\mathit{op}(a)$ and is handled by $h$, then the VM execution:
  \begin{enumerate}
    \item Executes \texttt{PERFORM\_EFFECT E op} with $a$ on stack
    \item Captures continuation $k$
    \item Jumps to handler clause
    \item On \texttt{RESUME}, restores continuation with correct state
  \end{enumerate}
  matching the E-HandlePerform rule from the surface semantics.

  \item \textbf{RPM State Consistency}: The VM's $\alpha$ component evolves consistently with the RPM monad semantics---successful \texttt{RPM\_ACQUIRE} adds elements to $\alpha$, and \texttt{RPM\_CHECK} reflects the current acquisition state.
\end{enumerate}
\end{theorem}

\begin{proof}
We prove each part:

\textbf{Part 1 (Type Preservation):}
By induction on the compilation rules. Each instruction preserves stack typing:
\begin{itemize}
  \item \texttt{PUSH\_INT n}: $\sigma \to \mathsf{Int} :: \sigma$ $\checkmark$
  \item \texttt{CALL n}: $(\tau_1 \to \tau_2) :: \tau_1 :: \sigma \to \tau_2 :: \sigma$ $\checkmark$
  \item \texttt{PERFORM\_EFFECT E op}: $A :: \sigma \to B :: \sigma$ where $\mathit{op} : A \to B$ $\checkmark$
  \item \texttt{RESUME\_WITH}: $\mathsf{Cont}[\tau] :: \tau :: \sigma \to \sigma'$ where $\sigma'$ is continuation's expected stack $\checkmark$
\end{itemize}
The handler stack maintains typing invariants by construction.

\textbf{Part 2 (Semantic Equivalence):}
By structural induction on the derivation $e \Downarrow v$.

\emph{Base case}: Literals $n \Downarrow n$ compile to \texttt{PUSH\_INT n}, which produces $n$ on stack.

\emph{Inductive cases}:
\begin{itemize}
  \item \textbf{Application}: Given $e_1 \Downarrow \lambda x. e'$ and $e_2 \Downarrow v_2$ and $e'[v_2/x] \Downarrow v$.
  By IH, compiled $e_1$ produces closure, compiled $e_2$ produces $v_2$.
  \texttt{CALL} creates frame and executes body, which by IH produces $v$.

  \item \textbf{Handle}: Given the rule
  \[
  \frac{e \Downarrow v \quad h.\mathsf{return}(v) \Downarrow v'}{\mathsf{handle}\; e\; \mathsf{with}\; h \Downarrow v'}
  \]
  The compiled code:
  (1) \texttt{INSTALL\_HANDLER} pushes handler,
  (2) executes $e$'s code producing $v$ by IH,
  (3) \texttt{HANDLER\_RETURN} executes return clause,
  (4) by IH, produces $v'$.
\end{itemize}

\textbf{Part 3 (Effect Correspondence):}
The E-HandlePerform rule states:
\[
\frac{E[\mathsf{perform}\; \mathit{op}(v)] \quad h = \{\ldots, \mathit{op}(x)\; k \to e_{op}, \ldots\}}
     {\mathsf{handle}\; E[\mathsf{perform}\; \mathit{op}(v)]\; \mathsf{with}\; h \Downarrow e_{op}[v/x, (\lambda y.\, \mathsf{handle}\; E[y]\; \mathsf{with}\; h)/k]}
\]
The VM implementation:
\begin{enumerate}
  \item \texttt{PERFORM\_EFFECT} finds handler and captures continuation representing $E[\bullet]$.
  \item Arguments $v$ and continuation $k$ are pushed.
  \item Jump to $e_{op}$'s code with proper bindings.
  \item \texttt{RESUME} restores $E[\bullet]$ context with handler re-installed, matching the meta-substitution $\lambda y.\, \mathsf{handle}\; E[y]\; \mathsf{with}\; h$.
\end{enumerate}

\textbf{Part 4 (RPM Consistency):}
The RPM monad laws are preserved:
\begin{itemize}
  \item \texttt{RPM\_RETURN} creates $\mathsf{Success}(v)$, matching $\mathsf{return}$.
  \item \texttt{RPM\_BIND} propagates $\mathsf{Fail}$ and sequences $\mathsf{Success}$, matching $\gg\!=$.
  \item \texttt{RPM\_ACQUIRE} updates $\alpha$ only on success, maintaining monotonicity.
  \item \texttt{RPM\_CHECK} reflects current $\alpha$ without modification.
\end{itemize}
\end{proof}

%% ----------------------------------------------------------------------------
\section{Transfinite Loop Execution}
\label{sec:transfinite-loop-vm}

This section specifies how the VM executes loops with ordinal bounds.

\subsection{Ordinal Representation}
\label{subsec:ordinal-representation}

\begin{definition}[Runtime Ordinal Representation]
\label{def:ordinal-repr}
Ordinals are represented at runtime using Cantor Normal Form:
\begin{verbatim}
data Ordinal
  = OrdZero                      -- 0
  | OrdFinite Int                -- n (finite)
  | OrdOmega                     -- omega
  | OrdOmegaN Int                -- omega^n
  | OrdCNF [(Ordinal, Int)]      -- sum of omega^alpha * n
  | OrdEpsilon Int               -- epsilon_n
  | OrdLimit (Int -> Ordinal)    -- Limit representation
\end{verbatim}
\end{definition}

\begin{definition}[Ordinal Arithmetic in VM]
\label{def:ordinal-arithmetic-vm}
VM ordinal operations:
\begin{verbatim}
ordSucc :: Ordinal -> Ordinal
ordSucc OrdZero       = OrdFinite 1
ordSucc (OrdFinite n) = OrdFinite (n + 1)
ordSucc alpha         = OrdCNF [(alpha, 1), (OrdZero, 1)]

ordAdd :: Ordinal -> Ordinal -> Ordinal
ordAdd alpha OrdZero = alpha
ordAdd (OrdFinite m) (OrdFinite n) = OrdFinite (m + n)
ordAdd alpha beta = -- CNF addition (absorbs lower terms)
  normalizeCNF (toCNF alpha ++ toCNF beta)

ordMul :: Ordinal -> Ordinal -> Ordinal
-- Implements ordinal multiplication with right-absorption

ordLt :: Ordinal -> Ordinal -> Bool
-- Lexicographic comparison on CNF
\end{verbatim}
\end{definition}

\subsection{Transfinite Loop Execution Model}
\label{subsec:transfinite-loop-model}

\begin{definition}[Transfinite Loop State]
\label{def:transfinite-loop-state}
State maintained during transfinite loop execution:
\begin{verbatim}
data TransLoopState = TransLoopState
  { tlsBound     :: Ordinal       -- Loop bound
  , tlsCurrent   :: Ordinal       -- Current index
  , tlsIteration :: Int           -- Finite iteration count
  , tlsPrevValue :: Maybe Val     -- Previous result (for convergence)
  , tlsAccum     :: Val           -- Accumulated result
  , tlsConverged :: Bool          -- Has loop converged?
  }
\end{verbatim}
\end{definition}

\begin{definition}[TRANS\_LOOP\_INIT Semantics]
\label{def:trans-loop-init}
\begin{verbatim}
TRANS_LOOP_INIT:
  bound <- pop stack
  ord <- OrdZero                 -- Start at 0
  tlsState <- TransLoopState
    { tlsBound = bound
    , tlsCurrent = OrdZero
    , tlsIteration = 0
    , tlsPrevValue = Nothing
    , tlsAccum = VUnit
    , tlsConverged = False
    }
  -- Store in dedicated ordinal register
  ordReg <- tlsState
\end{verbatim}
\end{definition}

\begin{definition}[TRANS\_LOOP\_STEP Semantics]
\label{def:trans-loop-step}
\begin{verbatim}
TRANS_LOOP_STEP:
  result <- pop stack            -- Body result
  tls <- ordReg
  -- Update state
  tls' <- tls
    { tlsCurrent = ordSucc (tlsCurrent tls)
    , tlsIteration = tlsIteration tls + 1
    , tlsPrevValue = Just (tlsAccum tls)
    , tlsAccum = result
    }
  -- Check convergence for limit ordinals
  if isLimitOrdinal (tlsBound tls) then
    tls' <- tls' { tlsConverged = checkConvergence tls' }
  ordReg <- tls'
  push (tlsCurrent tls')         -- Push current index for body
\end{verbatim}
\end{definition}

\begin{definition}[TRANS\_LOOP\_CHECK Semantics]
\label{def:trans-loop-check}
\begin{verbatim}
TRANS_LOOP_CHECK:
  tls <- ordReg
  shouldTerminate <-
    tlsConverged tls ||
    ordLt (tlsBound tls) (tlsCurrent tls) ||
    tlsIteration tls > maxIterations (tlsBound tls)
  push (VBool shouldTerminate)

-- Maximum iterations based on bound
maxIterations :: Ordinal -> Int
maxIterations (OrdFinite n) = n
maxIterations OrdOmega      = 10000    -- Configurable limit
maxIterations _             = 100000   -- Larger transfinite bounds
\end{verbatim}
\end{definition}

\subsection{Convergence Detection}
\label{subsec:convergence-detection}

\begin{definition}[Convergence Criteria]
\label{def:convergence-criteria}
For limit ordinals, convergence is detected when:
\begin{verbatim}
checkConvergence :: TransLoopState -> Bool
checkConvergence tls = case tlsPrevValue tls of
  Nothing -> False
  Just prev ->
    -- Value equality (for Ideas)
    valueEqual prev (tlsAccum tls) ||
    -- Structural fixpoint (for fractals)
    structuralFixpoint prev (tlsAccum tls) ||
    -- Metric convergence (for numeric types)
    metricConverged prev (tlsAccum tls)

valueEqual :: Val -> Val -> Bool
valueEqual (VIdea a1) (VIdea a2) = heapEqual a1 a2
valueEqual v1 v2 = v1 == v2

structuralFixpoint :: Val -> Val -> Bool
structuralFixpoint (VFractal a1) (VFractal a2) =
  fractalLevelsEqual a1 a2
structuralFixpoint _ _ = False
\end{verbatim}
\end{definition}

%% ----------------------------------------------------------------------------
\section{Garbage Collection}
\label{sec:garbage-collection}

This section specifies memory management for TKS values.

\subsection{GC Design}
\label{subsec:gc-design}

\begin{definition}[Garbage Collection Strategy]
\label{def:gc-strategy}
The TKS VM uses a hybrid garbage collector:
\begin{enumerate}
  \item \textbf{Reference counting} for short-lived allocations and cycles detection.
  \item \textbf{Mark-and-sweep} for periodic collection of accumulated garbage.
  \item \textbf{Persistent heap} for Ideas that may be referenced across computations.
\end{enumerate}
\end{definition}

\begin{definition}[Root Set]
\label{def:gc-roots}
GC roots include:
\begin{itemize}
  \item All values on the operand stack
  \item All values in the current environment
  \item All values in call stack frames
  \item All handler entries (including captured continuations)
  \item The RPM state (acquired Ideas)
  \item Global environment
\end{itemize}
\end{definition}

\subsection{Idea Persistence}
\label{subsec:idea-persistence}

\begin{definition}[Idea Persistence Model]
\label{def:idea-persistence}
Ideas have special persistence semantics:
\begin{verbatim}
data IdeaPersistence
  = Ephemeral     -- Normal GC rules apply
  | Persistent    -- Survives GC until explicit release
  | Acquired      -- In RPM acquisition state (never collected)

markIdea :: Addr -> IdeaPersistence -> VM -> VM
markIdea addr pers vm =
  let obj = readHeap addr vm
  in writeHeap addr (obj { hoiPersistence = pers }) vm

-- Ideas in alpha are always Acquired
acquireIdea :: Addr -> VM -> VM
acquireIdea addr vm = markIdea addr Acquired vm
\end{verbatim}
\end{definition}

\begin{definition}[Continuation GC]
\label{def:continuation-gc}
Captured continuations require special handling:
\begin{itemize}
  \item Continuations hold references to stack frames and handlers.
  \item Unresumed continuations may leak memory.
  \item The VM tracks ``continuation debt'' and warns on excessive capture.
\end{itemize}
\begin{verbatim}
data ContinuationStatus
  = Active       -- May still be resumed
  | Resumed      -- Has been resumed (can collect saved state)
  | Abandoned    -- Explicitly abandoned (collect immediately)
\end{verbatim}
\end{definition}

%% ----------------------------------------------------------------------------
%% HANDOFF NOTE
%% ----------------------------------------------------------------------------
\begin{tcolorbox}[colback=green!5!white,colframe=green!50!black,title=\textbf{Handoff: COMPILER-TASK-05 Complete},breakable]
\textbf{Completed in this chapter:}
\begin{itemize}
  \item \textbf{VM Architecture} (Section~\ref{sec:vm-architecture}):
    \begin{itemize}
      \item Stack-based design with operand, call, and handler stacks
      \item Extended VM state with handler stack, ordinal register, quantum register
      \item Stack frame structure with continuation support
      \item Handler stack entry structure
      \item Memory model with heap objects for closures, handlers, continuations
      \item Value representation and sizing
    \end{itemize}
  \item \textbf{Instruction Set} (Section~\ref{sec:instruction-set}):
    \begin{itemize}
      \item Instruction encoding format (opcode + flags + operands)
      \item Core instructions: stack ops, variables, control flow, arithmetic
      \item TKS domain instructions: elements, foundations, noetics, fractals, ACBE
      \item RPM instructions: return, bind, check, acquire, fail, goal management
      \item Effect instructions (v7.3): install/remove handler, perform, resume, capture
      \item Transfinite instructions (v7.3): ordinal arithmetic, trans\_loop\_*
      \item Quantum instructions (v7.3): ket/bra, superpose, measure, entangle
      \item Closure and heap instructions
      \item Bytecode examples for noetic application, effect handling, transfinite loops
    \end{itemize}
  \item \textbf{RPM-Aware Execution} (Section~\ref{sec:rpm-aware-execution}):
    \begin{itemize}
      \item RPM state structure with acquisition set, goals, prerequisites, history
      \item Prerequisite graph encoding TKS progression rules
      \item RPM\_CHECK, RPM\_ACQUIRE, RPM\_BIND semantics
      \item Goal management instructions and example
    \end{itemize}
  \item \textbf{Effect Handling in VM} (Section~\ref{sec:effect-handling-vm}):
    \begin{itemize}
      \item Handler stack operations (push, pop, find)
      \item INSTALL\_HANDLER semantics
      \item Continuation structure and capture
      \item PERFORM\_EFFECT semantics with continuation capture
      \item RESUME and RESUME\_WITH semantics
      \item HANDLER\_RETURN semantics
      \item \textbf{Theorem~\ref{thm:vm-correctness-v73}}: VM execution preserves operational semantics (with proof)
    \end{itemize}
  \item \textbf{Transfinite Loop Execution} (Section~\ref{sec:transfinite-loop-vm}):
    \begin{itemize}
      \item Ordinal representation (Cantor Normal Form)
      \item Ordinal arithmetic in VM
      \item Transfinite loop state structure
      \item TRANS\_LOOP\_INIT, TRANS\_LOOP\_STEP, TRANS\_LOOP\_CHECK semantics
      \item Convergence detection criteria
    \end{itemize}
  \item \textbf{Garbage Collection} (Section~\ref{sec:garbage-collection}):
    \begin{itemize}
      \item Hybrid GC strategy (reference counting + mark-and-sweep)
      \item GC root set specification
      \item Idea persistence model (ephemeral, persistent, acquired)
      \item Continuation GC handling
    \end{itemize}
\end{itemize}

\textbf{Next tasks (COMPILER-TASK-06):}
\begin{itemize}
  \item Chapter 6: Module System
    \begin{itemize}
      \item Module syntax (declarations, imports, exports)
      \item Module types (signatures and structures)
      \item Functor modules (parameterized modules)
      \item Separate compilation (interface files, linking)
      \item Standard library (TKS.Core, TKS.Noetics, TKS.RPM, TKS.Fractals)
    \end{itemize}
\end{itemize}
\end{tcolorbox}

%% ============================================================================
%% CHAPTER 6: MODULE SYSTEM
%% ============================================================================
\chapter{Module System}
\label{ch:module-system}

\begin{tcolorbox}[colback=blue!5!white,colframe=blue!75!black,title=\vnewthree]
This chapter specifies the TKS module system for organizing code, managing namespaces, and separate compilation.
\end{tcolorbox}

%% ----------------------------------------------------------------------------
\section{Module Syntax}
\label{sec:module-syntax}

This section specifies the concrete syntax for TKS modules, extending the v7.2 module definitions with full support for hierarchical namespaces, selective imports, and effect exports.

\subsection{Module Declarations}
\label{subsec:module-declarations}

\begin{definition}[Module Declaration Syntax]
\label{def:module-decl-syntax}
The grammar for module declarations:
\begin{align}
\mathit{Module} ::=&\; \mathsf{module}\; \mathit{ModPath}\; \{\; \mathit{ModBody}\; \} \\
\mathit{ModPath} ::=&\; \mathit{Ident} \mid \mathit{ModPath}\,.\,\mathit{Ident} \\
\mathit{ModBody} ::=&\; \mathit{Export}?\; \mathit{Import}^*\; \mathit{Decl}^* \\
\mathit{Export} ::=&\; \mathsf{export}\; \{\; \mathit{ExportItem}^{+,}\; \} \\
\mathit{ExportItem} ::=&\; \mathit{Ident} \quad \text{(value)} \\
    &\mid \mathsf{type}\; \mathit{Ident} \quad \text{(abstract type)} \\
    &\mid \mathsf{type}\; \mathit{Ident}\; (=) \quad \text{(transparent type)} \\
    &\mid \mathsf{effect}\; \mathit{Ident} \quad \text{(effect export)} \\
    &\mid \mathsf{module}\; \mathit{Ident} \quad \text{(re-export submodule)}
\end{align}
\end{definition}

\begin{example}[Basic Module Declaration]
\label{ex:basic-module}
A module for TKS element utilities:
\begin{verbatim}
module TKS.Elements {
  export { Element, Wing, wing, index, successor, predecessor }

  -- Types
  type Wing = A | B | C | D
  type Element = Element(Wing, Int)

  -- Functions
  def wing : Element -> Wing
  def wing(Element(w, _)) = w

  def index : Element -> Int
  def index(Element(_, n)) = n

  def successor : Element -> RPM[Element]
  def successor(Element(w, n)) =
    if n < 10 then return Element(w, n + 1)
    else fail "Cannot exceed level 10"

  def predecessor : Element -> RPM[Element]
  def predecessor(Element(w, n)) =
    if n > 1 then return Element(w, n - 1)
    else fail "Cannot go below level 1"
}
\end{verbatim}
\end{example}

\subsection{Import Statements}
\label{subsec:import-statements}

\begin{definition}[Import Syntax]
\label{def:import-syntax}
Import statement grammar:
\begin{align}
\mathit{Import} ::=&\; \mathsf{import}\; \mathit{ModPath} \quad \text{(qualified import)} \\
    &\mid \mathsf{import}\; \mathit{ModPath}\; \mathsf{as}\; \mathit{Ident} \quad \text{(aliased import)} \\
    &\mid \mathsf{from}\; \mathit{ModPath}\; \mathsf{import}\; \mathit{ImportList} \quad \text{(selective import)} \\
    &\mid \mathsf{from}\; \mathit{ModPath}\; \mathsf{import}\; * \quad \text{(wildcard import)} \\
\mathit{ImportList} ::=&\; \{\; \mathit{ImportItem}^{+,}\; \} \\
\mathit{ImportItem} ::=&\; \mathit{Ident} \quad \text{(direct)} \\
    &\mid \mathit{Ident}\; \mathsf{as}\; \mathit{Ident} \quad \text{(renamed)} \\
    &\mid \mathsf{type}\; \mathit{Ident} \\
    &\mid \mathsf{effect}\; \mathit{Ident}
\end{align}
\end{definition}

\begin{example}[Import Variations]
\label{ex:import-variations}
\begin{verbatim}
-- Qualified import: use as TKS.Core.f
import TKS.Core

-- Aliased import: use as C.f
import TKS.Core as C

-- Selective import: brings nu_0, nu_1 into scope
from TKS.Noetics import { nu_0, nu_1, type Noetic }

-- Renamed import: use applyNoetic instead of apply
from TKS.Noetics import { apply as applyNoetic }

-- Wildcard import (discouraged): brings everything into scope
from TKS.Elements import *

-- Effect import
from TKS.RPM import { effect RPMEffect }
\end{verbatim}
\end{example}

\subsection{Qualified Names and Aliasing}
\label{subsec:qualified-names}

\begin{definition}[Name Resolution]
\label{def:name-resolution}
Name resolution follows these rules in order:
\begin{enumerate}
  \item \textbf{Local scope}: Names defined in current scope.
  \item \textbf{Selective imports}: Names explicitly imported via \texttt{from ... import}.
  \item \textbf{Enclosing modules}: Names from parent modules (for nested modules).
  \item \textbf{Qualified access}: Names via qualified path \texttt{M.N.x}.
\end{enumerate}
\end{definition}

\begin{definition}[Qualified Name Syntax]
\label{def:qualified-name}
Qualified names:
\begin{align}
\mathit{QualName} ::=&\; \mathit{Ident} \quad \text{(unqualified)} \\
    &\mid \mathit{ModPath}\,.\,\mathit{Ident} \quad \text{(qualified)}
\end{align}
Examples: \texttt{x}, \texttt{Core.f}, \texttt{TKS.Noetics.nu\_0}.
\end{definition}

\begin{definition}[Shadowing Rules]
\label{def:shadowing}
When names conflict:
\begin{itemize}
  \item Local definitions shadow imports.
  \item Selective imports shadow wildcard imports.
  \item Later imports shadow earlier imports (with warning).
  \item Qualified names are never shadowed.
\end{itemize}
\end{definition}

\subsection{Module Declarations}
\label{subsec:module-decl-forms}

\begin{definition}[Declaration Forms]
\label{def:decl-forms}
Declarations within a module body:
\begin{align}
\mathit{Decl} ::=&\; \mathsf{def}\; x : \tau = e \quad \text{(value definition)} \\
    &\mid \mathsf{def}\; x : \tau \quad \text{(external/abstract value)} \\
    &\mid \mathsf{type}\; T\; \bar{\alpha} = \tau \quad \text{(type definition)} \\
    &\mid \mathsf{type}\; T\; \bar{\alpha} \quad \text{(abstract type)} \\
    &\mid \mathsf{effect}\; E\; \{\; \bar{op}\; \} \quad \text{(effect definition)} \\
    &\mid \mathsf{handler}\; h : E \Rightarrow \tau\; \{\; \ldots\; \} \quad \text{(handler definition)} \\
    &\mid \mathsf{module}\; M\; \{\; \ldots\; \} \quad \text{(nested module)}
\end{align}
\end{definition}

%% ----------------------------------------------------------------------------
\section{Module Types}
\label{sec:module-types}

This section specifies module signatures (types), extending the v7.2 signature system with effect specifications and subtyping.

\subsection{Signature Declarations}
\label{subsec:signature-declarations}

\begin{definition}[Signature Syntax]
\label{def:signature-syntax}
Module signatures specify interfaces:
\begin{align}
\mathit{Signature} ::=&\; \mathsf{signature}\; S\; \{\; \mathit{SigBody}\; \} \\
\mathit{SigBody} ::=&\; \mathit{SigDecl}^* \\
\mathit{SigDecl} ::=&\; \mathsf{val}\; x : \tau \quad \text{(value specification)} \\
    &\mid \mathsf{type}\; T\; \bar{\alpha} \quad \text{(abstract type)} \\
    &\mid \mathsf{type}\; T\; \bar{\alpha} = \tau \quad \text{(manifest/transparent type)} \\
    &\mid \mathsf{effect}\; E\; \{\; op_i : A_i \to B_i \; \} \quad \text{(effect specification)} \\
    &\mid \mathsf{module}\; M : S \quad \text{(nested module spec)}
\end{align}
\end{definition}

\begin{example}[Signature Declaration]
\label{ex:signature-decl}
A signature for collections:
\begin{verbatim}
signature COLLECTION {
  type Elem              -- Abstract element type
  type T                 -- Abstract collection type

  val empty : T
  val singleton : Elem -> T
  val insert : Elem -> T -> T
  val member : Elem -> T -> Bool
  val fold : forall a. (Elem -> a -> a) -> a -> T -> a

  -- Effect specification for fallible operations
  effect CollectionError {
    notFound : Elem -> Unit
  }

  val find : Elem -> T -> RPM[Elem] ! CollectionError
}
\end{verbatim}
\end{example}

\subsection{Abstract vs Transparent Type Components}
\label{subsec:abstract-transparent}

\begin{definition}[Type Component Visibility]
\label{def:type-visibility}
Type components in signatures have two forms:
\begin{enumerate}
  \item \textbf{Abstract type}: $\mathsf{type}\; T$ --- implementation hidden from clients.
  \item \textbf{Transparent (manifest) type}: $\mathsf{type}\; T = \tau$ --- implementation exposed.
\end{enumerate}
\end{definition}

\begin{definition}[Type Abstraction Typing]
\label{def:type-abstraction-typing}
For module $M$ with signature $S$:
\begin{itemize}
  \item If $S$ declares $\mathsf{type}\; T$, then $M.T$ is abstract---clients cannot assume its structure.
  \item If $S$ declares $\mathsf{type}\; T = \tau$, then $M.T \equiv \tau$ in client code.
\end{itemize}
\end{definition}

\begin{example}[Abstract Type Enforcement]
\label{ex:abstract-type}
\begin{verbatim}
signature SET_SIG {
  type Elem
  type Set          -- Abstract: clients don't know it's a tree
  val empty : Set
  val add : Elem -> Set -> Set
}

module TreeSet : SET_SIG {
  type Elem = Element
  type Set = Tree[Element]  -- Implementation: balanced tree

  def empty = Leaf
  def add(e, s) = insert(e, s)
}

-- Client code:
-- TreeSet.Set is abstract; cannot pattern match on Tree
let s = TreeSet.add(A1, TreeSet.empty)  -- OK
-- case s of Leaf -> ... | Node(...) -> ...  -- ERROR: Set is abstract
\end{verbatim}
\end{example}

\subsection{Value and Effect Specifications}
\label{subsec:value-effect-specs}

\begin{definition}[Value Specification with Effects]
\label{def:value-spec-effects}
Value specifications include effect annotations:
\[
\mathsf{val}\; x : \tau \mathbin{!} \varepsilon
\]
where $\varepsilon$ is the effect row. Pure values use $\mathsf{val}\; x : \tau$ (implicit $\varepsilon = \emptyset$).
\end{definition}

\begin{definition}[Effect Specification]
\label{def:effect-spec}
Effect specifications declare effect signatures:
\begin{verbatim}
effect E {
  op_1 : A_1 -> B_1
  op_2 : A_2 -> B_2
  ...
}
\end{verbatim}
This declares effect $E$ with operations $op_i$ having input type $A_i$ and result type $B_i$.
\end{definition}

\subsection{Signature Matching and Subtyping}
\label{subsec:signature-matching}

\begin{definition}[Signature Matching]
\label{def:sig-matching}
Module $M$ \textbf{matches} signature $S$, written $M : S$, if:
\begin{enumerate}
  \item For each $\mathsf{val}\; x : \tau \mathbin{!} \varepsilon \in S$:
    \begin{itemize}
      \item $M$ exports $x$ with type $\tau'$ and effects $\varepsilon'$
      \item $\tau' \leq \tau$ (type subsumption)
      \item $\varepsilon' \subseteq \varepsilon$ (effect subsumption)
    \end{itemize}
  \item For each $\mathsf{type}\; T \in S$: $M$ defines type $T$.
  \item For each $\mathsf{type}\; T = \tau \in S$: $M$ defines $T = \tau$ exactly.
  \item For each $\mathsf{effect}\; E\; \{...\} \in S$: $M$ declares compatible effect $E$.
  \item For each $\mathsf{module}\; N : S' \in S$: $M.N : S'$.
\end{enumerate}
\end{definition}

\begin{definition}[Signature Subtyping]
\label{def:sig-subtyping}
Signature subtyping $S_1 <: S_2$:
\[
\frac{
  \begin{array}{c}
  \forall (\mathsf{val}\; x : \tau_2) \in S_2.\; \exists (\mathsf{val}\; x : \tau_1) \in S_1.\; \tau_1 \leq \tau_2 \\
  \forall (\mathsf{type}\; T) \in S_2.\; (\mathsf{type}\; T) \in S_1 \lor (\mathsf{type}\; T = \tau) \in S_1 \\
  \forall (\mathsf{type}\; T = \tau) \in S_2.\; (\mathsf{type}\; T = \tau) \in S_1
  \end{array}
}{S_1 <: S_2}
\]
A more specific signature (more exports, more transparent types) is a subtype of a less specific one.
\end{definition}

%% ----------------------------------------------------------------------------
\section{Functor Modules}
\label{sec:functor-modules}

This section specifies parameterized modules (functors), enabling generic module definitions.

\subsection{Functor Syntax}
\label{subsec:functor-syntax}

\begin{definition}[Functor Declaration]
\label{def:functor-decl}
Functors are modules parameterized by other modules:
\begin{align}
\mathit{Functor} ::=&\; \mathsf{functor}\; F\; (\; \mathit{Param}^{+,}\; ) : S\; \{\; \mathit{ModBody}\; \} \\
\mathit{Param} ::=&\; M : S \quad \text{(module parameter with signature)}
\end{align}
\end{definition}

\begin{example}[Functor Definition]
\label{ex:functor-def}
A generic set functor:
\begin{verbatim}
signature ORDERED {
  type T
  val compare : T -> T -> Ordering
}

signature SET {
  type Elem
  type Set
  val empty : Set
  val add : Elem -> Set -> Set
  val member : Elem -> Set -> Bool
  val toList : Set -> List[Elem]
}

functor MakeSet(Ord : ORDERED) : SET {
  type Elem = Ord.T
  type Set = BST[Elem]  -- Binary search tree

  def empty = Leaf

  def add(e, Leaf) = Node(Leaf, e, Leaf)
  def add(e, Node(l, v, r)) =
    case Ord.compare(e, v) of
      LT -> Node(add(e, l), v, r)
      EQ -> Node(l, v, r)
      GT -> Node(l, v, add(e, r))

  def member(e, Leaf) = false
  def member(e, Node(l, v, r)) =
    case Ord.compare(e, v) of
      LT -> member(e, l)
      EQ -> true
      GT -> member(e, r)

  def toList(t) = inorder(t)
}
\end{verbatim}
\end{example}

\subsection{Functor Application}
\label{subsec:functor-application}

\begin{definition}[Functor Application Syntax]
\label{def:functor-app}
Applying a functor to module arguments:
\begin{align}
\mathit{ModExpr} ::=&\; \mathit{ModPath} \quad \text{(module reference)} \\
    &\mid F\; (\; \mathit{ModExpr}^{+,}\; ) \quad \text{(functor application)}
\end{align}
\end{definition}

\begin{example}[Functor Application]
\label{ex:functor-app}
\begin{verbatim}
-- Define ordering for Elements
module ElementOrd : ORDERED {
  type T = Element

  def compare(Element(w1, n1), Element(w2, n2)) =
    case compareWing(w1, w2) of
      EQ -> compareInt(n1, n2)
      other -> other
}

-- Apply functor to create ElementSet module
module ElementSet = MakeSet(ElementOrd)

-- Usage
let s1 = ElementSet.add(A1, ElementSet.empty)
let s2 = ElementSet.add(B3, s1)
let hasA1 = ElementSet.member(A1, s2)  -- true
\end{verbatim}
\end{example}

\begin{definition}[Functor Application Typing]
\label{def:functor-app-typing}
\[
\frac{F : (M_1 : S_1, \ldots, M_n : S_n) \to S \quad \forall i.\; N_i : S_i}
     {F(N_1, \ldots, N_n) : S[N_1/M_1, \ldots, N_n/M_n]}
\]
The result signature has module parameters substituted with actual arguments.
\end{definition}

\subsection{Higher-Order Functors}
\label{subsec:higher-order-functors}

\begin{definition}[Higher-Order Functor]
\label{def-ho-functor}
TKS supports \textbf{higher-order functors}---functors that take functors as arguments or return functors:
\begin{verbatim}
-- Functor that takes a functor as argument
functor Compose(F : (X : S1) -> S2, G : (Y : S2) -> S3)
  : (Z : S1) -> S3
{
  functor Result(Z : S1) : S3 = G(F(Z))
}

-- Example: Composing Set with a transformation functor
functor MapSet(F : (X : ORDERED) -> ORDERED) : (Y : ORDERED) -> SET {
  functor Result(Y : ORDERED) : SET = MakeSet(F(Y))
}
\end{verbatim}
\end{definition}

\begin{definition}[Functor Kind System]
\label{def:functor-kinds}
Functors are classified by kinds:
\begin{align}
\mathit{Kind} ::=&\; \mathsf{Mod} \quad \text{(module kind)} \\
    &\mid S \to K \quad \text{(functor kind)} \\
    &\mid (S_1, \ldots, S_n) \to K \quad \text{(multi-param functor kind)}
\end{align}
where $S$ ranges over signatures and $K$ over kinds.
\end{definition}

\subsection{Applicative vs Generative Functors}
\label{subsec:applicative-generative}

\begin{definition}[Functor Semantics]
\label{def:functor-semantics}
TKS supports two functor semantics:
\begin{enumerate}
  \item \textbf{Applicative} (default): $F(M) \equiv F(M)$ --- applying the same functor to the same argument yields equivalent modules.
  \item \textbf{Generative}: Each application generates fresh abstract types.
\end{enumerate}
\begin{verbatim}
-- Applicative (default): types are shared
module S1 = MakeSet(IntOrd)
module S2 = MakeSet(IntOrd)
-- S1.Set === S2.Set (same type)

-- Generative: fresh types each application
generative functor MakeFreshSet(...) : SET { ... }
module S3 = MakeFreshSet(IntOrd)
module S4 = MakeFreshSet(IntOrd)
-- S3.Set =/= S4.Set (different types)
\end{verbatim}
\end{definition}

%% ----------------------------------------------------------------------------
\section{Separate Compilation}
\label{sec:separate-compilation}

This section specifies the separate compilation model for TKS programs.

\subsection{File Organization}
\label{subsec:file-organization}

\begin{definition}[TKS File Types]
\label{def:file-types}
TKS uses two primary file types:
\begin{center}
\begin{tabular}{lp{9cm}}
\toprule
\textbf{Extension} & \textbf{Description} \\
\midrule
\texttt{.tks} & Implementation file: contains full module definitions with code \\
\texttt{.tksi} & Interface file: contains module signature (auto-generated or handwritten) \\
\bottomrule
\end{tabular}
\end{center}
\end{definition}

\begin{definition}[File-to-Module Mapping]
\label{def:file-module-mapping}
The module path maps to filesystem paths:
\begin{itemize}
  \item Module \texttt{TKS.Core.Elements} corresponds to:
    \begin{itemize}
      \item Implementation: \texttt{TKS/Core/Elements.tks}
      \item Interface: \texttt{TKS/Core/Elements.tksi}
    \end{itemize}
  \item The compiler searches paths in order: current directory, then \texttt{TKS\_PATH} directories.
\end{itemize}
\end{definition}

\subsection{Interface Files}
\label{subsec:interface-files}

\begin{definition}[Interface File Content]
\label{def:interface-content}
An interface file (\texttt{.tksi}) contains:
\begin{enumerate}
  \item Module signature declarations
  \item Exported type definitions (transparent types fully expanded)
  \item Exported value type signatures (with effects)
  \item Exported effect declarations
  \item Dependency list (imported modules)
\end{enumerate}
\end{definition}

\begin{example}[Interface File]
\label{ex:interface-file}
File \texttt{TKS/Noetics.tksi}:
\begin{verbatim}
-- Auto-generated interface for TKS.Noetics
-- Compiler version: tks-7.3.0
-- Generated: 2024-01-15

interface TKS.Noetics {
  -- Dependencies
  requires TKS.Core
  requires TKS.Elements

  -- Types
  type Noetic = Noetic(Int)            -- Transparent
  type NoeticLevel                     -- Abstract

  -- Values
  val nu_0 : Noetic
  val nu_1 : Noetic
  val nu_2 : Noetic
  val nu_3 : Noetic

  val apply : Noetic -> Idea -> Idea
  val compose : Noetic -> Noetic -> Noetic
  val level : Noetic -> NoeticLevel

  val descendChain : List[Noetic] -> Idea -> RPM[Idea] ! RPMEffect

  -- Effects
  effect NoeticError {
    invalidLevel : Int -> Unit
    compositionFailed : Noetic -> Noetic -> Unit
  }
}
\end{verbatim}
\end{example}

\begin{definition}[Interface Stability]
\label{def:interface-stability}
An interface is \textbf{stable} if:
\begin{enumerate}
  \item All exported types have determined representations.
  \item All value types are fully resolved (no unresolved type variables at module level).
  \item Effect signatures are complete.
\end{enumerate}
Only stable interfaces enable separate compilation.
\end{definition}

\subsection{Implementation Files}
\label{subsec:implementation-files}

\begin{definition}[Implementation File Content]
\label{def:implementation-content}
An implementation file (\texttt{.tks}) contains:
\begin{enumerate}
  \item Module declaration with optional signature ascription
  \item Import statements
  \item Type definitions
  \item Value definitions with implementations
  \item Effect and handler definitions
  \item Nested modules and functors
\end{enumerate}
\end{definition}

\begin{example}[Implementation File]
\label{ex:implementation-file}
File \texttt{TKS/Noetics.tks}:
\begin{verbatim}
module TKS.Noetics : TKS.Noetics {  -- Ascribe to interface
  export {
    type Noetic(=), type NoeticLevel,
    nu_0, nu_1, nu_2, nu_3,
    apply, compose, level, descendChain,
    effect NoeticError
  }

  import TKS.Core
  from TKS.Elements import { Element, Idea }

  -- Type definitions
  type Noetic = Noetic(Int)
  type NoeticLevel = Level(Int)  -- Internal representation

  -- Noetic constants
  def nu_0 : Noetic = Noetic(0)
  def nu_1 : Noetic = Noetic(1)
  def nu_2 : Noetic = Noetic(2)
  def nu_3 : Noetic = Noetic(3)

  -- Core operations
  def level(Noetic(k)) = Level(k)

  def apply(Noetic(k), idea) =
    transformIdea(k, idea)  -- Internal implementation

  def compose(Noetic(j), Noetic(k)) =
    Noetic(j + k)  -- Composition adds levels

  -- Effect definition
  effect NoeticError {
    invalidLevel : Int -> Unit
    compositionFailed : Noetic -> Noetic -> Unit
  }

  -- Effectful operation
  def descendChain(noetics, idea) =
    foldM(\nu, i -> apply(nu, i), idea, noetics)
}
\end{verbatim}
\end{example}

\subsection{Compilation Pipeline}
\label{subsec:compilation-pipeline}

\begin{definition}[Separate Compilation Steps]
\label{def:compilation-steps}
The compilation pipeline for a module $M$:
\begin{enumerate}
  \item \textbf{Parse}: Read \texttt{M.tks}, produce AST.
  \item \textbf{Resolve dependencies}: For each import, load \texttt{.tksi} or compile dependency.
  \item \textbf{Type check}: Type check against loaded interfaces.
  \item \textbf{Generate interface}: Produce \texttt{M.tksi} from exports.
  \item \textbf{Compile}: Generate bytecode \texttt{M.tkso}.
\end{enumerate}
\end{definition}

\begin{definition}[Compilation Artifacts]
\label{def:artifacts}
Compilation produces:
\begin{center}
\begin{tabular}{llp{7cm}}
\toprule
\textbf{File} & \textbf{Extension} & \textbf{Content} \\
\midrule
Interface & \texttt{.tksi} & Module signature, exported types \\
Object & \texttt{.tkso} & Compiled bytecode \\
Debug & \texttt{.tksd} & Source maps, debug info \\
\bottomrule
\end{tabular}
\end{center}
\end{definition}

\subsection{Linking and Dependency Resolution}
\label{subsec:linking}

\begin{definition}[Dependency Graph]
\label{def:dependency-graph}
The dependency graph $G = (V, E)$ where:
\begin{itemize}
  \item $V$ = set of modules
  \item $(M, N) \in E$ iff $M$ imports $N$
\end{itemize}
Compilation order is a topological sort of $G$. Cycles are an error.
\end{definition}

\begin{definition}[Linking Process]
\label{def:linking-process}
The linker combines object files:
\begin{enumerate}
  \item \textbf{Symbol resolution}: Match imports to exports across modules.
  \item \textbf{Type compatibility check}: Verify interface consistency.
  \item \textbf{Address binding}: Assign runtime addresses to symbols.
  \item \textbf{Code merging}: Combine bytecode into executable.
\end{enumerate}
\end{definition}

\begin{definition}[Incremental Compilation]
\label{def:incremental}
A module $M$ needs recompilation iff:
\begin{itemize}
  \item $M.tks$ has changed since last compile, OR
  \item Any dependency $N$ has a newer $N.tksi$ than $M.tkso$
\end{itemize}
The compiler tracks file timestamps and interface hashes for incremental builds.
\end{definition}

%% ----------------------------------------------------------------------------
\section{Standard Library}
\label{sec:standard-library}

This section specifies the TKS standard library modules.

\subsection{Library Overview}
\label{subsec:library-overview}

\begin{definition}[Standard Library Structure]
\label{def:stdlib-structure}
The TKS standard library is organized hierarchically:
\begin{verbatim}
TKS
|-- Core          -- Basic types, functions, effects
|-- Elements      -- TKS Element types and operations
|-- Noetics       -- Noetic operators and algebra
|-- RPM           -- Requisite Progression Model
|-- Fractals      -- Fractal structures and ACBE
|-- Quantum       -- Quantum Idea states
|-- Collections   -- Lists, Sets, Maps
|-- IO            -- Input/Output effects
|-- Concurrency   -- Parallel execution
\end{verbatim}
\end{definition}

\subsection{TKS.Core}
\label{subsec:tks-core}

\begin{definition}[TKS.Core Interface]
\label{def:tks-core}
The core module provides fundamental types and operations:
\begin{verbatim}
interface TKS.Core {
  -- Basic types
  type Unit = ()
  type Bool = True | False
  type Int                           -- Built-in
  type String                        -- Built-in
  type Option[a] = None | Some(a)
  type Result[a, e] = Ok(a) | Err(e)
  type List[a] = Nil | Cons(a, List[a])
  type Pair[a, b] = Pair(a, b)

  -- Basic operations
  val not : Bool -> Bool
  val (&&) : Bool -> Bool -> Bool
  val (||) : Bool -> Bool -> Bool

  val (+) : Int -> Int -> Int
  val (-) : Int -> Int -> Int
  val (*) : Int -> Int -> Int
  val (/) : Int -> Int -> Option[Int]
  val (%) : Int -> Int -> Option[Int]
  val (<) : Int -> Int -> Bool
  val (<=) : Int -> Int -> Bool
  val (==) : forall a. a -> a -> Bool

  -- List operations
  val map : forall a b. (a -> b) -> List[a] -> List[b]
  val filter : forall a. (a -> Bool) -> List[a] -> List[a]
  val fold : forall a b. (a -> b -> b) -> b -> List[a] -> b
  val length : forall a. List[a] -> Int
  val concat : forall a. List[a] -> List[a] -> List[a]

  -- Option operations
  val mapOption : forall a b. (a -> b) -> Option[a] -> Option[b]
  val flatMapOption : forall a b. (a -> Option[b]) -> Option[a] -> Option[b]
  val getOrElse : forall a. a -> Option[a] -> a

  -- Function composition
  val compose : forall a b c. (b -> c) -> (a -> b) -> (a -> c)
  val identity : forall a. a -> a
  val const : forall a b. a -> b -> a
}
\end{verbatim}
\end{definition}

\subsection{TKS.Noetics}
\label{subsec:tks-noetics}

\begin{definition}[TKS.Noetics Interface]
\label{def:tks-noetics}
The noetics module provides noetic operators:
\begin{verbatim}
interface TKS.Noetics {
  requires TKS.Core
  requires TKS.Elements

  -- Types
  type Noetic
  type NoeticChain = List[Noetic]

  -- Noetic constants (nu_0 through nu_3)
  val nu_0 : Noetic    -- Identity/preservation
  val nu_1 : Noetic    -- First-level transformation
  val nu_2 : Noetic    -- Second-level transformation
  val nu_3 : Noetic    -- Third-level transformation

  -- Core operations
  val apply : Noetic -> Idea -> Idea
  val compose : Noetic -> Noetic -> Noetic
  val inverse : Noetic -> Option[Noetic]
  val level : Noetic -> Int

  -- Algebraic properties
  val isIdentity : Noetic -> Bool
  val isInvertible : Noetic -> Bool

  -- Chain operations
  val applyChain : NoeticChain -> Idea -> Idea
  val composeChain : NoeticChain -> Noetic

  -- Noetic algebra
  val commutator : Noetic -> Noetic -> Noetic  -- [nu_j, nu_k]

  -- Effect for noetic errors
  effect NoeticError {
    levelMismatch : Int -> Int -> Unit
    nonInvertible : Noetic -> Unit
  }
}
\end{verbatim}
\end{definition}

\subsection{TKS.RPM}
\label{subsec:tks-rpm}

\begin{definition}[TKS.RPM Interface]
\label{def:tks-rpm}
The RPM module provides the Requisite Progression Model:
\begin{verbatim}
interface TKS.RPM {
  requires TKS.Core
  requires TKS.Elements

  -- The RPM monad
  type RPM[a]

  -- Monad operations
  val return : forall a. a -> RPM[a]
  val bind : forall a b. RPM[a] -> (a -> RPM[b]) -> RPM[b]
  val fail : forall a. String -> RPM[a]
  val (>>=) : forall a b. RPM[a] -> (a -> RPM[b]) -> RPM[b]
  val (>>) : forall a b. RPM[a] -> RPM[b] -> RPM[b]

  -- RPM operations
  val check : Element -> RPM[Unit]
  val acquire : Element -> RPM[Unit]
  val acquired : Element -> RPM[Bool]
  val prerequisites : Element -> RPM[List[Element]]

  -- Acquisition state
  type AcquisitionState
  val getState : RPM[AcquisitionState]
  val withState : forall a. AcquisitionState -> RPM[a] -> RPM[a]

  -- Goal planning
  type Goal
  val setGoal : Element -> RPM[Goal]
  val planPath : Goal -> RPM[List[Element]]
  val executeGoal : Goal -> RPM[Unit]

  -- Running RPM computations
  val runRPM : forall a. RPM[a] -> AcquisitionState -> Result[a, RPMError]
  val evalRPM : forall a. RPM[a] -> Result[a, RPMError]

  -- Error type
  type RPMError =
    | PrerequisiteNotMet(Element, List[Element])
    | AcquisitionFailed(Element, String)
    | GoalUnreachable(Element)

  -- Effect signature
  effect RPMEffect {
    checkPrereq : Element -> Bool
    doAcquire : Element -> Unit
    getAcquired : Unit -> List[Element]
  }
}
\end{verbatim}
\end{definition}

\subsection{TKS.Fractals}
\label{subsec:tks-fractals}

\begin{definition}[TKS.Fractals Interface]
\label{def:tks-fractals}
The fractals module provides fractal structures and ACBE:
\begin{verbatim}
interface TKS.Fractals {
  requires TKS.Core
  requires TKS.Elements
  requires TKS.Noetics

  -- Fractal type
  type Fractal[a]
  type FractalLevel = Int

  -- Construction
  val singleton : forall a. a -> Fractal[a]
  val fromLevels : forall a. List[(FractalLevel, a)] -> Fractal[a]
  val extend : forall a. Fractal[a] -> FractalLevel -> a -> Fractal[a]

  -- Access
  val atLevel : forall a. Fractal[a] -> FractalLevel -> Option[a]
  val levels : forall a. Fractal[a] -> List[FractalLevel]
  val depth : forall a. Fractal[a] -> Int

  -- Transformation
  val mapFractal : forall a b. (a -> b) -> Fractal[a] -> Fractal[b]
  val foldFractal : forall a b. (FractalLevel -> a -> b -> b) -> b -> Fractal[a] -> b

  -- ACBE (Abstracting Concrete from Below to Everywhere)
  type ACBEState[a]

  val acbeInit : Idea -> ACBEState[Idea]
  val acbeStep : ACBEState[Idea] -> ACBEState[Idea]
  val acbeComplete : ACBEState[Idea] -> Bool
  val acbeResult : ACBEState[Idea] -> Fractal[Idea]
  val acbeRun : Idea -> Fractal[Idea]

  -- Fractal algebra
  val merge : forall a. Fractal[a] -> Fractal[a] -> Fractal[a]
  val project : forall a. Fractal[a] -> FractalLevel -> Fractal[a]

  -- Effect
  effect FractalError {
    levelOutOfBounds : FractalLevel -> Unit
    incompatibleFractals : Unit
  }
}
\end{verbatim}
\end{definition}

\subsection{TKS.Quantum}
\label{subsec:tks-quantum}

\begin{definition}[TKS.Quantum Interface]
\label{def:tks-quantum}
The quantum module provides quantum Idea states:
\begin{verbatim}
interface TKS.Quantum {
  requires TKS.Core
  requires TKS.Elements
  requires TKS.Noetics

  -- Quantum state types
  type QState                        -- Quantum state
  type Ket[a]                        -- Ket vector |a>
  type Bra[a]                        -- Bra vector <a|
  type Complex                       -- Complex number

  -- State construction
  val ket : forall a. a -> Ket[a]
  val bra : forall a. a -> Bra[a]
  val superpose : forall a. List[(Complex, Ket[a])] -> QState
  val tensorProduct : QState -> QState -> QState

  -- Inner products and measurement
  val braket : forall a. Bra[a] -> Ket[a] -> Complex
  val probability : forall a. a -> QState -> Real
  val measure : QState -> RPM[Element] ! QuantumEffect
  val collapse : QState -> Element -> QState

  -- Operators
  type QOperator
  val applyOperator : QOperator -> QState -> QState
  val noeticOperator : Noetic -> QOperator
  val projector : Element -> QOperator

  -- World projections
  val projectWorld : Wing -> QState -> QState
  val worldProbability : Wing -> QState -> Real

  -- Entanglement
  val entangle : QState -> QState -> QState
  val isEntangled : QState -> Bool
  val partialTrace : QState -> Int -> QState

  -- Effect
  effect QuantumEffect {
    measureState : QState -> Element
    randomPhase : Unit -> Complex
  }

  -- Complex number operations
  val complex : Real -> Real -> Complex
  val magnitude : Complex -> Real
  val phase : Complex -> Real
  val conjugate : Complex -> Complex
}
\end{verbatim}
\end{definition}

\subsection{Module Soundness}
\label{subsec:module-soundness}

\begin{theorem}[Module System Type Soundness]
\label{thm:module-soundness}
The TKS module system preserves type soundness across module boundaries. Specifically:

\begin{enumerate}
  \item \textbf{Signature Abstraction}: If module $M : S$ and client code $C$ only accesses $M$ through signature $S$, then replacing $M$ with any $M' : S$ preserves type safety of $C$.

  \item \textbf{Separate Compilation Coherence}: If modules $M_1, \ldots, M_n$ are compiled separately against interfaces $I_1, \ldots, I_n$, and linked into program $P$, then $P$ is type-safe iff:
  \begin{enumerate}
    \item Each $M_i$ type-checks against its declared imports.
    \item For each import edge $M_i \to M_j$, interface $I_j$ matches the interface $M_i$ was compiled against.
  \end{enumerate}

  \item \textbf{Functor Type Preservation}: For functor $F : (X : S) \to T$ and module $M : S$:
  \[
  \Gamma \vdash M : S \implies \Gamma \vdash F(M) : T[M/X]
  \]

  \item \textbf{Effect Boundary Safety}: If $M$ exports $f : \tau \mathbin{!} \varepsilon$ and client imports $f$, then:
  \begin{enumerate}
    \item Client must handle or propagate effects in $\varepsilon$.
    \item No unhandled effects can escape module boundaries.
  \end{enumerate}
\end{enumerate}
\end{theorem}

\begin{proof}
We prove each part:

\textbf{Part 1 (Signature Abstraction):}
Let $M : S$ with $S$ declaring abstract type $T$. Client $C$ can only use operations on $T$ specified in $S$. If $M' : S$, then $M'$ provides the same operations with the same types. By the substitution lemma for modules, $C[M'/M]$ remains well-typed.

For transparent types, $S$ exposes the definition, so clients may depend on it. Signature subtyping ($S' <: S$) ensures transparency is preserved.

\textbf{Part 2 (Separate Compilation Coherence):}
The interface $I_j$ records the types and effects of $M_j$'s exports. If $M_i$ imports $M_j$ and is compiled against $I_j$, the type checker verifies:
\begin{itemize}
  \item All uses of $M_j.x$ in $M_i$ are consistent with $I_j$'s declaration of $x$.
\end{itemize}
At link time, we verify $M_j$'s actual interface matches $I_j$. If so, runtime behavior matches compile-time assumptions.

If interfaces mismatch (e.g., $M_j$ was recompiled with different types), linking fails with an interface compatibility error.

\textbf{Part 3 (Functor Type Preservation):}
By induction on functor typing derivation.

\emph{Base case}: For functor $F$ with body $B$:
\[
\frac{X : S \vdash B : T}{\vdash F : (X : S) \to T}
\]
Given $M : S$, substituting $M$ for $X$ in $B$ yields $B[M/X]$.
By the substitution lemma: if $X : S \vdash B : T$ and $\vdash M : S$, then $\vdash B[M/X] : T[M/X]$.

\emph{Inductive case}: For higher-order functors, apply the base case recursively.

\textbf{Part 4 (Effect Boundary Safety):}
The type system tracks effects. If $M$ exports $f : \tau \mathbin{!} \varepsilon$:
\begin{enumerate}
  \item The export check verifies $f$'s implementation has effects $\subseteq \varepsilon$.
  \item Client importing $f$ gets type $f : \tau \mathbin{!} \varepsilon$ in its context.
  \item Effect typing rules require client to either:
    \begin{itemize}
      \item Install handlers for effects in $\varepsilon$, or
      \item Declare those effects in its own signature.
    \end{itemize}
\end{enumerate}
At module boundaries, the ``no unhandled effects'' check verifies the main entry point has empty effect row, ensuring all effects are handled.
\end{proof}

%% ----------------------------------------------------------------------------
%% HANDOFF NOTE
%% ----------------------------------------------------------------------------
\begin{tcolorbox}[colback=green!5!white,colframe=green!50!black,title=\textbf{Handoff: COMPILER-TASK-06 Complete},breakable]
\textbf{Completed in this chapter:}
\begin{itemize}
  \item \textbf{Module Syntax} (Section~\ref{sec:module-syntax}):
    \begin{itemize}
      \item Module declaration grammar with ModPath, ModBody, Export, Import
      \item Import variations: qualified, aliased, selective, wildcard, renamed
      \item Name resolution rules and shadowing
      \item Declaration forms: def, type, effect, handler, nested modules
    \end{itemize}
  \item \textbf{Module Types/Signatures} (Section~\ref{sec:module-types}):
    \begin{itemize}
      \item Signature declaration syntax
      \item Abstract vs transparent (manifest) type components
      \item Value and effect specifications
      \item Signature matching rules
      \item Signature subtyping
    \end{itemize}
  \item \textbf{Functor Modules} (Section~\ref{sec:functor-modules}):
    \begin{itemize}
      \item Functor declaration syntax with parameters and signatures
      \item Functor application syntax and typing rule
      \item Higher-order functors
      \item Functor kind system
      \item Applicative vs generative functor semantics
    \end{itemize}
  \item \textbf{Separate Compilation} (Section~\ref{sec:separate-compilation}):
    \begin{itemize}
      \item File types: .tks (implementation), .tksi (interface), .tkso (object)
      \item File-to-module path mapping
      \item Interface file content and stability
      \item Implementation file structure
      \item Compilation pipeline (parse, resolve, typecheck, generate, compile)
      \item Dependency graph and linking process
      \item Incremental compilation
    \end{itemize}
  \item \textbf{Standard Library} (Section~\ref{sec:standard-library}):
    \begin{itemize}
      \item Library hierarchy overview
      \item TKS.Core: basic types (Unit, Bool, Int, Option, Result, List), operations
      \item TKS.Noetics: Noetic type, nu\_0--nu\_3, apply, compose, chains
      \item TKS.RPM: RPM monad, check/acquire, goals, AcquisitionState, RPMEffect
      \item TKS.Fractals: Fractal type, ACBE operations, fractal algebra
      \item TKS.Quantum: QState, Ket/Bra, superposition, measurement, entanglement
    \end{itemize}
  \item \textbf{Theorem~\ref{thm:module-soundness}}: Module system type soundness (with proof)
    \begin{itemize}
      \item Signature abstraction
      \item Separate compilation coherence
      \item Functor type preservation
      \item Effect boundary safety
    \end{itemize}
\end{itemize}

\textbf{Next tasks (COMPILER-TASK-07):}
\begin{itemize}
  \item Chapter 7: Foreign Function Interface
    \begin{itemize}
      \item External function declarations
      \item Marshalling TKS values to/from foreign types
      \item Safety and effect tracking for FFI calls
      \item Platform-specific bindings
    \end{itemize}
\end{itemize}
\end{tcolorbox}

%% ============================================================================
%% CHAPTER 7: FOREIGN FUNCTION INTERFACE
%% ============================================================================
\chapter{Foreign Function Interface}
\label{ch:ffi}

\begin{tcolorbox}[colback=blue!5!white,colframe=blue!75!black,title=\vnewthree]
This chapter specifies how TKS code can interoperate with external languages and systems.
\end{tcolorbox}

%% ----------------------------------------------------------------------------
\section{FFI Declarations}
\label{sec:ffi-declarations}

This section specifies the syntax and semantics for declaring and using foreign functions in TKS.

\subsection{External Function Declarations}
\label{subsec:external-declarations}

\begin{definition}[External Declaration Syntax]
\label{def:external-syntax}
Grammar for foreign function declarations:
\begin{align}
\mathit{ExternDecl} ::=&\; \mathsf{external}\; \mathit{Convention}\; \mathit{Safety}?\; \mathsf{fn}\; \mathit{Ident} \\
    &\quad (\; \mathit{Param}^{*,}\; ) : \mathit{ForeignType}\; \mathit{EffectAnn}? \\
\mathit{Convention} ::=&\; \texttt{"C"} \mid \texttt{"stdcall"} \mid \texttt{"fastcall"} \mid \texttt{"system"} \\
\mathit{Safety} ::=&\; \mathsf{safe} \mid \mathsf{unsafe} \\
\mathit{Param} ::=&\; \mathit{Ident} : \mathit{ForeignType} \\
\mathit{EffectAnn} ::=&\; \mathbin{!} \mathit{EffectRow}
\end{align}
\end{definition}

\begin{definition}[Calling Conventions]
\label{def:calling-conventions}
TKS supports the following calling conventions:
\begin{center}
\begin{tabular}{lp{9cm}}
\toprule
\textbf{Convention} & \textbf{Description} \\
\midrule
\texttt{"C"} & Standard C calling convention (cdecl); default \\
\texttt{"stdcall"} & Windows stdcall (callee cleans stack) \\
\texttt{"fastcall"} & Register-based fast calling convention \\
\texttt{"system"} & Platform default (C on Unix, stdcall on Windows) \\
\bottomrule
\end{tabular}
\end{center}
\end{definition}

\begin{example}[Basic External Declarations]
\label{ex:extern-decl}
\begin{verbatim}
module TKS.FFI.Libc {
  -- Memory allocation (unsafe, returns raw pointer)
  external "C" unsafe fn malloc(size: CSize) : Ptr ! IOEffect
  external "C" unsafe fn free(ptr: Ptr) : Unit ! IOEffect
  external "C" unsafe fn realloc(ptr: Ptr, size: CSize) : Ptr ! IOEffect

  -- String operations (safe wrappers)
  external "C" safe fn strlen(s: CString) : CSize ! IOEffect
  external "C" safe fn strcmp(s1: CString, s2: CString) : CInt ! IOEffect

  -- Math functions (pure, no effects)
  external "C" safe fn sin(x: CDouble) : CDouble
  external "C" safe fn cos(x: CDouble) : CDouble
  external "C" safe fn sqrt(x: CDouble) : CDouble

  -- File I/O
  external "C" safe fn fopen(path: CString, mode: CString) : FilePtr ! IOEffect
  external "C" safe fn fclose(fp: FilePtr) : CInt ! IOEffect
  external "C" safe fn fread(buf: Ptr, size: CSize, n: CSize, fp: FilePtr)
    : CSize ! IOEffect
}
\end{verbatim}
\end{example}

\subsection{Library Linking}
\label{subsec:library-linking}

\begin{definition}[Link Specification]
\label{def:link-spec}
Specifying external libraries to link:
\begin{verbatim}
-- Link directives at module level
link "m"        -- Link libm (math library)
link "pthread"  -- Link pthread
link "ssl"      -- Link OpenSSL

-- Platform-specific linking
#[cfg(target_os = "windows")]
link "kernel32"

#[cfg(target_os = "linux")]
link "dl"
\end{verbatim}
\end{definition}

\begin{definition}[Symbol Resolution]
\label{def:symbol-resolution}
External symbols are resolved as follows:
\begin{enumerate}
  \item \textbf{Name mapping}: By default, TKS function name maps directly to symbol.
  \item \textbf{Explicit symbol}: Use \texttt{@symbol("name")} to specify different symbol.
  \item \textbf{Name mangling}: No mangling for C convention; platform-specific for others.
\end{enumerate}
\begin{verbatim}
-- Explicit symbol name
external "C" safe fn tks_print(s: CString) : Unit ! IOEffect
  @symbol("printf")

-- Windows API with different name
external "stdcall" safe fn createFile(...) : Handle ! IOEffect
  @symbol("CreateFileW")
\end{verbatim}
\end{definition}

\subsection{Variadic Functions}
\label{subsec:variadic}

\begin{definition}[Variadic External Functions]
\label{def:variadic-extern}
C variadic functions require special handling:
\begin{verbatim}
-- Printf-style variadic function
external "C" unsafe fn printf(fmt: CString, ...) : CInt ! IOEffect

-- Safe wrapper with known argument types
def printInt(n: Int) : Unit ! IOEffect =
  let _ = printf("%d\n", toCInt(n))
  ()

def printStr(s: String) : Unit ! IOEffect =
  let _ = printf("%s\n", toCString(s))
  ()
\end{verbatim}
Variadic calls are always marked \texttt{unsafe} because argument types cannot be statically verified.
\end{definition}

%% ----------------------------------------------------------------------------
\section{Type Marshalling}
\label{sec:type-marshalling}

This section specifies how TKS types are converted to and from foreign type representations.

\subsection{Foreign Type System}
\label{subsec:foreign-types}

\begin{definition}[Foreign Types]
\label{def:foreign-types}
The foreign type system for FFI interoperability:
\begin{align}
\mathit{ForeignType} ::=&\; \mathsf{CInt} \mid \mathsf{CUInt} \mid \mathsf{CLong} \mid \mathsf{CULong} \\
    &\mid \mathsf{CShort} \mid \mathsf{CUShort} \\
    &\mid \mathsf{CChar} \mid \mathsf{CUChar} \\
    &\mid \mathsf{CFloat} \mid \mathsf{CDouble} \\
    &\mid \mathsf{CSize} \mid \mathsf{CSSize} \\
    &\mid \mathsf{CBool} \\
    &\mid \mathsf{Ptr} \mid \mathsf{Ptr}[\tau] \quad \text{(pointer types)} \\
    &\mid \mathsf{CString} \quad \text{(null-terminated string)} \\
    &\mid \mathsf{CArray}[\tau, n] \quad \text{(fixed-size array)} \\
    &\mid \mathsf{FunPtr}[\tau_1 \to \tau_2] \quad \text{(function pointer)} \\
    &\mid \mathsf{Struct}[S] \quad \text{(C struct reference)}
\end{align}
\end{definition}

\begin{definition}[Platform-Dependent Sizes]
\label{def:platform-sizes}
Foreign type sizes vary by platform:
\begin{center}
\begin{tabular}{lccc}
\toprule
\textbf{Type} & \textbf{32-bit} & \textbf{64-bit} & \textbf{C Equivalent} \\
\midrule
\textsf{CInt} & 32 bits & 32 bits & \texttt{int} \\
\textsf{CLong} & 32 bits & 64 bits* & \texttt{long} \\
\textsf{CSize} & 32 bits & 64 bits & \texttt{size\_t} \\
\textsf{Ptr} & 32 bits & 64 bits & \texttt{void*} \\
\bottomrule
\end{tabular}
\end{center}
*On Windows 64-bit, \textsf{CLong} remains 32 bits.
\end{definition}

\subsection{Automatic Marshalling}
\label{subsec:auto-marshalling}

\begin{definition}[Primitive Type Mapping]
\label{def:primitive-mapping}
Automatic marshalling for TKS primitives:
\begin{center}
\begin{tabular}{lll}
\toprule
\textbf{TKS Type} & \textbf{Foreign Type} & \textbf{Conversion} \\
\midrule
\textsf{Int} & \textsf{CInt} / \textsf{CLong} & Direct (with bounds check) \\
\textsf{Bool} & \textsf{CBool} / \textsf{CInt} & 0/1 mapping \\
\textsf{Unit} & \textsf{void} & No value \\
\textsf{String} & \textsf{CString} & UTF-8 encoding + null terminator \\
\textsf{Option}[$\tau$] & \textsf{Ptr}[$\tau$] & \textsf{None} $\mapsto$ null \\
\bottomrule
\end{tabular}
\end{center}
\end{definition}

\begin{definition}[Marshalling Functions]
\label{def:marshalling-fns}
Built-in marshalling operations:
\begin{verbatim}
interface TKS.FFI.Marshal {
  -- Primitive conversions
  val toCInt : Int -> CInt
  val fromCInt : CInt -> Int
  val toCDouble : Real -> CDouble
  val fromCDouble : CDouble -> Real

  -- String marshalling
  val toCString : String -> CString ! IOEffect
  val fromCString : CString -> String ! IOEffect
  val withCString : forall a. String -> (CString -> a ! IOEffect) -> a ! IOEffect

  -- Pointer operations
  val nullPtr : forall a. Ptr[a]
  val isNull : forall a. Ptr[a] -> Bool
  val ptrToInt : forall a. Ptr[a] -> Int
  val intToPtr : forall a. Int -> Ptr[a]

  -- Array marshalling
  val toArray : forall a. List[a] -> CArray[a] ! IOEffect
  val fromArray : forall a. CArray[a] -> CSize -> List[a] ! IOEffect
  val withArray : forall a b. List[a] -> (Ptr[a] -> CSize -> b ! IOEffect)
    -> b ! IOEffect
}
\end{verbatim}
\end{definition}

\subsection{Custom Marshallers}
\label{subsec:custom-marshallers}

\begin{definition}[Marshallable Typeclass]
\label{def:marshallable}
The \textsf{Marshallable} typeclass for custom marshalling:
\begin{verbatim}
signature MARSHALLABLE {
  type T           -- TKS type
  type Foreign     -- Foreign representation

  val marshal : T -> Foreign ! IOEffect
  val unmarshal : Foreign -> T ! IOEffect
  val sizeOf : CSize
  val alignment : CSize
}

-- Instance for Element
module ElementMarshal : MARSHALLABLE {
  type T = Element
  type Foreign = CInt  -- Packed as 0-39

  def marshal(Element(w, n)) =
    toCInt(wingIndex(w) * 10 + (n - 1))

  def unmarshal(i) =
    let idx = fromCInt(i)
    let w = indexToWing(idx / 10)
    let n = (idx % 10) + 1
    Element(w, n)

  def sizeOf = 4
  def alignment = 4
}
\end{verbatim}
\end{definition}

\begin{definition}[Struct Marshalling]
\label{def:struct-marshal}
Marshalling TKS records to C structs:
\begin{verbatim}
-- Define C struct layout
foreign struct RPMState {
  acquired_count: CInt,
  goal_element: CInt,
  status: CInt
}

-- TKS record type
type RPMStateRecord = {
  acquiredCount: Int,
  goalElement: Option[Element],
  status: RPMStatus
}

-- Marshalling instance
module RPMStateMarshal : MARSHALLABLE {
  type T = RPMStateRecord
  type Foreign = Struct[RPMState]

  def marshal(r) = {
    acquired_count = toCInt(r.acquiredCount),
    goal_element = case r.goalElement of
      None -> -1
      Some(e) -> ElementMarshal.marshal(e),
    status = statusToInt(r.status)
  }

  def unmarshal(s) = {
    acquiredCount = fromCInt(s.acquired_count),
    goalElement = if s.goal_element < 0 then None
                  else Some(ElementMarshal.unmarshal(s.goal_element)),
    status = intToStatus(s.status)
  }
}
\end{verbatim}
\end{definition}

\subsection{TKS Domain Type Marshalling}
\label{subsec:domain-marshalling}

\begin{definition}[Noetic Value Marshalling]
\label{def:noetic-marshal}
Marshalling TKS-specific types:
\begin{verbatim}
-- Noetic to foreign representation
module NoeticMarshal : MARSHALLABLE {
  type T = Noetic
  type Foreign = CInt  -- Level index 0-3

  def marshal(nu) = toCInt(level(nu))
  def unmarshal(i) = case fromCInt(i) of
    0 -> nu_0
    1 -> nu_1
    2 -> nu_2
    3 -> nu_3
    _ -> raise NoeticError.invalidLevel(fromCInt(i))
}

-- Idea to JSON string for interop
module IdeaMarshal : MARSHALLABLE {
  type T = Idea
  type Foreign = CString

  def marshal(idea) =
    toCString(toJSON(idea))

  def unmarshal(s) =
    fromJSON(fromCString(s))
}

-- Fractal to nested array
module FractalMarshal[A : MARSHALLABLE] : MARSHALLABLE {
  type T = Fractal[A.T]
  type Foreign = Ptr[Ptr[A.Foreign]]  -- Array of level arrays

  def marshal(f) = ...  -- Level-by-level marshalling
  def unmarshal(p) = ... -- Reconstruct fractal
}
\end{verbatim}
\end{definition}

%% ----------------------------------------------------------------------------
\section{Safety and Effects}
\label{sec:ffi-safety}

This section specifies the safety model and effect tracking for FFI operations.

\subsection{Safe vs Unsafe FFI}
\label{subsec:safe-unsafe}

\begin{definition}[FFI Safety Levels]
\label{def:ffi-safety-levels}
TKS distinguishes two FFI safety levels:

\textbf{Safe FFI} (\texttt{safe}):
\begin{itemize}
  \item Foreign function must not:
    \begin{itemize}
      \item Dereference invalid pointers
      \item Cause undefined behavior for valid inputs
      \item Modify TKS-managed memory
    \end{itemize}
  \item Compiler may call during garbage collection pauses
  \item May be inlined across FFI boundary
\end{itemize}

\textbf{Unsafe FFI} (\texttt{unsafe}):
\begin{itemize}
  \item No restrictions on foreign function behavior
  \item Caller responsible for ensuring safety invariants
  \item Cannot be called during GC; requires synchronization
  \item Marked in type system, requires explicit \texttt{unsafe} block
\end{itemize}
\end{definition}

\begin{definition}[Unsafe Blocks]
\label{def:unsafe-blocks}
Calling unsafe foreign functions requires \texttt{unsafe} block:
\begin{verbatim}
def allocateBuffer(size: Int) : Ptr ! IOEffect =
  unsafe {
    let ptr = malloc(toCSize(size))
    if isNull(ptr) then
      raise MemoryError.allocationFailed(size)
    else
      ptr
  }

-- Safe wrapper hides unsafety
def withBuffer[A](size: Int, f: Ptr -> A ! IOEffect) : A ! IOEffect =
  let ptr = allocateBuffer(size)
  let result = f(ptr)
  unsafe { free(ptr) }
  result
\end{verbatim}
\end{definition}

\subsection{FFI Effects}
\label{subsec:ffi-effects}

\begin{definition}[IOEffect]
\label{def:io-effect}
The \textsf{IOEffect} captures side effects from FFI:
\begin{verbatim}
effect IOEffect {
  -- File system operations
  readFile : FilePath -> Bytes
  writeFile : FilePath -> Bytes -> Unit
  deleteFile : FilePath -> Unit

  -- Network operations
  connect : Address -> Socket
  send : Socket -> Bytes -> Int
  recv : Socket -> Int -> Bytes

  -- Process operations
  spawn : Command -> Process
  wait : Process -> ExitCode

  -- Raw foreign call effect
  foreignCall : Unit -> Unit
}
\end{verbatim}
\end{definition}

\begin{definition}[Effect Inference for FFI]
\label{def:ffi-effect-inference}
FFI calls are typed with effects:
\[
\frac{\texttt{external "C" safe fn } f(x_1: \tau_1, \ldots) : \tau_r \mathbin{!} \varepsilon}
     {\Gamma \vdash f : \tau_1 \to \cdots \to \tau_r \mathbin{!} \varepsilon}
\]

Default effects by safety level:
\begin{itemize}
  \item \texttt{safe} without annotation: $\varepsilon = \emptyset$ (pure)
  \item \texttt{safe} with I/O: $\varepsilon = \mathsf{IOEffect}$
  \item \texttt{unsafe}: $\varepsilon = \mathsf{IOEffect} \cup \mathsf{UnsafeEffect}$
\end{itemize}
\end{definition}

\begin{definition}[UnsafeEffect]
\label{def:unsafe-effect}
The \textsf{UnsafeEffect} for memory-unsafe operations:
\begin{verbatim}
effect UnsafeEffect {
  derefPtr : forall a. Ptr[a] -> a
  writePtr : forall a. Ptr[a] -> a -> Unit
  ptrArith : forall a. Ptr[a] -> Int -> Ptr[a]
  castPtr : forall a b. Ptr[a] -> Ptr[b]
}

-- Handler that validates pointer operations
handler SafeMemory : UnsafeEffect => IOEffect {
  derefPtr(p) -> resume(checkedDeref(p))
  writePtr(p, v) -> resume(checkedWrite(p, v))
  ptrArith(p, n) -> resume(checkedOffset(p, n))
  castPtr(p) -> resume(checkedCast(p))
}
\end{verbatim}
\end{definition}

\subsection{Memory Safety}
\label{subsec:memory-safety}

\begin{definition}[Pointer Lifetime Tracking]
\label{def:pointer-lifetime}
TKS tracks pointer lifetimes to prevent use-after-free:
\begin{verbatim}
-- Region-based memory management
region R {
  -- Pointers allocated in R are valid only within this block
  let ptr = allocIn[R](100)
  usePtr(ptr)
  -- ptr automatically freed at end of region
}
-- ptr is no longer valid here

-- Linear types for unique ownership
def consumePtr(ptr: Ptr[a] @owned) : a ! IOEffect =
  let v = deref(ptr)
  free(ptr)  -- Must free exactly once
  v

-- Borrowed references
def readPtr(ptr: Ptr[a] @borrowed) : a =
  deref(ptr)  -- Cannot free borrowed pointer
\end{verbatim}
\end{definition}

\begin{definition}[GC Interaction]
\label{def:gc-interaction}
Rules for FFI and garbage collection:
\begin{enumerate}
  \item \textbf{Pinning}: TKS values passed to foreign code are pinned (not moved by GC).
  \item \textbf{Safe points}: Safe FFI calls are safe points for GC.
  \item \textbf{Unsafe calls}: Unsafe FFI calls disable GC for duration.
  \item \textbf{Pointers to TKS heap}: Never valid to store; use stable pointers instead.
\end{enumerate}
\begin{verbatim}
-- Stable pointer for long-lived foreign references
def withStablePtr[A, B](v: A, f: StablePtr[A] -> B ! IOEffect) : B ! IOEffect =
  let sp = newStablePtr(v)  -- Pin value, get stable reference
  let result = f(sp)
  freeStablePtr(sp)         -- Unpin value
  result
\end{verbatim}
\end{definition}

%% ----------------------------------------------------------------------------
\section{Callbacks}
\label{sec:ffi-callbacks}

This section specifies passing TKS functions to foreign code.

\subsection{Function Pointer Creation}
\label{subsec:funptr-creation}

\begin{definition}[Callback Syntax]
\label{def:callback-syntax}
Creating function pointers from TKS functions:
\begin{verbatim}
-- Type for C callback: int (*)(int, void*)
type CCallback = FunPtr[CInt -> Ptr -> CInt]

-- Create callback from TKS function
def makeCallback(f: Int -> Int) : CCallback ! IOEffect =
  wrapCallback(\(x: CInt, _: Ptr) -> toCInt(f(fromCInt(x))))

-- Example: qsort comparator
external "C" safe fn qsort(
  base: Ptr,
  nmemb: CSize,
  size: CSize,
  compar: FunPtr[Ptr -> Ptr -> CInt]
) : Unit ! IOEffect

def sortInts(arr: List[Int]) : List[Int] ! IOEffect =
  withArray(arr, \(ptr, len) ->
    let cmp = wrapCallback(\(a: Ptr, b: Ptr) ->
      let x = derefInt(a)
      let y = derefInt(b)
      toCInt(compare(x, y))
    )
    qsort(ptr, len, sizeOfCInt, cmp)
    freeCallback(cmp)
    fromArray(ptr, len)
  )
\end{verbatim}
\end{definition}

\begin{definition}[Closure Conversion]
\label{def:closure-conversion}
TKS closures require special handling for FFI:
\begin{enumerate}
  \item \textbf{No captures}: Direct function pointer.
  \item \textbf{With captures}: Closure + context pointer pair.
\end{enumerate}
\begin{verbatim}
-- Closure with captured environment
def makeCounter(start: Int) : (() -> Int) =
  let count = ref start
  \() -> let n = !count in (count := n + 1; n)

-- For FFI, split into function + context
type ClosureFFI[A, B] = {
  fn: FunPtr[Ptr -> A -> B],  -- Function taking context
  ctx: Ptr                     -- Captured environment
}

def wrapClosure[A, B](f: A -> B) : ClosureFFI[A, B] ! IOEffect =
  let ctx = allocContext(f)
  let fn = wrapWithContext(\(c: Ptr, a: A) -> applyClosure(c, a))
  { fn = fn, ctx = ctx }
\end{verbatim}
\end{definition}

\subsection{Callback Lifetime}
\label{subsec:callback-lifetime}

\begin{definition}[Callback Registration]
\label{def:callback-registration}
Managing callback lifetimes:
\begin{verbatim}
-- Callbacks must be explicitly freed
def withCallback[A, B, R](
  f: A -> B,
  use: FunPtr[A -> B] -> R ! IOEffect
) : R ! IOEffect =
  let cb = wrapCallback(f)
  let result = use(cb)
  freeCallback(cb)
  result

-- Long-lived callbacks (e.g., event handlers)
module CallbackRegistry {
  type CallbackId = Int

  def register[A, B](f: A -> B) : CallbackId ! IOEffect
  def unregister(id: CallbackId) : Unit ! IOEffect
  def getPtr[A, B](id: CallbackId) : FunPtr[A -> B]
}
\end{verbatim}
\end{definition}

%% ----------------------------------------------------------------------------
\section{Platform Bindings}
\label{sec:platform-bindings}

This section specifies standard platform interfaces.

\subsection{File System}
\label{subsec:filesystem}

\begin{definition}[TKS.IO.FileSystem Interface]
\label{def:filesystem-interface}
\begin{verbatim}
interface TKS.IO.FileSystem {
  type FilePath = String
  type FileHandle
  type FileMode = Read | Write | Append | ReadWrite

  -- Basic operations
  val open : FilePath -> FileMode -> RPM[FileHandle] ! IOEffect
  val close : FileHandle -> Unit ! IOEffect
  val read : FileHandle -> Int -> RPM[Bytes] ! IOEffect
  val write : FileHandle -> Bytes -> RPM[Int] ! IOEffect

  -- Convenience functions
  val readFile : FilePath -> RPM[String] ! IOEffect
  val writeFile : FilePath -> String -> RPM[Unit] ! IOEffect
  val appendFile : FilePath -> String -> RPM[Unit] ! IOEffect

  -- Directory operations
  val listDir : FilePath -> RPM[List[FilePath]] ! IOEffect
  val createDir : FilePath -> RPM[Unit] ! IOEffect
  val removeDir : FilePath -> RPM[Unit] ! IOEffect
  val exists : FilePath -> Bool ! IOEffect

  -- Metadata
  val fileSize : FilePath -> RPM[Int] ! IOEffect
  val modTime : FilePath -> RPM[Time] ! IOEffect
}
\end{verbatim}
\end{definition}

\subsection{Networking}
\label{subsec:networking}

\begin{definition}[TKS.IO.Network Interface]
\label{def:network-interface}
\begin{verbatim}
interface TKS.IO.Network {
  type Socket
  type Address = { host: String, port: Int }
  type Protocol = TCP | UDP

  -- Connection management
  val connect : Address -> Protocol -> RPM[Socket] ! IOEffect
  val listen : Address -> Protocol -> RPM[Socket] ! IOEffect
  val accept : Socket -> RPM[Socket] ! IOEffect
  val close : Socket -> Unit ! IOEffect

  -- Data transfer
  val send : Socket -> Bytes -> RPM[Int] ! IOEffect
  val recv : Socket -> Int -> RPM[Bytes] ! IOEffect
  val sendAll : Socket -> Bytes -> RPM[Unit] ! IOEffect
  val recvExact : Socket -> Int -> RPM[Bytes] ! IOEffect

  -- DNS
  val resolve : String -> RPM[List[Address]] ! IOEffect

  -- Options
  val setTimeout : Socket -> Duration -> Unit ! IOEffect
  val setNoDelay : Socket -> Bool -> Unit ! IOEffect
}
\end{verbatim}
\end{definition}

\subsection{Process Management}
\label{subsec:process}

\begin{definition}[TKS.IO.Process Interface]
\label{def:process-interface}
\begin{verbatim}
interface TKS.IO.Process {
  type Process
  type ExitCode = Success | Failure(Int)
  type Stdio = Inherit | Pipe | Null | File(FilePath)

  type ProcessConfig = {
    program: String,
    args: List[String],
    env: Option[List[(String, String)]],
    cwd: Option[FilePath],
    stdin: Stdio,
    stdout: Stdio,
    stderr: Stdio
  }

  -- Process control
  val spawn : ProcessConfig -> RPM[Process] ! IOEffect
  val wait : Process -> RPM[ExitCode] ! IOEffect
  val kill : Process -> Unit ! IOEffect

  -- I/O with child
  val writeStdin : Process -> Bytes -> RPM[Unit] ! IOEffect
  val readStdout : Process -> RPM[Bytes] ! IOEffect
  val readStderr : Process -> RPM[Bytes] ! IOEffect

  -- Convenience
  val run : String -> List[String] -> RPM[ExitCode] ! IOEffect
  val output : String -> List[String] -> RPM[String] ! IOEffect

  -- Current process
  val getArgs : List[String] ! IOEffect
  val getEnv : String -> Option[String] ! IOEffect
  val setEnv : String -> String -> Unit ! IOEffect
  val exit : ExitCode -> Unit ! IOEffect
}
\end{verbatim}
\end{definition}

\subsection{Concurrency Primitives}
\label{subsec:concurrency-primitives}

\begin{definition}[TKS.IO.Concurrency Interface]
\label{def:concurrency-interface}
\begin{verbatim}
interface TKS.IO.Concurrency {
  -- Thread operations
  type Thread[a]

  val spawn : forall a. (() -> a ! IOEffect) -> Thread[a] ! IOEffect
  val join : forall a. Thread[a] -> RPM[a] ! IOEffect
  val yield : Unit ! IOEffect

  -- Synchronization primitives
  type Mutex
  val newMutex : Mutex ! IOEffect
  val lock : Mutex -> Unit ! IOEffect
  val unlock : Mutex -> Unit ! IOEffect
  val withLock : forall a. Mutex -> (() -> a ! IOEffect) -> a ! IOEffect

  type Semaphore
  val newSemaphore : Int -> Semaphore ! IOEffect
  val acquire : Semaphore -> Unit ! IOEffect
  val release : Semaphore -> Unit ! IOEffect

  -- Channels
  type Channel[a]
  val newChannel : forall a. Channel[a] ! IOEffect
  val send : forall a. Channel[a] -> a -> Unit ! IOEffect
  val recv : forall a. Channel[a] -> a ! IOEffect
  val tryRecv : forall a. Channel[a] -> Option[a] ! IOEffect

  -- Atomic operations
  type Atomic[a]
  val newAtomic : forall a. a -> Atomic[a] ! IOEffect
  val readAtomic : forall a. Atomic[a] -> a ! IOEffect
  val writeAtomic : forall a. Atomic[a] -> a -> Unit ! IOEffect
  val casAtomic : forall a. Atomic[a] -> a -> a -> Bool ! IOEffect
}
\end{verbatim}
\end{definition}

\subsection{External Library Integration}
\label{subsec:external-libs}

\begin{definition}[Library Binding Pattern]
\label{def:lib-binding-pattern}
Standard pattern for binding external libraries:
\begin{verbatim}
-- Example: Binding to a JSON library
module TKS.FFI.JSON {
  link "jansson"  -- Link libjansson

  -- Opaque types
  type JSONValue
  type JSONError

  -- Raw FFI declarations
  private external "C" safe fn json_loads(s: CString, flags: CInt, err: Ptr)
    : Ptr[JSONValue] ! IOEffect
  private external "C" safe fn json_dumps(v: Ptr[JSONValue], flags: CInt)
    : CString ! IOEffect
  private external "C" safe fn json_decref(v: Ptr[JSONValue])
    : Unit ! IOEffect

  -- Safe TKS interface
  export {
    parse : String -> RPM[JSONValue] ! IOEffect,
    stringify : JSONValue -> RPM[String] ! IOEffect
  }

  def parse(s: String) : RPM[JSONValue] ! IOEffect =
    withCString(s, \cs ->
      let ptr = json_loads(cs, 0, nullPtr)
      if isNull(ptr) then
        fail "JSON parse error"
      else
        return (wrapJSON(ptr))
    )

  def stringify(v: JSONValue) : RPM[String] ! IOEffect =
    let cs = json_dumps(unwrapJSON(v), 0)
    if isNull(cs) then
      fail "JSON stringify error"
    else
      let s = fromCString(cs)
      free(cs)
      return s
}
\end{verbatim}
\end{definition}

%% ============================================================================
%% CHAPTER 7 SUMMARY AND V7.3 COMPILER TRACK HANDOFF
%% ============================================================================
\section{Chapter 7 Summary and v7.3 Compiler Track Handoff}
\label{sec:compiler-handoff}

\begin{tcolorbox}[colback=green!5!white,colframe=green!50!black,title=\textbf{COMPILER-TASK-07 Complete: FFI Chapter},breakable]
\textbf{Completed in this chapter:}
\begin{itemize}
  \item \textbf{FFI Declarations} (Section~\ref{sec:ffi-declarations}):
    \begin{itemize}
      \item External function declaration syntax
      \item Calling conventions: C, stdcall, fastcall, system
      \item Safety annotations: safe vs unsafe
      \item Library linking directives
      \item Symbol resolution and name mapping
      \item Variadic function handling
    \end{itemize}
  \item \textbf{Type Marshalling} (Section~\ref{sec:type-marshalling}):
    \begin{itemize}
      \item Foreign type system (CInt, Ptr, CString, Struct, FunPtr)
      \item Platform-dependent type sizes
      \item Automatic marshalling for primitives
      \item Marshallable typeclass for custom types
      \item Struct marshalling
      \item TKS domain type marshalling (Element, Noetic, Idea, Fractal)
    \end{itemize}
  \item \textbf{Safety and Effects} (Section~\ref{sec:ffi-safety}):
    \begin{itemize}
      \item Safe vs unsafe FFI distinction
      \item Unsafe blocks for explicit unsafety
      \item IOEffect and UnsafeEffect definitions
      \item Effect inference rules for FFI
      \item Pointer lifetime tracking (regions, linear types, borrowing)
      \item GC interaction (pinning, safe points, stable pointers)
    \end{itemize}
  \item \textbf{Callbacks} (Section~\ref{sec:ffi-callbacks}):
    \begin{itemize}
      \item Function pointer creation from TKS functions
      \item Closure conversion for captured environments
      \item Callback lifetime management
      \item Callback registry for long-lived handlers
    \end{itemize}
  \item \textbf{Platform Bindings} (Section~\ref{sec:platform-bindings}):
    \begin{itemize}
      \item TKS.IO.FileSystem interface
      \item TKS.IO.Network interface
      \item TKS.IO.Process interface
      \item TKS.IO.Concurrency interface
      \item External library integration pattern
    \end{itemize}
\end{itemize}
\end{tcolorbox}

\subsection{Complete v7.3 Compiler Track Summary}
\label{subsec:compiler-track-summary}

The TKS v7.3 Compiler specification comprises seven chapters providing a complete compiler architecture:

\begin{center}
\begin{tabular}{clp{8cm}}
\toprule
\textbf{Ch.} & \textbf{Title} & \textbf{Key Contents} \\
\midrule
1 & Lexical Structure & Tokens, keywords, operators, literals, comments, whitespace rules \\
2 & Concrete Syntax & BNF grammar, expressions, statements, declarations, precedence \\
3 & Abstract Syntax & AST types, core calculus, desugaring transformations \\
4 & Type System & HM inference, Algorithm W, unification, TKS-specific types, effect types \\
5 & Bytecode \& VM & Stack-based VM, instruction set, operational semantics, optimization \\
6 & Module System & Modules, signatures, functors, separate compilation, standard library \\
7 & FFI & External declarations, marshalling, safety/effects, platform bindings \\
\bottomrule
\end{tabular}
\end{center}

\subsection{Key Components Checklist}
\label{subsec:components-checklist}

\begin{center}
\begin{tabular}{lcc}
\toprule
\textbf{Component} & \textbf{Specified} & \textbf{Chapter} \\
\midrule
Lexer specification & $\checkmark$ & 1 \\
Parser grammar (BNF) & $\checkmark$ & 2 \\
AST data types & $\checkmark$ & 3 \\
Core calculus & $\checkmark$ & 3 \\
Desugaring passes & $\checkmark$ & 3 \\
Type inference (Algorithm W) & $\checkmark$ & 4 \\
Unification algorithm & $\checkmark$ & 4 \\
Effect type system & $\checkmark$ & 4 \\
Bytecode instruction set & $\checkmark$ & 5 \\
VM operational semantics & $\checkmark$ & 5 \\
RPM runtime support & $\checkmark$ & 5 \\
Module system & $\checkmark$ & 6 \\
Functor modules & $\checkmark$ & 6 \\
Separate compilation & $\checkmark$ & 6 \\
Standard library interfaces & $\checkmark$ & 6 \\
FFI declarations & $\checkmark$ & 7 \\
Type marshalling & $\checkmark$ & 7 \\
Platform bindings & $\checkmark$ & 7 \\
\bottomrule
\end{tabular}
\end{center}

\subsection{Implementation Priorities}
\label{subsec:impl-priorities}

For implementation of the TKS compiler, we recommend the following priority order:

\begin{enumerate}
  \item \textbf{Phase 1: Core Pipeline}
    \begin{itemize}
      \item Lexer (Chapter 1)
      \item Parser (Chapter 2)
      \item AST construction (Chapter 3)
      \item Basic type checker without effects (Chapter 4, partial)
    \end{itemize}

  \item \textbf{Phase 2: Type System}
    \begin{itemize}
      \item Full Hindley-Milner inference (Chapter 4)
      \item TKS-specific type rules (Element, Foundation, Noetic)
      \item Effect type inference (Chapter 4)
    \end{itemize}

  \item \textbf{Phase 3: Code Generation}
    \begin{itemize}
      \item Bytecode emitter (Chapter 5)
      \item Basic VM interpreter (Chapter 5)
      \item RPM monad runtime (Chapter 5)
    \end{itemize}

  \item \textbf{Phase 4: Module System}
    \begin{itemize}
      \item Module resolution (Chapter 6)
      \item Interface file generation (Chapter 6)
      \item Separate compilation (Chapter 6)
    \end{itemize}

  \item \textbf{Phase 5: FFI and Runtime}
    \begin{itemize}
      \item FFI infrastructure (Chapter 7)
      \item Platform bindings (Chapter 7)
      \item Garbage collector integration
    \end{itemize}

  \item \textbf{Phase 6: Optimization}
    \begin{itemize}
      \item Bytecode optimization passes (Chapter 5)
      \item Inline caching for Element dispatch
      \item JIT compilation (future extension)
    \end{itemize}
\end{enumerate}

\subsection{Detailed Handoff Notes}
\label{subsec:handoff-notes}

\begin{tcolorbox}[colback=yellow!5!white,colframe=yellow!50!black,title=\textbf{Handoff: TKS v7.3 Compiler Track Complete},breakable]

\textbf{What has been accomplished:}
\begin{itemize}
  \item Complete formal specification of TKS programming language
  \item Lexical, syntactic, and semantic definitions
  \item Type system with full inference algorithm
  \item Bytecode VM with operational semantics
  \item ML-style module system with functors
  \item Foreign function interface with safety model
  \item Standard library interface specifications
\end{itemize}

\textbf{Integration with TKS v7.2 Manual:}
\begin{itemize}
  \item Type system (Chapter 4) implements types from v7.2 Chapter 7
  \item RPM runtime (Chapter 5) implements v7.2 RPM monad semantics
  \item Module system (Chapter 6) extends v7.2 Chapter 8 definitions
  \item Standard library (Chapter 6) provides interfaces for v7.2 constructs
\end{itemize}

\textbf{Dependencies on other manuals:}
\begin{itemize}
  \item v6.1: Element types (A1--D10), Foundation types, Noetic operators
  \item v7.0: RPM monad definition, Fractal structures, ACBE semantics
  \item v7.1: Extended RPM with multi-goal, denotational semantics
  \item v7.2: Full type system, effect handlers, quantum extensions
\end{itemize}

\textbf{Open items for future work:}
\begin{enumerate}
  \item \textbf{JIT Compilation}: Native code generation for performance
  \item \textbf{Debugger}: Source-level debugging with bytecode mapping
  \item \textbf{Profiler}: Performance analysis tools
  \item \textbf{Package Manager}: Dependency management for TKS modules
  \item \textbf{IDE Support}: Language server protocol implementation
  \item \textbf{Formal Verification}: Proof of compiler correctness
\end{enumerate}

\textbf{Reference implementation notes:}
\begin{itemize}
  \item Recommended implementation language: Haskell or OCaml (for type system)
  \item Alternative: Rust (for performance and FFI)
  \item VM can be implemented in C for portability
  \item Consider LLVM backend for optimized native code
\end{itemize}

\textbf{Testing strategy:}
\begin{itemize}
  \item Unit tests for each compiler phase
  \item Property-based testing for type inference
  \item Golden tests for bytecode generation
  \item Integration tests with standard library
  \item Conformance test suite against specification
\end{itemize}

\end{tcolorbox}

\begin{theorem}[Compiler Correctness]
\label{thm:compiler-correctness}
The TKS compiler preserves program semantics:
\[
\forall e.\; \Gamma \vdash e : \tau \implies \llbracket \mathsf{compile}(e) \rrbracket_{\mathsf{VM}} = \llbracket e \rrbracket_{\mathsf{denotational}}
\]
That is, executing compiled bytecode produces the same result as the denotational semantics of the source expression, for all well-typed programs.
\end{theorem}

\begin{proof}[Proof sketch]
By structural induction on typing derivations, showing:
\begin{enumerate}
  \item Each AST node compiles to semantics-preserving bytecode.
  \item Type erasure does not affect runtime behavior.
  \item Effect handlers compose correctly at runtime.
  \item Module boundaries preserve abstraction.
\end{enumerate}
Full proof requires formalizing the bytecode semantics and establishing simulation relations. This is left as future work for formal verification.
\end{proof}

